% 电磁场课程笔记 —— 主文档（正文排版集中在 0.样式.tex，章节文件只放内容）
% 编译：xelatex 1.电磁场笔记.tex 两遍；成品输出见 build.ps1
\documentclass[11pt,a4paper,oneside,openany]{book}
\usepackage[UTF8,fontset=windows,scheme=plain]{ctex}
\usepackage{fontspec}
\setmainfont{Times New Roman}
\setsansfont{Microsoft YaHei}
\newcommand{\YHBbooktrue}{}

\newcommand{\YHBcourse}{电磁场}
\newcommand{\YHBversion}{v1.1}
\newcommand{\YHBquote}{电和磁的实验中最明显的现象，是物体之间的相互作用。\par\vspace{2mm}\normalfont —— 麦克斯韦}
\newcommand{\YHBcoverSize}{46pt}
\newcommand{\YHBcoverLead}{54pt}
\usepackage{xcolor}
\definecolor{yhbAccent}{HTML}{4A2C2A}
\definecolor{yhbAccent2}{HTML}{8A5A4E}
% ============================================================
% 0.样式.tex —— 课程笔记统一正文排版
% 用法：主文件先定义课程信息与主色，再 % ============================================================
% 0.样式.tex —— 课程笔记统一正文排版
% 用法：主文件先定义课程信息与主色，再 % ============================================================
% 0.样式.tex —— 课程笔记统一正文排版
% 用法：主文件先定义课程信息与主色，再 \input{0.样式.tex}
%   \newcommand{\YHBcourse}{微积分A}      % 课程名（页眉与书签用）
%   \newcommand{\YHBversion}{v1.2}        % 版本号（封面用）
%   \newcommand{\YHBquote}{……}            % 封面题句
%   \definecolor{yhbAccent}{HTML}{24465C} % 主色（深）
%   \definecolor{yhbAccent2}{HTML}{347A80}% 辅色（浅）
%   \YHBbooktrue / \YHBbookfalse          % 是否 book 文档类（有 \chapter）
% 章节文件只写内容，不要重复 \usepackage、不要改字号。
% ============================================================

% ---------- 页面与字体 ----------
\usepackage{geometry}
\geometry{a4paper,left=24mm,right=24mm,top=23mm,bottom=24mm,headheight=15pt,headsep=7mm,footskip=13mm}
\usepackage{microtype}
\usepackage{amsmath,amssymb,amsfonts,mathtools}
\usepackage{esint}
\usepackage{graphicx}
\usepackage{adjustbox}
\usepackage{array,booktabs,multirow,makecell,tabularx,longtable}
\usepackage{xcolor,colortbl}
\usepackage{enumitem}
\usepackage{needspace,etoolbox}
\usepackage{titlesec,titletoc}
\usepackage{fancyhdr}
\usepackage{caption}
\usepackage[most]{tcolorbox}

% ---------- 颜色别名（章节里用的旧名继续有效） ----------
\colorlet{mainblue}{yhbAccent}
\colorlet{secondcyan}{yhbAccent2}
\colorlet{thirdteal}{yhbAccent2}
\colorlet{highlight}{yhbAccent!6!white}
\colorlet{tablehead}{yhbAccent!12!white}
\definecolor{textgray}{HTML}{4A5057}

% ---------- 课程标签 ----------
\providecommand{\YHBdoclabel}{课程笔记}
% ---------- 正文节奏 ----------
\linespread{1.3}
\setlength{\parindent}{0pt}
\setlength{\parskip}{6pt plus 1.5pt minus 1.5pt}
\setlength{\emergencystretch}{2.5em}
\raggedbottom
\widowpenalty=10000
\clubpenalty=10000
\setlength{\tabcolsep}{7pt}
\renewcommand{\arraystretch}{1.3}
\allowdisplaybreaks[2]
\AtBeginDocument{%
  \setlength{\abovedisplayskip}{10pt plus 2pt minus 2pt}%
  \setlength{\belowdisplayskip}{10pt plus 2pt minus 2pt}%
  \setlength{\abovedisplayshortskip}{6pt}%
  \setlength{\belowdisplayshortskip}{6pt}%
}
\captionsetup{font=small,labelfont=bf,skip=6pt}

% ---------- 列表 ----------
\setlist{leftmargin=2em,labelsep=0.6em,itemsep=4pt,parsep=2pt,topsep=6pt,partopsep=0pt}
\setlist[itemize,2]{label=\textbullet,leftmargin=1.6em}

% ---------- 标题层级 ----------
\ifdefined\YHBbooktrue
  \titleformat{\chapter}[display]
    {\sffamily\bfseries\color{yhbAccent}}
    {\normalsize\sffamily 第\,\thechapter\,章}
    {6pt}
    {\LARGE}
    [{\vspace{5pt}{\color{yhbAccent!55}\titlerule[1pt]}}]
  \titlespacing*{\chapter}{0pt}{14pt}{24pt}
  \titleformat{\section}{\Large\sffamily\bfseries\color{yhbAccent}}{\thesection}{0.7em}{}
  \titlespacing*{\section}{0pt}{18pt plus 3pt minus 2pt}{8pt}
\else
  \titleformat{\section}{\Large\sffamily\bfseries\color{yhbAccent}}{\thesection}{0.7em}{}[{\vspace{3pt}{\color{yhbAccent!45}\titlerule[0.7pt]}}]
  \titlespacing*{\section}{0pt}{18pt plus 3pt minus 2pt}{9pt}
\fi
\titleformat{\subsection}{\large\sffamily\bfseries\color{yhbAccent!85!black}}{\thesubsection}{0.7em}{}
\titleformat{\subsubsection}{\normalsize\bfseries\color{black!85}}{\thesubsubsection}{0.7em}{}
\titlespacing*{\subsection}{0pt}{12pt plus 2pt minus 1pt}{5pt}
\titlespacing*{\subsubsection}{0pt}{9pt plus 2pt minus 1pt}{4pt}
\pretocmd{\section}{\Needspace{6\baselineskip}}{}{}
\pretocmd{\subsection}{\Needspace{5\baselineskip}}{}{}
\pretocmd{\subsubsection}{\Needspace{4\baselineskip}}{}{}

\ifdefined\YHBbooktrue
  \setcounter{secnumdepth}{2}
  \setcounter{tocdepth}{1}
\else
  \setcounter{secnumdepth}{3}
  \setcounter{tocdepth}{2}
\fi
\renewcommand{\contentsname}{目录}
% 目录编号颜色由 titletoc 的 titlecontents 控制，不再重定义 \numberline

% ---------- 目录样式 ----------
\ifdefined\YHBbooktrue
  \titlecontents{chapter}[0pt]
    {\addvspace{6pt}\sffamily\bfseries\color{yhbAccent}}
    {\contentslabel{2.2em}}{}{\titlerule*[0.6em]{.}\contentspage}
  \titlecontents{section}[2.2em]
    {\addvspace{2pt}}
    {\contentslabel{2.2em}}{}{\titlerule*[0.6em]{.}\contentspage}
\else
  \titlecontents{section}[0pt]
    {\addvspace{5pt}\sffamily\bfseries\color{yhbAccent}}
    {\contentslabel{2.2em}}{}{\titlerule*[0.6em]{.}\contentspage}
  \titlecontents{subsection}[2.2em]
    {\addvspace{2pt}}
    {\contentslabel{2.4em}}{}{\titlerule*[0.6em]{.}\contentspage}
\fi

% ---------- 页眉页脚 ----------
\pagestyle{fancy}
\fancyhf{}
\fancyhead[L]{\footnotesize\sffamily\color{textgray}\YHBcourse\ \YHBdoclabel}
\ifdefined\YHBbooktrue
  \fancyhead[R]{\footnotesize\sffamily\color{textgray}\nouppercase{\leftmark}}
  \renewcommand{\chaptermark}[1]{\markboth{#1}{}}
\else
  \fancyhead[R]{\footnotesize\sffamily\color{textgray}\nouppercase{\rightmark}}
  \renewcommand{\sectionmark}[1]{\markright{#1}}
\fi
\fancyfoot[C]{\footnotesize\color{textgray}\thepage}
\renewcommand{\headrulewidth}{0.35pt}
\renewcommand{\headrule}{{\color{yhbAccent!45}\hrule width\headwidth height\headrulewidth}}
\fancypagestyle{plain}{%
  \fancyhf{}%
  \fancyfoot[C]{\footnotesize\color{textgray}\thepage}%
  \renewcommand{\headrulewidth}{0pt}%
}

% ---------- 公式与考点框 ----------
\newcommand{\keyformula}[1]{%
  \begin{tcolorbox}[enhanced,breakable,colback=highlight,colframe=yhbAccent!30!white,boxrule=0.5pt,arc=2.5pt,left=11pt,right=11pt,top=9pt,bottom=9pt,boxsep=0pt,before skip=11pt,after skip=11pt,halign=center,width=\linewidth]
  \adjustbox{max width=\linewidth-2em}{$\displaystyle #1$}
  \end{tcolorbox}}
\newcommand{\corepoint}[1]{%
  \begin{tcolorbox}[enhanced,breakable,colback=highlight,colframe=yhbAccent2,boxrule=0pt,borderline west={2.4pt}{0pt}{yhbAccent2},arc=0pt,left=11pt,right=11pt,top=8pt,bottom=8pt,before skip=11pt,after skip=11pt]
  \textbf{\color{yhbAccent}核心考点}\quad #1
  \end{tcolorbox}}
\newcommand{\symb}[2]{\textbf{#1}：#2\par}

% ---------- 超链接（务必最后加载） ----------
\usepackage[unicode=true,colorlinks=true,linkcolor=yhbAccent,citecolor=yhbAccent,urlcolor=yhbAccent!85!black,pdfborder={0 0 0},bookmarksnumbered=true,pdfstartview=FitH]{hyperref}
\usepackage{hypcap}
\hypersetup{pdftitle={\YHBcourse 课程笔记},pdfauthor={浮萍},pdfsubject={课程学习与复习笔记},pdfcreator={XeLaTeX}}

%   \newcommand{\YHBcourse}{微积分A}      % 课程名（页眉与书签用）
%   \newcommand{\YHBversion}{v1.2}        % 版本号（封面用）
%   \newcommand{\YHBquote}{……}            % 封面题句
%   \definecolor{yhbAccent}{HTML}{24465C} % 主色（深）
%   \definecolor{yhbAccent2}{HTML}{347A80}% 辅色（浅）
%   \YHBbooktrue / \YHBbookfalse          % 是否 book 文档类（有 \chapter）
% 章节文件只写内容，不要重复 \usepackage、不要改字号。
% ============================================================

% ---------- 页面与字体 ----------
\usepackage{geometry}
\geometry{a4paper,left=24mm,right=24mm,top=23mm,bottom=24mm,headheight=15pt,headsep=7mm,footskip=13mm}
\usepackage{microtype}
\usepackage{amsmath,amssymb,amsfonts,mathtools}
\usepackage{esint}
\usepackage{graphicx}
\usepackage{adjustbox}
\usepackage{array,booktabs,multirow,makecell,tabularx,longtable}
\usepackage{xcolor,colortbl}
\usepackage{enumitem}
\usepackage{needspace,etoolbox}
\usepackage{titlesec,titletoc}
\usepackage{fancyhdr}
\usepackage{caption}
\usepackage[most]{tcolorbox}

% ---------- 颜色别名（章节里用的旧名继续有效） ----------
\colorlet{mainblue}{yhbAccent}
\colorlet{secondcyan}{yhbAccent2}
\colorlet{thirdteal}{yhbAccent2}
\colorlet{highlight}{yhbAccent!6!white}
\colorlet{tablehead}{yhbAccent!12!white}
\definecolor{textgray}{HTML}{4A5057}

% ---------- 课程标签 ----------
\providecommand{\YHBdoclabel}{课程笔记}
% ---------- 正文节奏 ----------
\linespread{1.3}
\setlength{\parindent}{0pt}
\setlength{\parskip}{6pt plus 1.5pt minus 1.5pt}
\setlength{\emergencystretch}{2.5em}
\raggedbottom
\widowpenalty=10000
\clubpenalty=10000
\setlength{\tabcolsep}{7pt}
\renewcommand{\arraystretch}{1.3}
\allowdisplaybreaks[2]
\AtBeginDocument{%
  \setlength{\abovedisplayskip}{10pt plus 2pt minus 2pt}%
  \setlength{\belowdisplayskip}{10pt plus 2pt minus 2pt}%
  \setlength{\abovedisplayshortskip}{6pt}%
  \setlength{\belowdisplayshortskip}{6pt}%
}
\captionsetup{font=small,labelfont=bf,skip=6pt}

% ---------- 列表 ----------
\setlist{leftmargin=2em,labelsep=0.6em,itemsep=4pt,parsep=2pt,topsep=6pt,partopsep=0pt}
\setlist[itemize,2]{label=\textbullet,leftmargin=1.6em}

% ---------- 标题层级 ----------
\ifdefined\YHBbooktrue
  \titleformat{\chapter}[display]
    {\sffamily\bfseries\color{yhbAccent}}
    {\normalsize\sffamily 第\,\thechapter\,章}
    {6pt}
    {\LARGE}
    [{\vspace{5pt}{\color{yhbAccent!55}\titlerule[1pt]}}]
  \titlespacing*{\chapter}{0pt}{14pt}{24pt}
  \titleformat{\section}{\Large\sffamily\bfseries\color{yhbAccent}}{\thesection}{0.7em}{}
  \titlespacing*{\section}{0pt}{18pt plus 3pt minus 2pt}{8pt}
\else
  \titleformat{\section}{\Large\sffamily\bfseries\color{yhbAccent}}{\thesection}{0.7em}{}[{\vspace{3pt}{\color{yhbAccent!45}\titlerule[0.7pt]}}]
  \titlespacing*{\section}{0pt}{18pt plus 3pt minus 2pt}{9pt}
\fi
\titleformat{\subsection}{\large\sffamily\bfseries\color{yhbAccent!85!black}}{\thesubsection}{0.7em}{}
\titleformat{\subsubsection}{\normalsize\bfseries\color{black!85}}{\thesubsubsection}{0.7em}{}
\titlespacing*{\subsection}{0pt}{12pt plus 2pt minus 1pt}{5pt}
\titlespacing*{\subsubsection}{0pt}{9pt plus 2pt minus 1pt}{4pt}
\pretocmd{\section}{\Needspace{6\baselineskip}}{}{}
\pretocmd{\subsection}{\Needspace{5\baselineskip}}{}{}
\pretocmd{\subsubsection}{\Needspace{4\baselineskip}}{}{}

\ifdefined\YHBbooktrue
  \setcounter{secnumdepth}{2}
  \setcounter{tocdepth}{1}
\else
  \setcounter{secnumdepth}{3}
  \setcounter{tocdepth}{2}
\fi
\renewcommand{\contentsname}{目录}
% 目录编号颜色由 titletoc 的 titlecontents 控制，不再重定义 \numberline

% ---------- 目录样式 ----------
\ifdefined\YHBbooktrue
  \titlecontents{chapter}[0pt]
    {\addvspace{6pt}\sffamily\bfseries\color{yhbAccent}}
    {\contentslabel{2.2em}}{}{\titlerule*[0.6em]{.}\contentspage}
  \titlecontents{section}[2.2em]
    {\addvspace{2pt}}
    {\contentslabel{2.2em}}{}{\titlerule*[0.6em]{.}\contentspage}
\else
  \titlecontents{section}[0pt]
    {\addvspace{5pt}\sffamily\bfseries\color{yhbAccent}}
    {\contentslabel{2.2em}}{}{\titlerule*[0.6em]{.}\contentspage}
  \titlecontents{subsection}[2.2em]
    {\addvspace{2pt}}
    {\contentslabel{2.4em}}{}{\titlerule*[0.6em]{.}\contentspage}
\fi

% ---------- 页眉页脚 ----------
\pagestyle{fancy}
\fancyhf{}
\fancyhead[L]{\footnotesize\sffamily\color{textgray}\YHBcourse\ \YHBdoclabel}
\ifdefined\YHBbooktrue
  \fancyhead[R]{\footnotesize\sffamily\color{textgray}\nouppercase{\leftmark}}
  \renewcommand{\chaptermark}[1]{\markboth{#1}{}}
\else
  \fancyhead[R]{\footnotesize\sffamily\color{textgray}\nouppercase{\rightmark}}
  \renewcommand{\sectionmark}[1]{\markright{#1}}
\fi
\fancyfoot[C]{\footnotesize\color{textgray}\thepage}
\renewcommand{\headrulewidth}{0.35pt}
\renewcommand{\headrule}{{\color{yhbAccent!45}\hrule width\headwidth height\headrulewidth}}
\fancypagestyle{plain}{%
  \fancyhf{}%
  \fancyfoot[C]{\footnotesize\color{textgray}\thepage}%
  \renewcommand{\headrulewidth}{0pt}%
}

% ---------- 公式与考点框 ----------
\newcommand{\keyformula}[1]{%
  \begin{tcolorbox}[enhanced,breakable,colback=highlight,colframe=yhbAccent!30!white,boxrule=0.5pt,arc=2.5pt,left=11pt,right=11pt,top=9pt,bottom=9pt,boxsep=0pt,before skip=11pt,after skip=11pt,halign=center,width=\linewidth]
  \adjustbox{max width=\linewidth-2em}{$\displaystyle #1$}
  \end{tcolorbox}}
\newcommand{\corepoint}[1]{%
  \begin{tcolorbox}[enhanced,breakable,colback=highlight,colframe=yhbAccent2,boxrule=0pt,borderline west={2.4pt}{0pt}{yhbAccent2},arc=0pt,left=11pt,right=11pt,top=8pt,bottom=8pt,before skip=11pt,after skip=11pt]
  \textbf{\color{yhbAccent}核心考点}\quad #1
  \end{tcolorbox}}
\newcommand{\symb}[2]{\textbf{#1}：#2\par}

% ---------- 超链接（务必最后加载） ----------
\usepackage[unicode=true,colorlinks=true,linkcolor=yhbAccent,citecolor=yhbAccent,urlcolor=yhbAccent!85!black,pdfborder={0 0 0},bookmarksnumbered=true,pdfstartview=FitH]{hyperref}
\usepackage{hypcap}
\hypersetup{pdftitle={\YHBcourse 课程笔记},pdfauthor={浮萍},pdfsubject={课程学习与复习笔记},pdfcreator={XeLaTeX}}

%   \newcommand{\YHBcourse}{微积分A}      % 课程名（页眉与书签用）
%   \newcommand{\YHBversion}{v1.2}        % 版本号（封面用）
%   \newcommand{\YHBquote}{……}            % 封面题句
%   \definecolor{yhbAccent}{HTML}{24465C} % 主色（深）
%   \definecolor{yhbAccent2}{HTML}{347A80}% 辅色（浅）
%   \YHBbooktrue / \YHBbookfalse          % 是否 book 文档类（有 \chapter）
% 章节文件只写内容，不要重复 \usepackage、不要改字号。
% ============================================================

% ---------- 页面与字体 ----------
\usepackage{geometry}
\geometry{a4paper,left=24mm,right=24mm,top=23mm,bottom=24mm,headheight=15pt,headsep=7mm,footskip=13mm}
\usepackage{microtype}
\usepackage{amsmath,amssymb,amsfonts,mathtools}
\usepackage{esint}
\usepackage{graphicx}
\usepackage{adjustbox}
\usepackage{array,booktabs,multirow,makecell,tabularx,longtable}
\usepackage{xcolor,colortbl}
\usepackage{enumitem}
\usepackage{needspace,etoolbox}
\usepackage{titlesec,titletoc}
\usepackage{fancyhdr}
\usepackage{caption}
\usepackage[most]{tcolorbox}

% ---------- 颜色别名（章节里用的旧名继续有效） ----------
\colorlet{mainblue}{yhbAccent}
\colorlet{secondcyan}{yhbAccent2}
\colorlet{thirdteal}{yhbAccent2}
\colorlet{highlight}{yhbAccent!6!white}
\colorlet{tablehead}{yhbAccent!12!white}
\definecolor{textgray}{HTML}{4A5057}

% ---------- 课程标签 ----------
\providecommand{\YHBdoclabel}{课程笔记}
% ---------- 正文节奏 ----------
\linespread{1.3}
\setlength{\parindent}{0pt}
\setlength{\parskip}{6pt plus 1.5pt minus 1.5pt}
\setlength{\emergencystretch}{2.5em}
\raggedbottom
\widowpenalty=10000
\clubpenalty=10000
\setlength{\tabcolsep}{7pt}
\renewcommand{\arraystretch}{1.3}
\allowdisplaybreaks[2]
\AtBeginDocument{%
  \setlength{\abovedisplayskip}{10pt plus 2pt minus 2pt}%
  \setlength{\belowdisplayskip}{10pt plus 2pt minus 2pt}%
  \setlength{\abovedisplayshortskip}{6pt}%
  \setlength{\belowdisplayshortskip}{6pt}%
}
\captionsetup{font=small,labelfont=bf,skip=6pt}

% ---------- 列表 ----------
\setlist{leftmargin=2em,labelsep=0.6em,itemsep=4pt,parsep=2pt,topsep=6pt,partopsep=0pt}
\setlist[itemize,2]{label=\textbullet,leftmargin=1.6em}

% ---------- 标题层级 ----------
\ifdefined\YHBbooktrue
  \titleformat{\chapter}[display]
    {\sffamily\bfseries\color{yhbAccent}}
    {\normalsize\sffamily 第\,\thechapter\,章}
    {6pt}
    {\LARGE}
    [{\vspace{5pt}{\color{yhbAccent!55}\titlerule[1pt]}}]
  \titlespacing*{\chapter}{0pt}{14pt}{24pt}
  \titleformat{\section}{\Large\sffamily\bfseries\color{yhbAccent}}{\thesection}{0.7em}{}
  \titlespacing*{\section}{0pt}{18pt plus 3pt minus 2pt}{8pt}
\else
  \titleformat{\section}{\Large\sffamily\bfseries\color{yhbAccent}}{\thesection}{0.7em}{}[{\vspace{3pt}{\color{yhbAccent!45}\titlerule[0.7pt]}}]
  \titlespacing*{\section}{0pt}{18pt plus 3pt minus 2pt}{9pt}
\fi
\titleformat{\subsection}{\large\sffamily\bfseries\color{yhbAccent!85!black}}{\thesubsection}{0.7em}{}
\titleformat{\subsubsection}{\normalsize\bfseries\color{black!85}}{\thesubsubsection}{0.7em}{}
\titlespacing*{\subsection}{0pt}{12pt plus 2pt minus 1pt}{5pt}
\titlespacing*{\subsubsection}{0pt}{9pt plus 2pt minus 1pt}{4pt}
\pretocmd{\section}{\Needspace{6\baselineskip}}{}{}
\pretocmd{\subsection}{\Needspace{5\baselineskip}}{}{}
\pretocmd{\subsubsection}{\Needspace{4\baselineskip}}{}{}

\ifdefined\YHBbooktrue
  \setcounter{secnumdepth}{2}
  \setcounter{tocdepth}{1}
\else
  \setcounter{secnumdepth}{3}
  \setcounter{tocdepth}{2}
\fi
\renewcommand{\contentsname}{目录}
% 目录编号颜色由 titletoc 的 titlecontents 控制，不再重定义 \numberline

% ---------- 目录样式 ----------
\ifdefined\YHBbooktrue
  \titlecontents{chapter}[0pt]
    {\addvspace{6pt}\sffamily\bfseries\color{yhbAccent}}
    {\contentslabel{2.2em}}{}{\titlerule*[0.6em]{.}\contentspage}
  \titlecontents{section}[2.2em]
    {\addvspace{2pt}}
    {\contentslabel{2.2em}}{}{\titlerule*[0.6em]{.}\contentspage}
\else
  \titlecontents{section}[0pt]
    {\addvspace{5pt}\sffamily\bfseries\color{yhbAccent}}
    {\contentslabel{2.2em}}{}{\titlerule*[0.6em]{.}\contentspage}
  \titlecontents{subsection}[2.2em]
    {\addvspace{2pt}}
    {\contentslabel{2.4em}}{}{\titlerule*[0.6em]{.}\contentspage}
\fi

% ---------- 页眉页脚 ----------
\pagestyle{fancy}
\fancyhf{}
\fancyhead[L]{\footnotesize\sffamily\color{textgray}\YHBcourse\ \YHBdoclabel}
\ifdefined\YHBbooktrue
  \fancyhead[R]{\footnotesize\sffamily\color{textgray}\nouppercase{\leftmark}}
  \renewcommand{\chaptermark}[1]{\markboth{#1}{}}
\else
  \fancyhead[R]{\footnotesize\sffamily\color{textgray}\nouppercase{\rightmark}}
  \renewcommand{\sectionmark}[1]{\markright{#1}}
\fi
\fancyfoot[C]{\footnotesize\color{textgray}\thepage}
\renewcommand{\headrulewidth}{0.35pt}
\renewcommand{\headrule}{{\color{yhbAccent!45}\hrule width\headwidth height\headrulewidth}}
\fancypagestyle{plain}{%
  \fancyhf{}%
  \fancyfoot[C]{\footnotesize\color{textgray}\thepage}%
  \renewcommand{\headrulewidth}{0pt}%
}

% ---------- 公式与考点框 ----------
\newcommand{\keyformula}[1]{%
  \begin{tcolorbox}[enhanced,breakable,colback=highlight,colframe=yhbAccent!30!white,boxrule=0.5pt,arc=2.5pt,left=11pt,right=11pt,top=9pt,bottom=9pt,boxsep=0pt,before skip=11pt,after skip=11pt,halign=center,width=\linewidth]
  \adjustbox{max width=\linewidth-2em}{$\displaystyle #1$}
  \end{tcolorbox}}
\newcommand{\corepoint}[1]{%
  \begin{tcolorbox}[enhanced,breakable,colback=highlight,colframe=yhbAccent2,boxrule=0pt,borderline west={2.4pt}{0pt}{yhbAccent2},arc=0pt,left=11pt,right=11pt,top=8pt,bottom=8pt,before skip=11pt,after skip=11pt]
  \textbf{\color{yhbAccent}核心考点}\quad #1
  \end{tcolorbox}}
\newcommand{\symb}[2]{\textbf{#1}：#2\par}

% ---------- 超链接（务必最后加载） ----------
\usepackage[unicode=true,colorlinks=true,linkcolor=yhbAccent,citecolor=yhbAccent,urlcolor=yhbAccent!85!black,pdfborder={0 0 0},bookmarksnumbered=true,pdfstartview=FitH]{hyperref}
\usepackage{hypcap}
\hypersetup{pdftitle={\YHBcourse 课程笔记},pdfauthor={浮萍},pdfsubject={课程学习与复习笔记},pdfcreator={XeLaTeX}}

\usepackage{empheq}
% 圈码 ①-⑳ 与箭头 → ⇒ ⇄ 在 Times New Roman 下无字形，交给中文字体渲染
\xeCJKDeclareCharClass{CJK}{"2460->"24FF, "2190->"21FF}
\renewcommand{\tablename}{表}
\renewcommand{\figurename}{图}
\newtcolorbox{keyformulabox}{colback=yhbAccent!6!white,colframe=yhbAccent!25!white,boxrule=0.4pt,arc=2pt,left=3mm,right=3mm,top=2mm,bottom=2mm,halign=center,breakable}
\usepackage{tikz}
\begin{document}
% 封面
\begin{titlepage}
    \centering
    \vspace*{2cm}
    
    % 主标题
    {\Huge \bfseries 电磁场作业集 \par}
    \vspace{0.8cm}
    
    % 副标题+课程信息
    {\Large \bfseries 考前复习 \par}
    \vspace{0.4cm}
    {\Large 课程代码：\textbf{EE2003} \quad 课程名称：\textbf{电磁场} \par}
    \vspace{0.6cm}
    
    \rule{\linewidth}{0.2mm}
    \vspace{1cm}
    
    % 麦克斯韦名言
    {\itshape \large 电和磁的实验中最明显的现象是，物体之间的相互作用。\par}
    \vspace{0.4cm}
    {\large —— 麦克斯韦 \par}
    \vspace{1.2cm}
    
    % 作者信息
    {\Large 作者：\textbf{浮萍} \par}
    \vspace{0.3cm}
    {\Large 版本：\textbf{1.0} \par}
    \vspace{0.3cm}
    {\Large 邮箱：\textbf{3071974740@qq.com} \par}
    \vspace{0.3cm}
    {\Large GitHub：\textbf{github.com/Duckweed-yhb} \par}
    \vspace{1cm}
    
    % 版权信息
    {\large \copyright\ 2026 \textbf{浮萍} 版权所有 \par}
    {\large \textbf{未经授权，不得复制、传播或修改} \par}
    \vspace{0.5cm}
    
    \vfill
    {\large 生成日期：\textbf{\today} \par}
\end{titlepage}

\section*{版本说明}
\noindent\textbf{版本：V1.0}
\begin{itemize}
    \item \textbf{配套课程：}课程代码\textbf{EE2003}，课程名称\textbf{电磁场}；
    \item \textbf{参考教材：}\textbf{《工程电磁场导论》}，冯慈璋、马西奎主编，西安交通大学出版社2000版，全书内容严格对标教材目录与期末考点；
    \item \textbf{覆盖范围：}全书六大模块：\textbf{第0章矢量分析、第1章静电场、第2章恒定电场、第3章恒定磁场、第4章时变电磁场、第5章准静态电磁场}，对应教材$\S0\sim\S5$全部课内知识点；
    \item \textbf{内容编排：}
    \begin{itemize}
        \item 0.矢量分析：梯度、散度、旋度、高斯公式、斯托克斯公式、正交坐标系算子；
        \item 1.静电场：电场/电位、高斯定理、边值问题、有限差分法(实验内容)、镜像法、电容、静电能量与力（采用库仑定律+虚位移法求电场力）；
        \item 2.恒定电场：导电媒质电流、局外场强、基本方程、分界面衔接条件、恒定电场与静电场类比、电导；
        \item 3.恒定磁场：磁感应强度、安培环路定律、基本方程、磁矢位、标量磁位、磁场能量与力（洛伦兹公式+虚位移法求磁场力）；
        \item 4.时变电磁场：法拉第电磁感应+全电流定律、麦克斯韦方程组、动态位及其积分解、坡印廷矢量、正弦电磁场、电磁辐射；
        \item 5.准静态场：EQS电准静态/MQS磁准静态场、似稳条件、集肤效应、导体交流内阻抗；
    \end{itemize}
    \item \textbf{现存缺点：全篇仅文字与公式，无配套几何示意图、原理插图。}
\end{itemize}

% 目录
\newpage
\tableofcontents
\newpage

% 章节内容（矢量分析按课件编为第0章，故章号从 -1 起，使静电场为第1章）
\setcounter{chapter}{-1}
% -------------------------- 矢量分析（教材 附录一 / 课件 第0章） --------------------------
% 说明：本文件为子文件，由主文档 \input 引入，不含 document 环境与重复宏包
% 依赖宏包（主文档 preamble 已装）：tcolorbox、adjustbox、booktabs、amsmath
%
% 彩色公式框：本文件定义一次，供 4.静电场.tex 等同目录子文件复用（勿重复定义）
\newenvironment{formulaBox}{%
    \par\begin{adjustbox}{max width=\linewidth}%
    \begin{tcolorbox}[colback=blue!8!white,colframe=transparent,boxrule=0pt,halign=center,boxsep=2pt,left=4pt,right=4pt,top=3pt,bottom=3pt]%
}{%
    \end{tcolorbox}\end{adjustbox}\par
}

% -------------------------- 正文 --------------------------
\chapter{矢量分析}

矢量分析是研究电磁场理论的重要数学工具。它对简化公式、明确概念、掌握观察规律的实质有极大裨益，
也是学习电磁场入门的数学基础。本章（教材列为附录一，课件称“第 0 章”）只讨论梯度、散度、旋度的概念、物理意义
与直角坐标系下的计算，并给出三种正交坐标系中的常用公式；不追求纯数学论证，侧重阐明物理概念与
应用手段。

\section{正交坐标系}

\subsection{场的描述与坐标系的选择}

电磁场的场量是空间点与时间的函数。为了同时刻画“大小”“方向”和“变化率”，必须选一套坐标基。
当坐标系的三个坐标基矢量互相垂直、且都取单位矢量时，称为\textbf{正交坐标系}，它同时满足
正交性与归一性：
\begin{formulaBox}
\[
\vec{e}_i\cdot\vec{e}_j=\delta_{ij}=
\begin{cases}1, & i=j\\ 0, & i\neq j\end{cases}
\qquad |\vec{e}_i|=1
\]
\end{formulaBox}

题目的边界形状决定坐标系的选择——边界若能与某坐标系的坐标面相重合，问题往往能大幅简化。
这正是后面“分离变量法”能否奏效的前提。

\subsection{三种常用正交坐标系}

\begin{table}[htbp]
\centering
{\small\textbf{表 1-1\quad 三种正交坐标系的坐标、基矢量与适用场域}}\\[2pt]
\small
\begin{tabular}{@{}llll@{}}
\toprule
坐标系 & 坐标 & 单位矢量 & 基矢量是否随位置变化 \\
\midrule
直角（笛卡儿） & $(x,y,z)$ & $\vec{e}_x,\vec{e}_y,\vec{e}_z$ & 常矢量，不随位置变化 \\
圆柱 & $(\rho,\varphi,z)$ & $\vec{e}_\rho,\vec{e}_\varphi,\vec{e}_z$ & $\vec{e}_\rho,\vec{e}_\varphi$ 随 $\varphi$ 变 \\
球 & $(r,\theta,\varphi)$ & $\vec{e}_r,\vec{e}_\theta,\vec{e}_\varphi$ & 三者均随位置变化 \\
\bottomrule
\end{tabular}
\end{table}

名称约定：$\rho$ 为圆柱坐标的径向距离，$\varphi$ 为方位角；$r$ 为球坐标的矢径长度，
$\theta$ 为极角（与 $z$ 轴夹角），$\varphi$ 为方位角。三者取值范围
$\rho\ge 0$、$r\ge 0$、$0\le\theta\le\pi$、$0\le\varphi<2\pi$。

基矢量是否随位置变化，决定了求导与 $\nabla$ 表达式的繁简：直角坐标的三个基矢量方向固定，
故可直接对分量求导，运算最简；圆柱坐标的 $\vec{e}_\rho$、$\vec{e}_\varphi$ 随 $\varphi$ 改变，
是“变矢量”，对它们求导不为零，这正是圆柱坐标下 $\nabla$ 表达式出现 $1/\rho$ 因子的根源；
球坐标的三个基矢量都随位置改变。相应地，直角坐标适用于无限大平面、矩形边界、平板电极；
圆柱坐标适用于同轴电缆、长直导线；球坐标适用于点电荷、同心球壳。

\subsection{坐标变换与单位矢量的变换}

同一场点在不同坐标系中的坐标之间可以互换。设 $x,y,z$ 为直角坐标，则
\begin{formulaBox}
\[
\begin{aligned}
&\text{直角}\to\text{圆柱}:\quad \rho=\sqrt{x^2+y^2},\quad \varphi=\arctan\frac{y}{x},\quad z=z\\
&\text{圆柱}\to\text{直角}:\quad x=\rho\cos\varphi,\quad y=\rho\sin\varphi,\quad z=z\\
&\text{直角}\to\text{球}:\quad r=\sqrt{x^2+y^2+z^2},\quad \theta=\arccos\frac{z}{r},\quad \varphi=\arctan\frac{y}{x}\\
&\text{球}\to\text{直角}:\quad x=r\sin\theta\cos\varphi,\quad y=r\sin\theta\sin\varphi,\quad z=r\cos\theta
\end{aligned}
\]
\end{formulaBox}

反三角函数的象限由 $x$、$y$ 的符号共同确定（只用 $\arctan(y/x)$ 会丢失象限信息，这是坐标变换
最常见的错误）。单位矢量之间的变换为
\begin{formulaBox}
\[
\begin{aligned}
\vec{e}_\rho&=\cos\varphi\,\vec{e}_x+\sin\varphi\,\vec{e}_y, &
\vec{e}_\varphi&=-\sin\varphi\,\vec{e}_x+\cos\varphi\,\vec{e}_y\\[2pt]
\vec{e}_r&=\sin\theta\cos\varphi\,\vec{e}_x+\sin\theta\sin\varphi\,\vec{e}_y+\cos\theta\,\vec{e}_z, &
\vec{e}_\theta&=\cos\theta\cos\varphi\,\vec{e}_x+\cos\theta\sin\varphi\,\vec{e}_y-\sin\theta\,\vec{e}_z
\end{aligned}
\]
\end{formulaBox}

线元、面元、体元是推导管积分的基础，必须记准：
\begin{formulaBox}
\[
\begin{aligned}
\text{直角}:\quad &\mathrm{d}l=\mathrm{d}x\,\vec{e}_x+\mathrm{d}y\,\vec{e}_y+\mathrm{d}z\,\vec{e}_z,\quad \mathrm{d}V=\mathrm{d}x\,\mathrm{d}y\,\mathrm{d}z\\
\text{圆柱}:\quad &\mathrm{d}l=\mathrm{d}\rho\,\vec{e}_\rho+\rho\,\mathrm{d}\varphi\,\vec{e}_\varphi+\mathrm{d}z\,\vec{e}_z,\quad \mathrm{d}V=\rho\,\mathrm{d}\rho\,\mathrm{d}\varphi\,\mathrm{d}z\\
\text{球}:\quad &\mathrm{d}l=\mathrm{d}r\,\vec{e}_r+r\,\mathrm{d}\theta\,\vec{e}_\theta+r\sin\theta\,\mathrm{d}\varphi\,\vec{e}_\varphi,\quad \mathrm{d}V=r^2\sin\theta\,\mathrm{d}r\,\mathrm{d}\theta\,\mathrm{d}\varphi
\end{aligned}
\]
\end{formulaBox}

\subsection{源点、场点与距离矢量}

计算场量时要区分两个点：产生场的电荷所在处称\textbf{源点}，坐标带撇
（$x',y',z'$ 或 $\vec{r}'$）；待求场量的点称\textbf{场点}，坐标不带撇（$x,y,z$ 或 $\vec{r}$）。
两者之间的\textbf{距离矢量}规定由源点指向场点：
\begin{formulaBox}
\[
\vec{R}=\vec{r}-\vec{r}'=R\,\vec{e}_R,\qquad R=|\vec{r}-\vec{r}'|=\sqrt{(x-x')^2+(y-y')^2+(z-z')^2}
\]
\end{formulaBox}

由此导出电磁场中使用频率最高的两个微分关系（$\nabla$ 对场点作用，$\nabla'$ 对源点作用）：
\begin{formulaBox}
\[
\nabla R=\frac{\vec{R}}{R}=\vec{e}_R,\qquad
\nabla'\left(\frac{1}{R}\right)=-\nabla\left(\frac{1}{R}\right)=\frac{\vec{R}}{R^3}=\frac{\vec{e}_R}{R^2}
\]
\end{formulaBox}

\corepoint{注意 $\nabla$ 与 $\nabla'$ 只差一个负号，但作用对象完全不同：前者对场点求导，后者对源点求导。
写积分式时若把 $\nabla$ 提到积分号外，必须确认被积函数中只有源点坐标与积分变量有关。}

\section{矢量代数}

\subsection{标量与矢量}

\textbf{标量}：坐标轴转动时数值不变的物理量，如质量、温度、电位、能量、电阻、电荷量。
\textbf{矢量}：坐标轴转动时按转动矩阵变换的物理量，如位移、速度、力、电场强度、电流密度。
矢量在直角坐标系中写成三个分量与基矢量之积的和，其模为
\begin{formulaBox}
\[
\vec{A}=A_x\vec{e}_x+A_y\vec{e}_y+A_z\vec{e}_z,\qquad
|\vec{A}|=\sqrt{A_x^2+A_y^2+A_z^2},\qquad
\vec{e}_A=\frac{\vec{A}}{|\vec{A}|}
\]
\end{formulaBox}

\subsection{矢量点乘（标量积）}

定义为两矢量模与夹角余弦之积，结果是一个标量：
\begin{formulaBox}
\[
\vec{A}\cdot\vec{B}=|\vec{A}||\vec{B}|\cos(\vec{A},\vec{B})
=A_xB_x+A_yB_y+A_zB_z
\]
\end{formulaBox}

\textbf{理解要点}：点乘等于一个矢量在另一矢量方向上的投影与另一矢量模的乘积——即“$\vec{A}$ 在
$\vec{B}$ 上的投影乘以 $|\vec{B}|$”。因此

\begin{itemize}
    \item $\vec{A}\cdot\vec{B}=0\iff \vec{A}\perp\vec{B}$（正交判据）；
    \item 满足交换律 $\vec{A}\cdot\vec{B}=\vec{B}\cdot\vec{A}$ 与分配律
    $\vec{A}\cdot(\vec{B}+\vec{C})=\vec{A}\cdot\vec{B}+\vec{A}\cdot\vec{C}$；
    \item 物理原型：力做功 $W=\vec{F}\cdot\vec{l}$，电通量 $\Phi=\vec{D}\cdot\vec{S}$。
\end{itemize}

\subsection{矢量叉乘（矢量积）}

结果仍为矢量，大小等于两矢量张成的平行四边形面积，方向由右手螺旋法则确定：
\begin{formulaBox}
\[
\vec{A}\times\vec{B}=
\begin{vmatrix}
\vec{e}_x & \vec{e}_y & \vec{e}_z\\
A_x & A_y & A_z\\
B_x & B_y & B_z
\end{vmatrix}
=|\vec{A}||\vec{B}|\sin(\vec{A},\vec{B})\,\vec{e}_n
\]
\end{formulaBox}

\begin{itemize}
    \item $\vec{A}\times\vec{B}=0\iff \vec{A}\parallel\vec{B}$（平行判据）；
    \item 反交换律 $\vec{A}\times\vec{B}=-\vec{B}\times\vec{A}$，且不满足结合律；
    \item 物理原型：力矩 $\vec{M}=\vec{r}\times\vec{F}$，洛伦兹力 $\vec{F}=q\vec{v}\times\vec{B}$，
    以及“绕行方向与曲面法向成右手螺旋”的有向面积规定。
\end{itemize}

展开行列式时\textbf{$y$ 分量带负号}，这是叉乘最容易出错的一处：
\begin{formulaBox}
\[
\vec{A}\times\vec{B}=(A_yB_z-A_zB_y)\vec{e}_x+(A_zB_x-A_xB_z)\vec{e}_y+(A_xB_y-A_yB_x)\vec{e}_z
\]
\end{formulaBox}

\subsection{混合积与矢量的分解}

三个矢量的混合积 $\vec{A}\cdot(\vec{B}\times\vec{C})$ 是一个标量，其绝对值等于三矢量张成的
平行六面体体积，且满足轮换对称性；二重叉乘则服从“展开法则”：
\begin{formulaBox}
\[
\vec{A}\cdot(\vec{B}\times\vec{C})=\vec{B}\cdot(\vec{C}\times\vec{A})=\vec{C}\cdot(\vec{A}\times\vec{B}),\qquad
\vec{A}\times(\vec{B}\times\vec{C})=(\vec{A}\cdot\vec{C})\vec{B}-(\vec{A}\cdot\vec{B})\vec{C}
\]
\end{formulaBox}

利用正交基，任一矢量都可沿三个互相垂直的方向唯一分解，这正是把矢量方程化为三个标量方程的依据。

\section{标量场与矢量场}

\subsection{场的概念}

分布着某种物理量的空间区域称为该物理量的\textbf{场}。若物理量是标量，称为\textbf{标量场}，
记作 $f(x,y,z;t)$，如温度场、电位场；若是矢量，称为\textbf{矢量场}，记作
$\vec{A}(x,y,z;t)$，如速度场、电场、磁场。不随时间变化的场称\textbf{静态场}，
否则称\textbf{时变场}（动态场）。

\subsection{标量场的等值面与等值线}

等值面是函数值相等的空间曲面，其方程为
\begin{formulaBox}
\[
f(x,y,z)=C
\]
\end{formulaBox}

取不同的 $C$ 得到一族等值面，它们充满整个场域。由于 $f$ 一般是单值函数，\textbf{等值面互不相交}。
电位场中的等位面、温度场中的等温面都是等值面；平行平面场（二维场）中等值面退化为\textbf{等值线}。

等值面的疏密可直接读出场的强弱：相邻等值面间差值相同时，面愈密处场量变化愈剧烈，梯度愈大。

\subsection{矢量场的矢量线}

矢量线是这样的曲线：线上每一点的切线方向与该点矢量场的方向相同，如静电场中的电力线、
流速场中的流线。矢量线充满场域但互不相交（否则该点场方向不唯一）。

因为线元 $\mathrm{d}\vec{l}$ 与 $\vec{A}$ 同向，故矢量线的微分方程（亦称矢量线方程）为
\begin{formulaBox}
\[
\vec{A}\times\mathrm{d}\vec{l}=0
\quad\Longleftrightarrow\quad
\frac{\mathrm{d}x}{A_x}=\frac{\mathrm{d}y}{A_y}=\frac{\mathrm{d}z}{A_z}
\]
\end{formulaBox}

\corepoint{矢量线方程是“成比例”的微分方程组，求解时任取其中一对作常微分方程积分，得到两个积分曲面，
其交线即为矢量线。若给定过某点，用该点坐标定出积分常数。}

作过某曲面 $S$ 上全部点的矢量线，这些线围成的管状区域称\textbf{矢量管}。矢量线只表征方向；
若要同时表征大小，则规定\textbf{穿过单位垂直面积的矢量线条数正比于场的大小}——线愈密处场愈强。

\section{标量场的梯度}

\subsection{方向导数}

标量场 $f(P)$ 从给定点沿某方向 $l$ 的变化率称为\textbf{方向导数}，记作 $\partial f/\partial l$。
设 $l$ 的方向余弦为 $\cos\alpha,\cos\beta,\cos\gamma$，$\vec{e}_l$ 为 $l$ 方向的单位矢量，则
\begin{formulaBox}
\[
\frac{\partial f}{\partial l}=\frac{\partial f}{\partial x}\cos\alpha
+\frac{\partial f}{\partial y}\cos\beta
+\frac{\partial f}{\partial z}\cos\gamma,
\qquad
\vec{e}_l=\vec{e}_x\cos\alpha+\vec{e}_y\cos\beta+\vec{e}_z\cos\gamma
\]
\end{formulaBox}

方向导数大于零表示沿该方向 $f$ 增大，小于零表示减小，等于零表示沿该方向 $f$ 不变
（即沿等值面的切向）。它是一个\textbf{与方向有关}的量：同一点沿不同方向，变化率不同。

\subsection{梯度的定义、大小与方向}

既然从一点出发有无穷多个方向，自然要问：哪个方向变化率最大？最大值是多少？把上式的三个偏导数
记成一个矢量 $\vec{g}$，则方向导数就是 $\vec{g}$ 在 $\vec{e}_l$ 上的投影：
\begin{formulaBox}
\[
\frac{\partial f}{\partial l}=\vec{g}\cdot\vec{e}_l=|\vec{g}|\cos(\vec{g},\vec{e}_l)
\]
\end{formulaBox}

当 $\vec{e}_l$ 与 $\vec{g}$ 同向时方向导数取最大值 $|\vec{g}|$，反向时取最小值 $-|\vec{g}|$，
垂直时为零。于是定义该矢量为标量场的\textbf{梯度}：
\begin{formulaBox}
\[
\operatorname{grad} f=\nabla f
=\frac{\partial f}{\partial x}\vec{e}_x+\frac{\partial f}{\partial y}\vec{e}_y+\frac{\partial f}{\partial z}\vec{e}_z
\]
\end{formulaBox}

\begin{itemize}
    \item 梯度是\textbf{矢量}，且是空间坐标点的函数；其\textbf{方向}为该点最大变化率的方向，
    其\textbf{模}等于该点的最大方向导数。
    \item 已知梯度即可求任一方向的方向导数：$\partial f/\partial l=\nabla f\cdot\vec{e}_l$。
    \item 梯度\textbf{垂直于等值面}，且指向函数值增大的方向；等值面法向单位矢量
    $\vec{e}_n=\pm\nabla f/|\nabla f|$。
    \item $\nabla f=0$ 的点是 $f$ 的驻点；若在场域内处处 $\nabla f=0$，则 $f$ 为常数。
\end{itemize}

在圆柱、球坐标系中，梯度的表达式多出 $1/\rho$、$1/r$ 因子：
\begin{formulaBox}
\[
\nabla f=\frac{\partial f}{\partial\rho}\vec{e}_\rho+\frac{1}{\rho}\frac{\partial f}{\partial\varphi}\vec{e}_\varphi+\frac{\partial f}{\partial z}\vec{e}_z,
\qquad
\nabla f=\frac{\partial f}{\partial r}\vec{e}_r+\frac{1}{r}\frac{\partial f}{\partial\theta}\vec{e}_\theta+\frac{1}{r\sin\theta}\frac{\partial f}{\partial\varphi}\vec{e}_\varphi
\]
\end{formulaBox}

\subsection{哈密顿算子}

为书写与运算方便，引入\textbf{矢量性微分算子}（哈密顿算子，读作 nabla）：
\begin{formulaBox}
\[
\nabla=\vec{e}_x\frac{\partial}{\partial x}+\vec{e}_y\frac{\partial}{\partial y}+\vec{e}_z\frac{\partial}{\partial z}
\]
\end{formulaBox}

它同时具备两种性质：既有\textbf{微分性}（对跟随其后的函数求导），又有\textbf{矢量性}
（可像矢量一样作点乘、叉乘，但不能与之交换顺序）。因此同一符号可以表示三种运算：
\begin{formulaBox}
\[
\begin{aligned}
\nabla f&:\quad \text{梯度（标量场}\to\text{矢量场）}\\
\nabla\cdot\vec{A}&:\quad \text{散度（矢量场}\to\text{标量场）}\\
\nabla\times\vec{A}&:\quad \text{旋度（矢量场}\to\text{矢量场）}
\end{aligned}
\]
\end{formulaBox}

算子作用于不同对象时遵守普通求导法则，例如对两个标量场之积：
\begin{formulaBox}
\[
\nabla(au+bv)=a\nabla u+b\nabla v,\qquad \nabla(uv)=u\nabla v+v\nabla u
\]
\end{formulaBox}

（$a,b$ 为常数。）注意 $\nabla$ 不能写成 $\vec{A}\cdot\nabla$ 与 $\nabla\cdot\vec{A}$ 混用：
前者是标量算子，后者是散度，两者含义不同。

\subsection{拉普拉斯算子}

两个 $\nabla$ 的标积是一个标量算子，称\textbf{拉普拉斯算子}：
\begin{formulaBox}
\[
\nabla^2=\nabla\cdot\nabla=\frac{\partial^2}{\partial x^2}+\frac{\partial^2}{\partial y^2}+\frac{\partial^2}{\partial z^2}
\]
\end{formulaBox}

它作用在标量函数上得标量（拉普拉斯运算），作用在矢量函数上则对三个分量分别运算：
\begin{formulaBox}
\[
\nabla^2 f=\frac{\partial^2 f}{\partial x^2}+\frac{\partial^2 f}{\partial y^2}+\frac{\partial^2 f}{\partial z^2},
\qquad
\nabla^2\vec{A}=(\nabla^2A_x)\vec{e}_x+(\nabla^2A_y)\vec{e}_y+(\nabla^2A_z)\vec{e}_z
\]
\end{formulaBox}

圆柱、球坐标系中的标量拉普拉斯算子为
\begin{formulaBox}
\[
\nabla^2 f=\frac{1}{\rho}\frac{\partial}{\partial\rho}\left(\rho\frac{\partial f}{\partial\rho}\right)
+\frac{1}{\rho^2}\frac{\partial^2 f}{\partial\varphi^2}+\frac{\partial^2 f}{\partial z^2}
\]
\end{formulaBox}
\begin{formulaBox}
\[
\nabla^2 f=\frac{1}{r^2}\frac{\partial}{\partial r}\left(r^2\frac{\partial f}{\partial r}\right)
+\frac{1}{r^2\sin\theta}\frac{\partial}{\partial\theta}\left(\sin\theta\frac{\partial f}{\partial\theta}\right)
+\frac{1}{r^2\sin^2\theta}\frac{\partial^2 f}{\partial\varphi^2}
\]
\end{formulaBox}

\corepoint{圆柱坐标下 $\nabla^2 f$ 的第一项必须写成 $\frac{1}{\rho}\frac{\partial}{\partial\rho}\left(\rho\frac{\partial f}{\partial\rho}\right)$，
不能写成 $\frac{\partial^2 f}{\partial\rho^2}$；球坐标下也不等于 $\frac{\partial^2 f}{\partial r^2}$。
少了 $\rho$、$r^2$ 因子是解泊松/拉普拉斯方程时最常见的错误。}

\section{矢量场的通量与散度}

\subsection{通量}

在矢量场中取一个有向曲面 $S$（$\vec{e}_n$ 为其单位法向），矢量 $\vec{A}$ 在 $S$ 上的面积分
称为 $\vec{A}$ 穿过 $S$ 的\textbf{通量}：
\begin{formulaBox}
\[
\Phi=\iint_S\vec{A}\cdot\mathrm{d}\vec{S}=\iint_S\vec{A}\cdot\vec{e}_n\,\mathrm{d}S
=\iint_S\left(A_x\,\mathrm{d}y\mathrm{d}z+A_y\,\mathrm{d}z\mathrm{d}x+A_z\,\mathrm{d}x\mathrm{d}y\right)
\]
\end{formulaBox}

通量的数学意义是\textbf{穿过有向曲面的矢量线条数}；物理意义随场而定：流速场的通量是单位时间
穿过 $S$ 的体积流量，电流密度 $\vec{J}$ 的通量是电流，电位移 $\vec{D}$ 的通量是自由电荷量。

对\textbf{闭合曲面}，习惯取外法向为正，此时通量表示穿出的净通量，可据其正负判断面内源的性质：
\begin{formulaBox}
\[
\Phi=\oiint_S\vec{A}\cdot\mathrm{d}\vec{S}
=\begin{cases}
>0, & \text{面内有正源（发出矢量线）}\\
<0, & \text{面内有负源（汇集矢量线）}\\
=0, & \text{无源，或正源与负源相抵}
\end{cases}
\]
\end{formulaBox}

通量是积分量，只反映曲面所围区域内的\textbf{整体}情况，说不清面内某一点的性质——因此引出散度。

\subsection{散度的定义与计算}

散度是“通量的体密度”，即单位体积发出的通量：
\begin{formulaBox}
\[
\operatorname{div}\vec{A}=\nabla\cdot\vec{A}
=\lim_{\Delta V\to 0}\frac{1}{\Delta V}\oiint_S\vec{A}\cdot\mathrm{d}\vec{S}
\]
\end{formulaBox}

在直角坐标系中可展开为三个偏导数之和：
\begin{formulaBox}
\[
\nabla\cdot\vec{A}=\frac{\partial A_x}{\partial x}+\frac{\partial A_y}{\partial y}+\frac{\partial A_z}{\partial z}
\]
\end{formulaBox}

柱坐标与球坐标下的散度为
\begin{formulaBox}
\[
\nabla\cdot\vec{A}=\frac{1}{\rho}\frac{\partial(\rho A_\rho)}{\partial\rho}
+\frac{1}{\rho}\frac{\partial A_\varphi}{\partial\varphi}+\frac{\partial A_z}{\partial z}
\]
\end{formulaBox}
\begin{formulaBox}
\[
\nabla\cdot\vec{A}=\frac{1}{r^2}\frac{\partial(r^2A_r)}{\partial r}
+\frac{1}{r\sin\theta}\frac{\partial(\sin\theta A_\theta)}{\partial\theta}
+\frac{1}{r\sin\theta}\frac{\partial A_\varphi}{\partial\varphi}
\]
\end{formulaBox}

\begin{itemize}
    \item 散度是\textbf{标量}，是空间坐标点的函数，描述该点通量源的强度与正负。
    \item $\nabla\cdot\vec{A}>0$：该点有发出通量线的正源；$<0$：有汇集通量线的负源；
    $=0$：该点无源。
    \item 若场中处处 $\nabla\cdot\vec{A}=0$，称\textbf{无源场}（如恒定磁场）。
    \item 静电场中 $\nabla\cdot\vec{D}=\rho_f$：静止电荷就是电场的散度源。
\end{itemize}

\subsection{高斯散度定理}

把散度的定义对体积积分，就得到电磁场理论中使用最广的积分变换定理：
\begin{formulaBox}
\[
\oiint_S\vec{A}\cdot\mathrm{d}\vec{S}=\iiint_V\nabla\cdot\vec{A}\,\mathrm{d}V
\qquad (S \text{ 为包围 } V \text{ 的闭合曲面})
\]
\end{formulaBox}

它把\textbf{闭合曲面上的面积分}与\textbf{体积分}互相转化，是“微分形式 $\Leftrightarrow$ 积分形式”场方程的桥梁：
静电场的高斯定理、恒定磁场的磁通连续性方程、时变场的电荷守恒方程都由它沟通。

\corepoint{应用散度定理前须检查两件事：（1） $S$ 必须是闭合面且取外法向；（2） $\vec{A}$ 在 $V$ 内须有连续的一阶偏导数。
若场域内有奇异点（如点电荷所在处），必须把这些点挖去，用“挖补法”在多连通域上分区应用。}

\section{矢量场的环量与旋度}

\subsection{环量}

矢量场沿有向闭合曲线的线积分称为\textbf{环量}：
\begin{formulaBox}
\[
\Gamma=\oint_L\vec{A}\cdot\mathrm{d}\vec{l}
\]
\end{formulaBox}

环量的物理意义随场而定：力场中的环量是力沿闭合路径做的功；流速场中的环量反映流体是否作涡旋运动。
$\Gamma=0$ 表示沿该回路无涡旋源；$\Gamma\neq 0$ 表示回路包围的区域内必有产生涡旋的源。
但环量和通量一样是积分量，说不清回路内每一点的涡旋强度——因此引出旋度。

\subsection{环量面密度与旋度}

围绕一点可取无穷多个方向的面元 $S$，每个面元对应一个环量面密度 $\Gamma/\Delta S$。其中最大者
的方向与大小就定义了旋度。在直角坐标系中：
\begin{formulaBox}
\[
\operatorname{rot}\vec{A}=\nabla\times\vec{A}=
\begin{vmatrix}
\vec{e}_x & \vec{e}_y & \vec{e}_z\\[2pt]
\dfrac{\partial}{\partial x} & \dfrac{\partial}{\partial y} & \dfrac{\partial}{\partial z}\\[6pt]
A_x & A_y & A_z
\end{vmatrix}
\]
\end{formulaBox}

\begin{itemize}
    \item 旋度是\textbf{矢量}：其\textbf{方向}为使环量面密度取最大值时面元法线的方向，
    其\textbf{模}等于该点环量面密度的最大值。
    \item $\nabla\times\vec{A}\neq 0$：该点矢量线有旋（存在涡旋源）；$=0$：该点无旋。
    若场中处处 $\nabla\times\vec{A}=0$，称\textbf{无旋场}（如静电场）。
    \item 若 $\nabla\times\vec{A}=\vec{J}\neq 0$，称 $\vec{J}$ 为\textbf{旋度源}（涡旋源）。
\end{itemize}

柱坐标与球坐标下的旋度为（注意行列式中的 $\rho$、$r$、$\sin\theta$ 因子）
\begin{formulaBox}
\[
\nabla\times\vec{A}=\frac{1}{\rho}
\begin{vmatrix}
\vec{e}_\rho & \rho\vec{e}_\varphi & \vec{e}_z\\[2pt]
\dfrac{\partial}{\partial\rho} & \dfrac{\partial}{\partial\varphi} & \dfrac{\partial}{\partial z}\\[6pt]
A_\rho & \rho A_\varphi & A_z
\end{vmatrix}
\]
\end{formulaBox}
\begin{formulaBox}
\[
\nabla\times\vec{A}=\frac{1}{r^2\sin\theta}
\begin{vmatrix}
\vec{e}_r & r\vec{e}_\theta & r\sin\theta\,\vec{e}_\varphi\\[2pt]
\dfrac{\partial}{\partial r} & \dfrac{\partial}{\partial\theta} & \dfrac{\partial}{\partial\varphi}\\[6pt]
A_r & rA_\theta & r\sin\theta\,A_\varphi
\end{vmatrix}
\]
\end{formulaBox}

\subsection{斯托克斯定理}

把旋度的定义沿曲面叠加，得到另一个基本积分变换定理：
\begin{formulaBox}
\[
\oint_L\vec{A}\cdot\mathrm{d}\vec{l}=\iint_S(\nabla\times\vec{A})\cdot\mathrm{d}\vec{S}
\]
\end{formulaBox}

其中 $S$ 是以闭合曲线 $L$ 为周界的曲面，$L$ 的绕行方向与 $S$ 的法向成\textbf{右手螺旋}关系。
它把\textbf{闭合曲线上的线积分}与\textbf{曲面积分}互相转化，是推导分界面衔接条件、建立
“环量 $\Leftrightarrow$ 旋度”场方程的工具。

\corepoint{斯托克斯定理中 $L$ 与 $S$ 的方向必须成右手螺旋；若算出的环量符号与预期相反，八成是绕行方向
或法向取反了。散度定理与斯托克斯定理合起来，构成“微分场方程 $\Leftrightarrow$ 积分场方程”的完整通道。}

\section{无源场和无旋场}

\subsection{无源场}

若在场域中每一点恒有 $\nabla\cdot\vec{A}=0$，称 $\vec{A}$ 为\textbf{无源场}。它有两个重要性质。

\textbf{性质 1（矢量管通量守恒）}：无源场中穿过任一个矢量管的所有截面的通量都相等。
证明要点是取矢量管的一段，其上、下端面 $S_1$、$S_2$ 与侧面 $S_3$ 构成闭合面，由散度定理
$\oiint\vec{A}\cdot\mathrm{d}\vec{S}=0$，而侧面上 $\vec{A}\cdot\vec{e}_n=0$，故穿过 $S_1$ 与
$S_2$ 的通量大小相等、符号相反。在磁场中这一性质称为\textbf{磁通连续性原理}。

\textbf{性质 2（存在矢势）}：任一矢量场旋度再取散度恒为零，即
$\nabla\cdot(\nabla\times\vec{F})\equiv 0$，所以旋度场一定是无源场；反之，无源场 $\vec{A}$ 必存在
另一矢量场 $\vec{F}$ 使
\begin{formulaBox}
\[
\vec{A}=\nabla\times\vec{F}
\]
\end{formulaBox}

$\vec{F}$ 称为 $\vec{A}$ 的\textbf{矢势}（矢量位）。例如恒定磁场中 $\nabla\cdot\vec{B}=0$，
故 $\vec{B}=\nabla\times\vec{A}$，$\vec{A}$ 即磁矢位。

\subsection{无旋场}

若在场域中每一点恒有 $\nabla\times\vec{A}=0$，称 $\vec{A}$ 为\textbf{无旋场}。

\textbf{性质 1}：无旋场沿任意闭合路径的环量等于零，$\oint_L\vec{A}\cdot\mathrm{d}\vec{l}=0$。
因此无旋场中场力做功与路径无关，故又称\textbf{保守场}。静电场 $\oint\vec{E}\cdot\mathrm{d}\vec{l}=0$ 即属此类。

\textbf{性质 2}：无旋场可表示为某一标量函数的梯度场。因为任一标量场的梯度再取旋度恒为零，
即 $\nabla\times(\nabla\varphi)\equiv 0$，所以梯度场一定是无旋场；反之无旋场 $\vec{A}$ 可写成
\begin{formulaBox}
\[
\vec{A}=-\nabla\varphi
\]
\end{formulaBox}

这样的场称\textbf{位势场}，$\varphi$ 称 $\vec{A}$ 的\textbf{位函数}。静电场中 $\varphi$ 就是电位，
满足 $\vec{E}=-\nabla\varphi$。（式中负号来自物理约定：电场指向电位降低的方向。）

\subsection{调和场}

散度与旋度都等于零的矢量场称\textbf{调和场}，它是既无源又无旋的场。设 $\vec{A}=-\nabla\varphi$
且 $\nabla\cdot\vec{A}=0$，则
\begin{formulaBox}
\[
\nabla\cdot(-\nabla\varphi)=-\nabla^2\varphi=0
\quad\Longrightarrow\quad
\nabla^2\varphi=0
\]
\end{formulaBox}

即调和场的位函数必满足\textbf{拉普拉斯方程}；满足拉普拉斯方程且具有二阶连续偏导数的函数称
\textbf{调和函数}。若矢量场仅无旋而不无源，其位函数满足\textbf{泊松方程}
$\nabla^2\varphi=-\rho/\varepsilon_0$。这正是静电场边值问题的方程来源。

\section{亥姆霍兹定理}

\textbf{定理}：在空间有限区域 $V$ 内的任一矢量场 $\vec{A}$，由它的\textbf{散度}、\textbf{旋度}
和\textbf{边界条件}（即包围 $V$ 的闭合面 $S$ 上矢量场的分布）唯一地确定。

等价表述：给定了矢量场的通量源密度 $\nabla\cdot\vec{A}$、涡旋源密度 $\nabla\times\vec{A}$ 以及
场域的边界条件，就可唯一确定该矢量场。

\begin{itemize}
    \item 该定理说明：一个矢量场可以分解为无旋部分与无散部分之和，
    \begin{formulaBox}
    \[
    \vec{A}=-\nabla\varphi+\nabla\times\vec{F}
    \]
    \end{formulaBox}
    前者由散度源（通量源）产生，后者由旋度源（涡旋源）产生。
    \item 它是电磁场理论中“场的唯一性”与“场的分解”的基础：静电场由 $\nabla\cdot\vec{D}=\rho_f$
    与 $\nabla\times\vec{E}=0$ 唯一确定；恒定磁场由 $\nabla\times\vec{H}=\vec{J}$ 与
    $\nabla\cdot\vec{B}=0$ 唯一确定；时变场则由完整的麦克斯韦方程组确定。
    \item 由此也看出为什么静电场可以用标量电位描述（无旋，只剩散度源），而恒定磁场必须用矢量位
    描述（无散，只剩旋度源）。
\end{itemize}

\corepoint{本章五组关系必须烂熟：$\nabla\varphi$（梯度，标量$\to$矢量）、$\nabla\cdot\vec{A}$（散度，矢量$\to$标量）、
$\nabla\times\vec{A}$（旋度，矢量$\to$矢量）、$\nabla^2$（拉普拉斯，标量$\to$标量），以及两条恒等式
$\nabla\times(\nabla\varphi)\equiv 0$、$\nabla\cdot(\nabla\times\vec{A})\equiv 0$。
后面每一章的场方程都由这几个算子写成。}

% 静电场（教材 第1章 §1-1～§1-9）- 子文件（仅用于 \input 引入）
% 依赖：formulaBox 环境在 3.矢量分析.tex 中定义，本文件不再重复定义
% 编号：章号由主文件 \chapter 自动给出，节内不再手写 §1-x

% -------------------------- 正文 --------------------------
\chapter{静电场}

相对于观察者静止、电荷量不随时间变化的电荷所产生的电场称为静电场。本章由库仑定律引入电场强度
$\vec{E}$ 与标量电位 $\varphi$，讨论导体与电介质的影响，导出衔接条件与泊松、拉普拉斯方程，把问题
归结为边值问题，再介绍直接积分、分离变量、有限差分、镜像法与电轴法，最后讨论电容、静电能量与力。

\section{电场强度、电位}

\subsection{库仑定律与电场强度}

无限大真空中，两个静止小带电体间距离远大于其几何尺寸时，作用力为
\begin{formulaBox}
\[
\vec{F}_{12}=\frac{1}{4\pi\varepsilon_0}\frac{q_1q_2}{R^2}\vec{e}_R,
\qquad
\varepsilon_0=\frac{10^{-9}}{36\pi}\ \mathrm{F/m}\approx 8.85\times10^{-12}\ \mathrm{F/m}
\]
\end{formulaBox}

适用条件：带电体几何尺寸远小于其间距离，即视为\textbf{点电荷}（“点”只是相对意义）。电场强度定义
为试验电荷所受的力与其电荷量之比的极限
\begin{formulaBox}
\[
\vec{E}=\lim_{q_0\to0}\frac{\vec{F}}{q_0}\ (\mathrm{V/m}),
\qquad
\vec{E}(\vec{r})=\frac{1}{4\pi\varepsilon_0}\frac{q}{|\vec{r}-\vec{r}'|^3}\left(\vec{r}-\vec{r}'\right)
\]
\end{formulaBox}

$\vec{E}$ 仅由该点电场决定，\textbf{与试验电荷无关}。场与源成线性关系，故连续分布电荷的场强按
叠加原理写成矢量积分：
\begin{formulaBox}
\[
\vec{E}=\frac{1}{4\pi\varepsilon_0}\int\frac{\vec{R}}{R^3}\,\mathrm{d}q,
\qquad
\mathrm{d}q=\rho\,\mathrm{d}V'\ \text{或}\ \sigma\,\mathrm{d}S'\ \text{或}\ \tau\,\mathrm{d}l'
\]
\end{formulaBox}

面电荷、线电荷是体电荷的极限形式，体电荷公式通用。电荷分布具有平面对称、柱对称、球对称时，
应改用电场的高斯定律而不做硬性矢量积分；积分的实用技巧是“配对消去”——利用对称性成对取元电荷，
使垂直合成方向的分量相消。

\subsection{电位与电压}

单位正电荷由 $A$ 移至 $B$，电场力作功与路径无关；沿闭合路径作功恒为零：
\begin{formulaBox}
\[
\oint_L\vec{E}\cdot\mathrm{d}\vec{l}=0
\quad\Longleftrightarrow\quad
\nabla\times\vec{E}=0
\]
\end{formulaBox}

前者为\textbf{环路定律}，后者为其微分形式（由斯托克斯定理得到），二者等价，说明静电场是
\textbf{无旋场}、即\textbf{保守场}。由“梯度的旋度恒为零”，无旋的 $\vec{E}$ 可表为标量电位的负梯度：
\begin{formulaBox}
\[
\vec{E}=-\nabla\varphi,
\qquad
\int_A^B\vec{E}\cdot\mathrm{d}\vec{l}=\varphi_A-\varphi_B=U_{AB}
\]
\end{formulaBox}

负号表示 $\vec{E}$ 指向电位\textbf{最大减小率}方向。$\varphi$ 与 $\varphi+C$ 代表同一 $\vec{E}$，
故电位是相对的，须选参考点 $Q$（$\varphi_Q=0$）：
\begin{formulaBox}
\[
\varphi_P=\int_P^{Q}\vec{E}\cdot\mathrm{d}\vec{l}
\]
\end{formulaBox}

工程上常取大地为参考点；电荷全在有限空间内时取无穷远为零电位最方便。\textbf{电荷分布延拓到无限远}
时（无限长线电荷、无限大带电平面）无穷远处电位为无限大，必须改在有限空间内选参考点。取无穷远为
零电位时
\begin{formulaBox}
\[
\varphi(\vec{r})=\frac{1}{4\pi\varepsilon_0}\frac{q}{R},
\qquad
\varphi(\vec{r})=\frac{1}{4\pi\varepsilon_0}\int\frac{\mathrm{d}q}{R}
\]
\end{formulaBox}

一对相距很近的等量异号电荷构成\textbf{电偶极子}，电偶极矩 $\vec{p}=q\vec{l}$（由负指向正，
$\mathrm{C\cdot m}$）。$r\gg l$ 处
\begin{formulaBox}
\[
\varphi=\frac{1}{4\pi\varepsilon_0}\frac{\vec{p}\cdot\vec{e}_r}{r^2},
\qquad
\vec{E}=\frac{p}{4\pi\varepsilon_0r^3}\left(2\cos\theta\,\vec{e}_r+\sin\theta\,\vec{e}_\theta\right)
\]
\end{formulaBox}

即电位与 $r^2$ 成反比、场强与 $r^3$ 成反比，衰减比点电荷快，是分析介质极化与电磁辐射的基础模型。

\subsection{电场线与等位面}

$E$ 线上每点切线方向与该点场强方向一致，其微分方程为
\begin{formulaBox}
\[
\vec{E}\times\mathrm{d}\vec{l}=0
\quad\Longleftrightarrow\quad
\frac{\mathrm{d}x}{E_x}=\frac{\mathrm{d}y}{E_y}=\frac{\mathrm{d}z}{E_z}
\]
\end{formulaBox}

等位面方程为 $\varphi(x,y,z)=C$，它与平面的截迹为等位线。由 $\vec{E}=-\nabla\varphi$ 得：等位面与
$E$ 线\textbf{处处正交}；$E$ 线指向\textbf{电位降低}方向；相邻等位面电位差相等时，
\textbf{等位面愈密处场强愈大}。

\corepoint{等位面上电位处处相同，但场强一般\textbf{不}相同（场强取决于电位的变化率）。故“电位为零处场强
必为零”“电位高处场强必大”“场强相等处电位必相等”均为错误判断。}

\section{高斯定律}

\subsection{静电场中的导体}

导体引入外电场后，自由电荷移动、积聚于表面，其附加场与外场在导体内处处抵消而达静电平衡。此时：
（1） 导体内 $\vec{E}=0$；（2） 导体为\textbf{等位体}、表面为\textbf{等位面}；（3） 表面 $\vec{E}$
\textbf{垂直于}表面；（4） 电荷只能分布在导体\textbf{表面}。

\subsection{电介质的极化}

介质中电子被束缚，外场作用下正负电荷作用中心不再重合，形成小电偶极子，称\textbf{介质极化}。
取体积元 $\Delta V$，定义\textbf{电极化强度}
\begin{formulaBox}
\[
\vec{P}=\lim_{\Delta V\to0}\frac{\sum\vec{p}_i}{\Delta V}\ (\mathrm{C/m^2}),
\qquad
\rho_p=-\nabla\cdot\vec{P},
\qquad
\sigma_p=\vec{P}\cdot\vec{e}_n,
\qquad
\vec{P}=\chi_e\varepsilon_0\vec{E}
\]
\end{formulaBox}

$\rho_p$、$\sigma_p$ 为极化电荷的体密度与面密度（$\vec{e}_n$ 为介质表面外法向），$\chi_e$ 为电极化率。

\corepoint{均匀介质极化时体极化电荷为零（$\vec{P}$ 均匀则 $\nabla\cdot\vec{P}=0$），只在表面出现面极化电荷；
不均匀介质即使无自由电荷也可能有体极化电荷。判断“何处有极化电荷”看 $\vec{P}$ 是否随位置变化，
而不是看介质是否被极化。}

\subsection{电通密度与高斯定律}

把极化电荷用 $\vec{P}$ 表示后引入\textbf{电通（量）密度}（电位移矢量）：
\begin{formulaBox}
\[
\vec{D}=\varepsilon_0\vec{E}+\vec{P},
\qquad
\oint_S\vec{D}\cdot\mathrm{d}\vec{S}=q_f,
\qquad
\nabla\cdot\vec{D}=\rho_f
\]
\end{formulaBox}

即任意闭合面上 $\vec{D}$ 的面积分等于面内\textbf{总自由电荷}，与极化电荷及面外电荷无关——引入
$\vec{D}$ 的便利正在于此。对各向同性线性介质，代入 $\vec{P}=\chi_e\varepsilon_0\vec{E}$ 得构成方程
\begin{formulaBox}
\[
\vec{D}=\varepsilon_0(1+\chi_e)\vec{E}=\varepsilon\vec{E}=\varepsilon_r\varepsilon_0\vec{E}
\]
\end{formulaBox}

\subsection{用高斯定律计算电场}

场分布具有平面对称、柱对称或球对称时，用高斯定律求场很直接。步骤：（1） 由电荷分布对称性判断
$\vec{D}$ 的方向与等值面（球对称$\to$径向，柱对称$\to$径向，面对称$\to$法向）；（2） 过待求点作合适的高斯面
（球面、同轴圆柱面加两端面、扁圆柱面），使面上 $\vec{D}$ 与法向平行且等值、或与法向垂直；
（3） 由 $\oint_S\vec{D}\cdot\mathrm{d}\vec{S}=q_f$ 解出 $D$，再由 $\vec{D}=\varepsilon\vec{E}$ 得
$\vec{E}$，必要时积分求 $\varphi$。

\corepoint{高斯定律永远成立，但只有对称性足够时才可用来“解出”场强。对称性不足（有限长线段、偏心圆柱）只能
回到叠加积分或解边值问题；多层介质按分界面分段，各段用高斯定律，再用衔接条件与电压约束联立。}

\section{静电场基本方程、分界面上的衔接条件}

\subsection{静电场基本方程}

静电场是\textbf{无旋场}，同时又是\textbf{有散场}——静止电荷是它的散度源：
\begin{formulaBox}
\[
\begin{aligned}
\text{积分形式：}\quad &\oint_S\vec{D}\cdot\mathrm{d}\vec{S}=q_f,
\qquad \oint_L\vec{E}\cdot\mathrm{d}\vec{l}=0\\
\text{微分形式：}\quad &\nabla\cdot\vec{D}=\rho_f,
\qquad \nabla\times\vec{E}=0\\
\text{构成方程：}\quad &\vec{D}=\varepsilon\vec{E}
\qquad(\text{各向同性线性介质})
\end{aligned}
\]
\end{formulaBox}

积分形式刻画每条回路、每个闭合面上场量的整体情况，微分形式刻画各点及其邻域，便于分析计算。

\subsection{分界面上的衔接条件}

分界面两侧物性突变，场量可能突变，基本方程的\textbf{微分形式失效}，须回到\textbf{积分形式}推导。
设 $\vec{e}_n$ 由媒质 1 指向媒质 2。

（1）围绕界面点作高 $\Delta h\to0$ 的小扁圆柱体（两端面 $\Delta S$ 分居两侧），用高斯定律，侧面通量
趋零，得电通密度的法向条件（$\sigma_f$ 为界面自由电荷面密度）：
\begin{formulaBox}
\[
D_{2n}-D_{1n}=\sigma_f
\qquad\Longleftrightarrow\qquad
\vec{e}_n\cdot\left(\vec{D}_2-\vec{D}_1\right)=\sigma_f
\]
\end{formulaBox}

（2）围绕界面点作与界面垂直的狭小矩形环路（垂直边长 $\Delta h\to0$），用环路定律，得电场强度的
切向条件：
\begin{formulaBox}
\[
E_{1t}=E_{2t}
\qquad\Longleftrightarrow\qquad
\vec{e}_n\times\left(\vec{E}_2-\vec{E}_1\right)=0
\]
\end{formulaBox}

\subsection{两种特殊分界面与电位的衔接条件}

（1）两种\textbf{理想电介质}分界面（$\sigma_f=0$）：$D_{1n}=D_{2n}$、$E_{1t}=E_{2t}$。设场与法线
夹角为 $\alpha_1,\alpha_2$，由 $D_n=\varepsilon E\cos\alpha$、$E_t=E\sin\alpha$ 相除得\textbf{折射定律}
\begin{formulaBox}
\[
\frac{\tan\alpha_1}{\tan\alpha_2}=\frac{\varepsilon_1}{\varepsilon_2}
\]
\end{formulaBox}

即 $\varepsilon$ 大的一侧场与法线夹角大，电场线“折向”介电常数大的一侧。

（2）\textbf{导体与电介质}分界面：静电平衡时导体内 $\vec{E}_1=0$、$\vec{D}_1=0$，故
\begin{formulaBox}
\[
E_{2t}=0,\qquad D_{2n}=\sigma_f,\qquad E_2=\frac{\sigma_f}{\varepsilon_2}
\]
\end{formulaBox}

介质侧电场只有法向分量，垂直于导体表面，大小由表面自由电荷面密度决定。

（3）\textbf{用电位表示}：在界面两侧取相距 $d\to0$ 的两点，由
$\varphi_A-\varphi_B=\int_A^B\vec{E}\cdot\mathrm{d}\vec{l}$ 得
\begin{formulaBox}
\[
\varphi_1=\varphi_2,
\qquad
\varepsilon_1\frac{\partial\varphi_1}{\partial n}-\varepsilon_2\frac{\partial\varphi_2}{\partial n}=\sigma_f
\]
\end{formulaBox}

即界面两侧\textbf{电位连续}（与 $E_{1t}=E_{2t}$ 等价），而电位法向导数有跃变。

\corepoint{衔接条件必须与 $\vec{e}_n$ 的指向配套：规定 $\vec{e}_n$ 由媒质 1 指向媒质 2，才是
$D_{2n}-D_{1n}=\sigma_f$；方向取反则右端变号，全题符号皆错。角度 $\alpha$ 一律与\textbf{法线}起算。}

\section{静电场边值问题、唯一性定理}

\subsection{泊松方程与拉普拉斯方程}

把 $\vec{D}=\varepsilon\vec{E}$、$\vec{E}=-\nabla\varphi$ 代入 $\nabla\cdot\vec{D}=\rho_f$，对均匀介质
（$\varepsilon$ 为常数）得
\begin{formulaBox}
\[
\nabla\cdot(\varepsilon\nabla\varphi)=-\rho_f
\quad\Longrightarrow\quad
\nabla^2\varphi=-\frac{\rho_f}{\varepsilon},
\qquad
\nabla^2\varphi=0\ (\rho_f=0)
\]
\end{formulaBox}

前者为\textbf{泊松方程}，后者为\textbf{拉普拉斯方程}。$\nabla^2$ 在直角、圆柱、球坐标系中的展开式
见“矢量分析”部分（拉普拉斯算子一节），解题时按场域对称性选用相应形式。

\subsection{边值问题与三类边界条件}

只靠微分方程不能唯一确定解（积分过程中出现待定常数），须由边界条件定。故
\begin{formulaBox}
\[
\text{边值问题}=\text{泛定方程（泊松或拉普拉斯方程）}+\text{定解条件（边界条件）}
\]
\end{formulaBox}

\begin{itemize}
    \item \textbf{第一类（狄里赫利问题）}：已知边界电位值 $\varphi|_S=g_1(\vec{r})$；
    \item \textbf{第二类（纽曼问题）}：已知电位法向导数值
    $\left.\dfrac{\partial\varphi}{\partial n}\right|_S=g_2(\vec{r})$；给定导体总电荷量也属此类
    （由 $\oint_S\varepsilon\dfrac{\partial\varphi}{\partial n}\mathrm{d}S=-q$ 给出积分形式的条件）；
    \item \textbf{第三类（混合问题）}：一部分边界给定 $\varphi$，另一部分给定 $\partial\varphi/\partial n$。
\end{itemize}

场域伸展到无限远时须补充\textbf{无穷远边界条件}：电荷分布在有限区域时
$\lim_{r\to\infty}\varphi=$ 有限值，称\textbf{自然边界条件}。若介质分区均匀，则各子区域分别列方程，
并\textbf{引入分界面衔接条件}作为定解条件的一部分。

\subsection{唯一性定理}

\textbf{定理}：凡满足下列条件的电位函数是唯一解——（1） 在场域 $V$ 中满足电位的微分方程（分区均匀时
满足各子区域的方程）；（2） 在不同介质分界面上符合衔接条件；（3） 在场域边界面 $S$ 上满足给定的边界条件。
简述为：\textbf{满足电位微分方程和给定边界条件的解，就是唯一解}。

\textbf{证明要点（反证法）}：设 $\varphi_1,\varphi_2$ 都满足泊松方程，则 $u=\varphi_1-\varphi_2$
满足 $\nabla^2u=0$。对 $u$ 用高斯散度定理与 $\nabla\cdot(u\nabla u)=|\nabla u|^2+u\nabla^2u$，得
\begin{formulaBox}
\[
\oint_S u\frac{\partial u}{\partial n}\,\mathrm{d}S=\int_V|\nabla u|^2\,\mathrm{d}V
\]
\end{formulaBox}

第一类边界下全边界 $u=0$，第二类下全边界 $\partial u/\partial n=0$，两种情形左端均为零，故
$\int_V|\nabla u|^2\mathrm{d}V=0$；被积函数非负，则 $V$ 内 $\nabla u=0$、$u=$ 常数，再由边界条件或
参考点定该常数为零，于是 $\varphi_1=\varphi_2$。

唯一性定理的意义：它给出唯一解的条件，并可判定所得解是否正确——\textbf{只要能找到满足全部条件的解，
它就是唯一正确的解}。镜像法、电轴法、分离变量法都在利用这一许可。

\corepoint{唯一性定理是解题许可而非理论摆设：凡“凑”出的解，只需验过 （1） 场内方程、（2） 分界面衔接条件、
（3） 边界条件三关即为唯一解；反之，只满足边界条件而不满足场内方程的解一定不是解。}

\section{分离变量法}

\subsection{基本思想与直角坐标下的解}

一维问题可直接积分；二维、三维问题不能。\textbf{当场域边界面能与某正交坐标系的坐标面重合时}，
分离变量法是有效解法：把电位写成单变量函数的\textbf{乘积}，代入偏微分方程，借“分离常数”化为几个
常微分方程，解之并由边界条件定常数；所得解常为傅里叶级数形式，故又称\textbf{傅里叶法}。

设电位只是 $x,y$ 的函数（沿 $z$ 无变化），令 $\varphi=X(x)Y(y)$，代入
$\partial^2\varphi/\partial x^2+\partial^2\varphi/\partial y^2=0$ 并除以 $XY$：
\begin{formulaBox}
\[
\frac{1}{X}\frac{\mathrm{d}^2X}{\mathrm{d}x^2}+\frac{1}{Y}\frac{\mathrm{d}^2Y}{\mathrm{d}y^2}=0
\quad\Longrightarrow\quad
\frac{\mathrm{d}^2X}{\mathrm{d}x^2}+k^2X=0,\quad
\frac{\mathrm{d}^2Y}{\mathrm{d}y^2}-k^2Y=0
\]
\end{formulaBox}

左端两项分属 $x$、$y$ 而恒等，故必各等于同一常数，记作 $-k^2$，$k$ 称\textbf{分离常数}。$k=0$ 时
解为线性函数；$k\neq0$ 时为三角函数与双曲函数。由线性叠加得一般解
\begin{formulaBox}
\[
\varphi=(A_0x+B_0)(C_0y+D_0)
+\sum_{n}\left(A_n\cos k_nx+B_n\sin k_nx\right)
\left(C_n\,\mathrm{ch}\,k_ny+D_n\,\mathrm{sh}\,k_ny\right)
\]
\end{formulaBox}

即 $x$ 的周期函数、$y$ 的双曲函数；分离常数换号则得另一个一般解（$x$ 双曲、$y$ 周期）。
\textbf{如何选取分离常数由具体边界条件决定}，这是本法的关键。

\subsection{圆柱坐标系中的解}

设电位沿 $z$ 无变化，则
$\rho\frac{\partial}{\partial\rho}\left(\rho\frac{\partial\varphi}{\partial\rho}\right)
+\frac{\partial^2\varphi}{\partial\varphi^2}=0$。令 $\varphi=R(\rho)Q(\varphi)$，分离得
\begin{formulaBox}
\[
n=0:\quad R_0=A_0\ln\rho+B_0,
\qquad
Q_0=C_0\varphi+D_0
\]
\end{formulaBox}
\begin{formulaBox}
\[
n\neq0:\quad R_n=A_n\rho^{n}+B_n\rho^{-n},
\qquad
Q_n=C_n\cos n\varphi+D_n\sin n\varphi
\]
\end{formulaBox}

相应乘积叠加即为一般解。$\ln\rho$ 与 $\rho^{-n}$ 在 $\rho\to0$ 时发散，故圆柱内部（含轴线）的解中
不能保留这两项——用“有限性”筛项。

\subsection{求解步骤}

（1） 按场域形状选坐标系，使边界面与坐标面吻合，写出边值问题；（2） 分离变量化为常微分方程；（3） 解之并
组成通解（含分离常数与待定常数）；（4） 由边界条件定常数得唯一解。

\corepoint{成败在第 （1） 步：坐标系选错则边界无法与坐标面重合，级数不能逐项比较系数。定系数时用“边界上处处
为零 $\Rightarrow$ 该项系数为零”与三角函数正交性（乘 $\sin$ 或 $\cos$ 后积分）；判形式：有两条平行边界（周期）
的方向用三角函数，延伸到无穷远的方向用衰减项。}

\section{有限差分法}

\subsection{基本思想与差分格式}

边界形状复杂时一般只能求近似解。有限差分法把\textbf{场域用网格离散化}，把拉普拉斯方程用以各节点
电位为未知数的\textbf{差分方程}代替，把偏微分方程问题化为联立差分方程组问题。

用两组平行于坐标轴的直线把二维场域划成足够多的\textbf{正方形网格}，交点称节点，相邻平行网格线
间距称\textbf{步距} $h$。用中心差商代替二阶偏导，得拉普拉斯方程的\textbf{差分格式}（五点格式）
\begin{formulaBox}
\[
\frac{\partial^2\varphi}{\partial x^2}\approx\frac{\varphi_{i+1,j}-2\varphi_{i,j}+\varphi_{i-1,j}}{h^2}
\quad\Longrightarrow\quad
\varphi_{i,j}=\frac{1}{4}\left(\varphi_{i+1,j}+\varphi_{i-1,j}+\varphi_{i,j+1}+\varphi_{i,j-1}\right)
\]
\end{formulaBox}

含义直观：\textbf{任一节点电位近似等于相邻四点电位的平均值}，对场域内每一节点均可列出一个差分方程。
若边界不落在节点上（紧邻边界的节点），须按不等步距 $ph$、$qh$ 处理，差分格式中四邻项系数不再相等。

\subsection{差分方程组的解法}

节点数可达几百上千，直接解法不适用，但每方程只含少数几项，可用迭代法。

（1）\textbf{高斯—塞德尔迭代法}：先选初值，再按
\begin{formulaBox}
\[
\varphi^{(k+1)}_{i,j}=\frac{1}{4}\left(
\varphi^{(k)}_{i+1,j}+\varphi^{(k+1)}_{i-1,j}
+\varphi^{(k)}_{i,j+1}+\varphi^{(k+1)}_{i,j-1}\right)
\]
\end{formulaBox}
反复迭代（本轮新值立即代入，比雅可比法快一倍）。遇边界点须代入边界值。迭代到对所有内节点满足
$\max\limits_{i,j}\left|\varphi^{(k+1)}_{i,j}-\varphi^{(k)}_{i,j}\right|<W$ 为止，$W$ 为预定最大
允许误差；节点按“自然顺序”（先左到右、再下到上）排列。

（2）\textbf{逐次超松弛法}：把每次迭代的变化量加权后再代入以加速收敛
\begin{formulaBox}
\[
\varphi^{(k+1)}_{i,j}=(1-\alpha)\varphi^{(k)}_{i,j}
+\frac{\alpha}{4}\left(
\varphi^{(k)}_{i+1,j}+\varphi^{(k+1)}_{i-1,j}
+\varphi^{(k)}_{i,j+1}+\varphi^{(k+1)}_{i,j-1}\right),
\qquad 0<\alpha<2
\]
\end{formulaBox}
$\alpha$ 称\textbf{加速收敛因子}：$\alpha=1$ 退化为高斯—塞德尔法，$\alpha>1$ 为超松弛。选得好可
明显加快收敛，选得不好会振荡甚至发散，最佳因子需经验或试算确定。

\corepoint{有限差分法优点是网格灵活、适应复杂边界、程序简单；缺点是只给节点上的数值解，精度受步距限制
（步距减半则节点数增为四倍、计算量剧增），边界附近用不等步距格式时误差较大。其迭代格式与实验课
“静电场有限差分法编程实现”一致。}

\section{镜像法和电轴法}

\subsection{原理与适用区域}

两种方法的\textbf{实质}是：把实际不均匀媒质看成均匀媒质，而把分界面上的电荷分布（导体表面感应
电荷、介质界面极化电荷）用虚设的\textbf{镜像电荷}代替。根据唯一性定理，只要虚设电荷与实际电荷
共同产生的电场能满足给定边界条件，结果就正确。

\corepoint{镜像法两条铁律：（1） 镜像电荷必须放在\textbf{待求场域之外}（否则引入新奇异点）；（2） 镜像解\textbf{只在
原场域内}有效，边界另一侧的场不由这组电荷产生。违反任一条，答案必错。}

\subsection{点电荷与无限大接地导体平面}

接地无限大导体平面上方 $d$ 处有点电荷 $q$，上半空间须满足：除 $q$ 处外 $\nabla^2\varphi=0$；导体
平面及无穷远处 $\varphi=0$。在平面下方对称点放 $-q$、撤去导体板、全空间充满 $\varepsilon_0$ 介质即可：
\begin{formulaBox}
\[
\varphi=\frac{q}{4\pi\varepsilon_0}\left(\frac{1}{R_1}-\frac{1}{R_2}\right),
\qquad
R_1=\sqrt{x^2+y^2+(z-d)^2},\quad R_2=\sqrt{x^2+y^2+(z+d)^2}
\]
\end{formulaBox}

$-q$ 即 $q$ 对导体平面的镜像，它代替了平面感应电荷的作用；\textbf{适用区域为平面上半空间}，下半
空间实际无电场。由 $\sigma_f=D_n=\varepsilon_0E_n$ 得感应电荷面密度
$\sigma_f=-qd/\left[2\pi(x^2+y^2+d^2)^{3/2}\right]$，总量为 $-q$。

\subsection{点电荷与导体球}

半径 $R$ 的接地导体球外、距球心 $d$ 处有点电荷 $q$。镜像电荷 $-q'$ 放在球心与 $q$ 的联线上、
距球心 $b$ 处（球内），由球面电位为零可定
\begin{formulaBox}
\[
b=\frac{R^2}{d},
\qquad
q'=\frac{R}{d}q,
\qquad
\varphi=\frac{1}{4\pi\varepsilon_0}\left(\frac{q}{R_1}-\frac{q'}{R_2}\right)
\]
\end{formulaBox}

球面感应电荷总量为 $-Rq/d$。对\textbf{不接地、不带电}的孤立导体球，球面为未知常数电位且总电荷为
零，须再在\textbf{球心}放 $+q'=Rq/d$，使其与球内镜像电荷总量为零。若为导体球壳内的点电荷，可反用
结果（像电荷在球外，适用区域为球腔）。

\subsection{点电荷与两种介质的分界平面}

界面两侧介电常数为 $\varepsilon_1,\varepsilon_2$，左半空间有点电荷 $q$。镜像系统：左半空间的场由
$q$ 与像点处的 $q'$ 共同产生（$\varepsilon_1$ 充满全空间）；右半空间的场由位于原电荷处的 $q''$
单独产生（$\varepsilon_2$ 充满全空间）。由界面衔接条件解得
\begin{formulaBox}
\[
q'=\frac{\varepsilon_1-\varepsilon_2}{\varepsilon_1+\varepsilon_2}q,
\qquad
q''=\frac{2\varepsilon_2}{\varepsilon_1+\varepsilon_2}q
\]
\end{formulaBox}

$q'$ 在界面对称点，适用区域为含 $q$ 的一侧；$q''$ 在原电荷处，适用另一侧。若
$\varepsilon_1=\varepsilon_2$ 则 $q'=0,q''=q$；若 $\varepsilon_2\to\infty$（相当于导体面）则
$q'\to-q,q''\to2q$。

\subsection{电轴法}

两平行长直圆柱导体表面电荷分布不均匀且未知，已知的只是单位长度总电荷 $\pm\tau$。设想撤去圆柱、
以两根相距 $2b$、线密度 $\pm\tau$ 的长直带电细线代替其表面电荷效应，这两根细线所在轴线即
\textbf{电轴}。两圆柱外部任一点的电位为
\begin{formulaBox}
\[
\varphi=\frac{\tau}{2\pi\varepsilon_0}\ln\frac{\rho_2}{\rho_1}+C
\]
\end{formulaBox}

$\rho_1,\rho_2$ 为场点到 $+\tau$、$-\tau$ 的距离；参考点取在对称轴上时 $C=0$。由 $\rho_2/\rho_1=K$
（常数）得等位线，为\textbf{一族偏心圆}：圆心 $x_0=\frac{K^2+1}{K^2-1}b$、半径
$R=\left|\frac{2bK}{K^2-1}\right|$，且半径 $R$、圆心到原点距离 $d$、电轴到原点距离 $b$ 满足
\begin{formulaBox}
\[
R^2=d^2-b^2
\quad\Longrightarrow\quad
b^2=h^2-a^2
\]
\end{formulaBox}

后者是把圆柱半径 $a$、轴心距 $h$ 代入后的形式：由它确定电轴位置，代入电位式即得外部电位分布。
对已知两圆柱间电压 $U_0$ 的情况，易得 $\tau$ 与 $U_0$ 的关系及单位长度电容：
\begin{formulaBox}
\[
U_0=\frac{\tau}{\pi\varepsilon_0}\ln\frac{h+b}{a},
\qquad
C_0=\frac{\tau}{U_0}=\frac{\pi\varepsilon_0}{\ln\left[(h+b)/a\right]}
\]
\end{formulaBox}

对两根不同半径、轴线距离为 $d$ 的非对称传输线，对每个圆柱分别写出 $b^2=h_1^2-a_1^2$、
$b^2=h_2^2-a_2^2$，再结合 $h_1+h_2=d$ 联立解出电轴位置。

\corepoint{电轴法只适用于\textbf{带等量异号电荷的两平行圆柱导体}（同号带电、三根以上导线不适用）。顺序：
由 $b^2=h^2-a^2$ 定电轴位置 $\to$ 写 $\varphi=\frac{\tau}{2\pi\varepsilon_0}\ln\frac{\rho_2}{\rho_1}$
$\to$ 由电压或电荷定 $\tau$ $\to$ 求 $C_0$。}

\section{电容和部分电容}

\subsection{电容及其计算}

两带等量异号电荷的导体组成电容器，其电容定义为电荷与导体间电压之比：
\begin{formulaBox}
\[
C=\frac{q}{U}
=\frac{\oint_S\vec{D}\cdot\mathrm{d}\vec{S}}{\int_A^B\vec{E}\cdot\mathrm{d}\vec{l}}
=\frac{\oint_S\varepsilon\vec{E}\cdot\mathrm{d}\vec{S}}{\int_A^B\vec{E}\cdot\mathrm{d}\vec{l}}
\]
\end{formulaBox}

单位 $\mathrm{F}$（工程常用 $\mathrm{\mu F}$、$\mathrm{pF}$）。电容只与两导体的形状、尺寸、相互位置
及其间介质有关，\textbf{与电荷量、电压无关}；孤立导体的电容指它与无限远处的电容。电容计算实质就是
静电场计算：
\begin{formulaBox}
\[
\text{设电荷 }q\ \xrightarrow{\oint_S\vec{D}\cdot\mathrm{d}\vec{S}}\ \vec{D},\vec{E}
\ \xrightarrow{\int_A^B\vec{E}\cdot\mathrm{d}\vec{l}}\ U
\ \xrightarrow{C=q/U}\ C
\]
\end{formulaBox}

三个必背结果（介质为 $\varepsilon$）：
\begin{formulaBox}
\[
C_{\text{平行板}}=\frac{\varepsilon S}{d},
\qquad
C_{0(\text{同轴圆柱})}=\frac{2\pi\varepsilon}{\ln(b/a)},
\qquad
C_{\text{同心球}}=\frac{4\pi\varepsilon ab}{b-a}
\]
\end{formulaBox}

同轴圆柱式为\textbf{单位长度}电容；同心球式当 $b\to\infty$ 时趋于 $4\pi\varepsilon a$，即\textbf{孤立
导体球}的电容。多层介质可按分界面分段、用衔接条件联立，或按串并联处理。

\subsection{静电独立系统与电位系数}

若系统内场的分布只与系统内带电体的形状、尺寸、相互位置及介质分布有关、与外部带电体无关，且所有
电通密度线都从系统内带电体发出并终止于系统内带电体，则称\textbf{静电独立系统}。$n+1$ 个导体的
静电独立系统必有
\begin{formulaBox}
\[
q_0+q_1+\cdots+q_n=0
\]
\end{formulaBox}

线性介质下由叠加原理得电位与电荷的关系（取 0 号导体为参考点，$\varphi_0=0$）：
\begin{formulaBox}
\[
\varphi_k=\sum_{i=1}^{n}\alpha_{ki}q_i
\qquad\Longleftrightarrow\qquad
[\varphi]=[\alpha][q]
\]
\end{formulaBox}

$\alpha_{kk}$ 为\textbf{自有电位系数}，$\alpha_{ki}$（$k\neq i$）为\textbf{互有电位系数}，性质：
（1） 全为正值；（2） 自有系数大于与之有关的互有系数；（3） 只与几何形状、尺寸、相互位置及介电常数有关；
（4） $\alpha_{ki}=\alpha_{ik}$（$[\alpha]$ 对称，是静电场互易原理的表现）。反解得以电位表示电荷的
\textbf{感应系数}形式 $q_k=\sum_i\beta_{ki}\varphi_i$：自有感应系数为\textbf{正}、互有感应系数为
\textbf{负}，且自有系数大于与之有关的互有系数的绝对值。

\subsection{部分电容与电容网络}

把上式改写为用该导体与其他各导体间电压表示，得
\begin{formulaBox}
\[
q_k=C_{k1}U_{k1}+C_{k2}U_{k2}+\cdots+C_{k0}U_{k0}+\cdots+C_{kn}U_{kn}
\]
\end{formulaBox}

系数 $C$ 称\textbf{部分电容}：$C_{k0}$ 为\textbf{自有部分电容}，$C_{ki}$ 为\textbf{互有部分电容}，
定义为
\begin{formulaBox}
\[
C_{k0}=\left.\frac{q_k}{U_{k0}}\right|_{\text{其余电压为零}},
\qquad
C_{ki}=\left.\frac{q_k}{U_{ki}}\right|_{\text{其余电压为零}}
\]
\end{formulaBox}

性质：（1） 全为\textbf{正值}；（2） 只与导体几何、介质有关，\textbf{与电荷量无关}；（3） 互易性
$C_{ki}=C_{ik}$。$n+1$ 个导体的静电独立系统中共有 $\frac{n(n+1)}{2}$ 个部分电容，它们构成一个
\textbf{电容网络}，从而把静电场问题变为电容电路问题，把“场”与“路”联系起来。

\corepoint{部分电容与等效电容含义不同：部分电容是网络支路参数，只与结构有关；等效电容是把两导体当电容器的两
极板、加压后极板电荷与电压之比，还包含其他导体的影响（如二线传输线的等效电容受大地影响）。已知两个
等效电容的测量值时，须借对称性与网络方程反解各支路部分电容。}

\subsection{静电屏蔽}

三导体系统中 1 号导体被接地的 0 号导体完全包围。由部分电容方程，0 号导体内无电场、$\varphi_{10}=0$；
因 $U_{12}$ 可取任意值（导体 2 的电荷可任意），必有 $C_{12}=0$，于是
\begin{formulaBox}
\[
q_1=C_{10}U_{10},
\qquad
q_2=C_{20}U_{20}
\]
\end{formulaBox}

即 $q_1$ 只与 $U_{10}$ 有关、$q_2$ 只与 $U_{20}$ 有关，导体 1、2 之间\textbf{无静电联系}，达到
静电屏蔽。工程上把不宜受外界电场影响（或不宜影响外界）的带电体用接地金属壳罩起来，如高压设备周围
的屏蔽网。

\section{静电能量与力}

\subsection{带电体系统中的静电能量}

静电能量是在电场建立过程中由\textbf{外力作功}转化而来的。设某一瞬时场中某点电位为 $\varphi$，
把电荷增量 $\delta q$ 从无穷远移至该点，外力作功 $\delta A=\varphi\,\delta q$。对线性介质，达到最终
分布所需作的功与过程无关，故可取各处电荷密度按同一比例 $m$（$0\le m\le1$）同步增长的方式，把
$m$ 从 0 积到 1，得
\begin{formulaBox}
\[
W_e=\frac{1}{2}\int_{V'}\rho\varphi\,\mathrm{d}V'
=\frac{1}{2}\int_{S'}\sigma\varphi\,\mathrm{d}S'
=\frac{1}{2}\sum_{k=1}^{n}q_k\varphi_k
\]
\end{formulaBox}

第三式适用于只有带电导体的系统。系数 $1/2$ 来自充电过程中电荷与电位同时按比例增长（积两次）。
对点电荷系统，能量可分为\textbf{自有能量}（只与各电荷自身有关，电荷相互移近时不变；点电荷的自有
能量趋于无限大）与\textbf{互有能量}（由电荷间相互作用引起，是可测量的那部分）。

\subsection{能量分布与能量密度}

上述公式只给总能量，易误解为“能量集中在电荷上”。实际上静电能量分布在整个电场空间中。应用
$\vec{D}=\varepsilon\vec{E}$、$\vec{E}=-\nabla\varphi$、矢量恒等式与散度定理，并把 $V$ 扩展到整个
无限空间（$S$ 为 $r\to\infty$ 的球面，$\varphi\sim1/r$、$\vec{D}\sim1/r^2$、$\mathrm{d}S\sim r^2$，
面积分项随 $1/r$ 趋于零），得
\begin{formulaBox}
\[
W_e=\frac{1}{2}\int_{V_\infty}\vec{D}\cdot\vec{E}\,\mathrm{d}V,
\qquad
w_e=\frac{1}{2}\vec{D}\cdot\vec{E}=\frac{1}{2}\varepsilon E^2\ (\mathrm{J/m^3})
\]
\end{formulaBox}

即\textbf{凡静电场不为零的空间都储存着静电能量，场是能量的载体}；$w_e$ 为\textbf{静电能量密度}。
两式均为体积分但积分区域不同：$\frac12\int\rho\varphi\mathrm{d}V$ 为\textbf{电荷积分式}，区域是有
电荷分布处；$\frac12\int\vec{D}\cdot\vec{E}\mathrm{d}V$ 为\textbf{电场积分式}，区域是整个空间。
只有对全部电场空间积分，两式才相等。

\corepoint{能量密度是\textbf{场}的属性，不是“某点电荷的能量密度”。用 $\frac12\int\rho\varphi\mathrm{d}V$ 时
若只对部分区域（如只对某导体积分）会得到错误结果——这正是必须分清两式积分区域的原因。}

\subsection{虚位移法计算静电力}

小带电体受力可按定义 $\vec{F}=q\vec{E}$ 计算，其中 $\vec{E}$ 应理解为\textbf{除 $q$ 以外}其他电荷
在 $q$ 处引起的场强；对连续分布电荷这样算相当复杂，常用基于虚功原理的\textbf{虚位移法}。

\textbf{广义坐标}是确定各导体形状、尺寸与位置的一组独立几何量（距离、面积、体积、角度），企图改变
某一广义坐标的力即对应的\textbf{广义力}。设只有一个广义坐标 $g$ 变化，功能过程为
\begin{formulaBox}
\[
\mathrm{d}W_{\text{源}}=\mathrm{d}W_e+f\,\mathrm{d}g
\]
\end{formulaBox}

（1）\textbf{各带电体总电荷不变}（不与外源相联，$\mathrm{d}q_k=0$）：$\mathrm{d}W_{\text{源}}=0$，
$f\mathrm{d}g=-\mathrm{d}W_e$，即电场力作功靠\textbf{减少}静电能量实现：
\begin{formulaBox}
\[
f=-\left.\frac{\partial W_e}{\partial g}\right|_{q=\text{常量}}
\]
\end{formulaBox}

（2）\textbf{各带电体电位不变}（均接电源）：由 $W_e=\frac12\sum q_k\varphi_k$ 得
$\mathrm{d}W_e=\frac12\mathrm{d}W_{\text{源}}$，故 $f\mathrm{d}g=\mathrm{d}W_e$，即外源提供的能量
一半用于增加静电能量、一半用于电场力作功：
\begin{formulaBox}
\[
f=+\left.\frac{\partial W_e}{\partial g}\right|_{\varphi=\text{常量}}
\]
\end{formulaBox}

两种假设所得结果\textbf{应相等}——带电体实际未动（虚位移），电场力分布未变，求得的是当时的电荷与
电位下的力。以 $W_e=\frac{q^2}{2C}$ 与 $\frac12CU^2$ 分别代入均得
$f=\frac{q^2}{2C^2}\frac{\partial C}{\partial g}$，可见\textbf{电场力总有使电容 $C$ 增大的趋势}。
对板面积 $S$、间距 $d$、介电常数 $\varepsilon$ 的平行板电容器（取 $d$ 为广义坐标）：
\begin{formulaBox}
\[
f=-\frac{q^2}{2\varepsilon S}=-\frac{\varepsilon SU^2}{2d^2},
\qquad
p=\frac{f}{S}=\frac{\sigma^2}{2\varepsilon}=\frac{1}{2}\sigma E
\]
\end{formulaBox}

负号表示电场力有使 $d$ 缩短的趋势（两极板相吸）。

\corepoint{这里的系数 $1/2$ 是最大陷阱：极板受力只由\textbf{另一极板（其余电荷）}产生的场决定，而那部分场只占
总场强的一半，故 $p=\frac12\sigma E$ 而非 $\sigma E$。球形带电薄膜的膨胀力同法得
$p=\sigma^2/(2\varepsilon_0)$。}

\subsection{法拉第对静电力的观点}

法拉第认为：静电场中每一段电通密度管沿轴向受\textbf{纵张力}、垂直于轴向受\textbf{侧压力}，两者量值
相等，都是
\begin{formulaBox}
\[
\frac{1}{2}D\cdot E=\frac{1}{2}DE\qquad(\mathrm{N/m^2})
\]
\end{formulaBox}

形象地说，电位移管像被拉紧的橡皮筋：沿轴向有收缩趋势、垂直方向有扩张趋势。据此可由场图判断受力：
孤立带电体表面有扩张趋势；带等量异号电荷的两导体相吸；带等量同号电荷的两导体相斥；不带电导体有被
吸向带电体的趋势。

对两种介质的分界面，用极短的电通密度管取极限，得单位面积受力
\begin{formulaBox}
\[
\text{场垂直于界面：}\ p=\frac{1}{2}D^2\left|\frac{1}{\varepsilon_1}-\frac{1}{\varepsilon_2}\right|,
\qquad
\text{场平行于界面：}\ p=\frac{1}{2}E^2\left|\varepsilon_1-\varepsilon_2\right|
\]
\end{formulaBox}

结论（常考）：无论场与界面垂直还是平行，作用在界面上的力总\textbf{垂直于界面}，且总是\textbf{由
介电常数大的一边指向小的一边}。

\corepoint{本章主线：§1-1 定义场量（$\vec{E},\varphi$）$\to$ §1-2 引入介质与 $\vec{D}$、得散度方程 $\to$ §1-3 凑齐
基本方程与衔接条件 $\to$ §1-4 化为边值问题并证唯一性 $\to$ §1-5～§1-7 三种直接解法与两种间接解法 $\to$
§1-8～§1-9 由场求路参数（$C$）与力学量（$W_e,f$）。大题多落在“衔接条件＋边值问题＋镜像法＋电容/
能量”这条链上，动手前先判断题目属于哪一环。}

% ==================== 第2章 恒定电场 ====================
% 子文件：仅正文，由 1.电磁场笔记.tex \input 引入
% 教材：《工程电磁场导论》第2章 §2-1 ~ §2-5
% 公式统一用 formulaBox（本目录 3.矢量分析.tex 中定义）
% ==========================================================

\chapter{恒定电场}

\textbf{研究对象}：导电媒质中的恒定电场，即导电媒质内有恒定电流通过时，媒质内、外建立的电场。

\textbf{基本特征}：电荷分布不随时间变化，处处保持动态平衡；场量（$\vec{E}$、$\vec{J}$、$\rho$）均与时间无关，
\begin{formulaBox}
\[
\frac{\partial}{\partial t}(\bullet)=0
\]
\end{formulaBox}
这一条件把恒定电场与静电场、时变电磁场区分开：静电场由\emph{静止}电荷产生、场域内无电流；恒定电场由\emph{恒定运动}的电荷（恒定电流）维持，导电媒质内部有电流通过。因此恒定电场中\textbf{导体不再是等位体}，导体内部场强不为零（$E=J/\gamma \neq 0$），这是它与静电场最容易被混淆的区别。

本章按教材节序展开：导电媒质中的电流（\S2-1）$\to$ 电源电动势与局外场强（\S2-2）$\to$ 基本方程与分界面衔接条件（\S2-3）$\to$ 恒定电场与静电场的比拟（\S2-4）$\to$ 电导和部分电导（\S2-5）。

% ---------------------------------------------------------
\section{导电媒质中的电流}

\subsection{传导电流与运流电流}
按电流形成的物理机制分类：
\begin{itemize}
  \item \textbf{传导电流}：导电媒质（导体、电解液、大地等）中自由电荷的定向运动形成，服从欧姆定律；
  \item \textbf{运流电流}：自由空间（真空、气体）中带电体的机械运动形成，\textbf{不服从欧姆定律}。
\end{itemize}
本章只讨论传导电流。描述电流有两个层次的物理量：电流强度（整体量）与电流密度（分布量）。

\subsection{电流强度与电流密度矢量}
\textbf{电流强度}定义为单位时间内通过某一曲面的电量，
\begin{formulaBox}
\[
I=\frac{\mathrm{d}q}{\mathrm{d}t}\qquad(\mathrm{A})
\]
\end{formulaBox}
它是\textbf{通量}概念，只说明每秒通过某面积的总电量，无法反映电荷在截面各点的流动情况。

\textbf{电流密度矢量} $\vec{J}$ 描述电流的空间分布：方向取正电荷运动方向；大小为单位时间内穿过与该方向垂直的单位面积的净电量。对以速度 $\vec{v}$ 运动的体电荷分布，
\begin{formulaBox}
\[
\vec{J}=\rho\vec{v}\qquad(\mathrm{A/m^2})
\]
\end{formulaBox}
其中 $\rho$ 为电荷体密度。穿过任意曲面 $S$ 的电流是 $\vec{J}$ 的通量：
\begin{formulaBox}
\[
I=\iint_S \vec{J}\cdot\mathrm{d}\vec{S}
\]
\end{formulaBox}
按电荷分布形态，电流分为三类（$\vec{J}$、$\vec{K}$、$I$ 分别对应体、面、线分布）：
\begin{center}
\small
\begin{tabularx}{\linewidth}{@{}l l X@{}}
\toprule
\textbf{分布} & \textbf{密度} & \textbf{与电流强度的关系} \\
\midrule
体电流 & $\vec{J}=\rho\vec{v}$\ ($\mathrm{A/m^2}$) & $I=\displaystyle\iint_S \vec{J}\cdot\mathrm{d}\vec{S}$ \\
面电流 & $\vec{K}=\sigma\vec{v}$\ ($\mathrm{A/m}$) & $I=\displaystyle\int_l \vec{K}\cdot\vec{e}_n\,\mathrm{d}l$（$\vec{e}_n$ 为面法向） \\
线电流 & $I=\tau v$（$\tau$ 为单位长度电荷） & 沿导线流动，$I$ 即电流强度 \\
\bottomrule
\end{tabularx}
\end{center}

\corepoint{三层关系不要写反：$\vec{J}$ 是\emph{点}函数（$\mathrm{A/m^2}$），$\vec{K}$ 是\emph{面}上的密度（$\mathrm{A/m}$，沿面法向线元积分），$I$ 是\emph{积分}结果（$\mathrm{A}$）。凡是题目给"电流线密度 $K$"，就要用 $\int K\,\mathrm{d}l$ 而不是 $\iint J\,\mathrm{d}S$。}

\subsection{欧姆定律的微分形式}
在线性、各向同性导电媒质中，任一点传导电流密度与该点电场强度成正比：
\begin{formulaBox}
\[
\vec{J}=\gamma\vec{E}
\]
\end{formulaBox}
\begin{itemize}
  \item $\gamma$ 为媒质的\textbf{电导率}（$\mathrm{S/m}$），$\rho_r=1/\gamma$ 为电阻率（$\Omega\cdot\mathrm{m}$）；
  \item 适用条件：线性、各向同性导电媒质，且为\emph{点}的方程，与导体形状无关；
  \item 物理意义：导电媒质的物性方程，地位与电介质的 $\vec{D}=\varepsilon\vec{E}$、磁介质的 $\vec{B}=\mu\vec{H}$ 相当。
\end{itemize}
对一段均匀导体积分即得电路形式的欧姆定律 $U=IR$、$I=GU$，故 $U=IR$ 只是微分形式的宏观平均结果。

\subsection{焦耳定律的微分形式}
单位体积内的功率损耗（热功率体密度）为
\begin{formulaBox}
\[
p=\lim_{\Delta V\to 0}\frac{\Delta P}{\Delta V}=\vec{J}\cdot\vec{E}
\qquad(\mathrm{W/m^3})
\]
\end{formulaBox}
在线性、各向同性媒质中，代入 $\vec{J}=\gamma\vec{E}$ 得常用的三种等价写法：
\begin{formulaBox}
\[
p=\gamma E^2=\frac{J^2}{\gamma}
\]
\end{formulaBox}
对整个导体体积分得总功率损耗
\begin{formulaBox}
\[
P=\int_V p\,\mathrm{d}V=\int_V \vec{E}\cdot\vec{J}\,\mathrm{d}V=I^2R
\]
\end{formulaBox}
\textbf{物理意义}：恒定电场中电流的流动必然伴随能量损耗，要维持恒定电流必须持续由电源补充能量；$p$ 是\emph{点}的损耗，$P=I^2R$ 是整体结果，两者不能混用。

% ---------------------------------------------------------
\section{电源电动势与局外场强}

\subsection{局外力与局外场强}
要维持恒定电流，必须有非静电力把正电荷从低电位处"搬"回高电位处，这种力称为\textbf{局外力} $\vec{f}_e$（不是电磁力）。把作用于单位正电荷的局外力等效为一个场强，称为\textbf{局外场强}：
\begin{formulaBox}
\[
\vec{E}_e=\lim_{q\to 0}\frac{\vec{f}_e}{q}
\]
\end{formulaBox}
\begin{itemize}
  \item 方向：由电源负极指向正极（电源内部）；
  \item 它\textbf{不是}真正的电场强度，只在"引起电流"这一点上与电场强度等效；它不满足保守场条件；
  \item 电源外部不存在局外场强，只有库仑场强 $\vec{E}_c$。
\end{itemize}

\subsection{含源导电媒质中的电流与电动势}
电源内部总场强为 $\vec{E}_c+\vec{E}_e$，故通过含源导电媒质的电流密度为
\begin{formulaBox}
\[
\vec{J}=\gamma(\vec{E}_e+\vec{E}_c)
\]
\end{formulaBox}
\textbf{电动势}衡量电源做功能力，定义为单位正电荷从负极移到正极时局外力所做的功：
\begin{formulaBox}
\[
\varepsilon=\int_{B}^{A}\vec{E}_e\cdot\mathrm{d}\vec{l}
\]
\end{formulaBox}
（$B$ 为负极、$A$ 为正极，积分只在电源内部进行。）

\subsection{电源外恒定电场的守恒性}
对不经过电源的任意闭合路径 $L$，
\begin{formulaBox}
\[
\oint_L \vec{E}\cdot\mathrm{d}\vec{l}=\oint_L \vec{E}_c\cdot\mathrm{d}\vec{l}=0
\]
\end{formulaBox}
即\textbf{电源外恒定电场仍是无旋（保守）场}，可以引入标量电位 $\vec{E}=-\nabla\varphi$。

\corepoint{$\oint_L\vec{E}_e\cdot\mathrm{d}\vec{l}\neq 0$（局外场非保守），但 $\oint_L\vec{E}_c\cdot\mathrm{d}\vec{l}=0$。讨论"恒定电场是否无旋"时，指的是库仑场强 $\vec{E}_c$，且积分路径\emph{不穿过电源}；若路径穿过电源，总场强的环路积分等于电动势而不是零。}

% ---------------------------------------------------------
\section{恒定电场的基本方程与分界面上的衔接条件}

\subsection{电荷守恒定律与电流连续性方程}
单位时间流出闭合面 $S$ 的电荷等于体积 $V$ 内电量的减少率：
\begin{formulaBox}
\[
\oint_S \vec{J}\cdot\mathrm{d}\vec{S}=-\frac{\mathrm{d}q}{\mathrm{d}t}
\]
\end{formulaBox}
由高斯散度定理得\textbf{电流连续性方程}（微分形式）
\begin{formulaBox}
\[
\nabla\cdot\vec{J}+\frac{\partial\rho}{\partial t}=0
\]
\end{formulaBox}
恒定条件下 $\partial\rho/\partial t=0$，退化为
\begin{formulaBox}
\[
\nabla\cdot\vec{J}=0,\qquad \oint_S \vec{J}\cdot\mathrm{d}\vec{S}=0
\]
\end{formulaBox}
\textbf{物理意义}：恒定电流"无头无尾"，电流线是闭合曲线；任一闭合面净流出的传导电流为零（流入等于流出）。

\subsection{恒定电场的基本方程组}
电源外导电媒质中的恒定电场由下列方程完整描述（三种形式并列，缺一不可）：
\begin{formulaBox}
\[
\begin{cases}
\nabla\times\vec{E}=0 & \text{（无旋）}\\[2pt]
\nabla\cdot\vec{J}=0 & \text{（无源）}\\[2pt]
\vec{J}=\gamma\vec{E} & \text{（物性方程）}
\end{cases}
\qquad\Longleftrightarrow\qquad
\begin{cases}
\displaystyle\oint_L \vec{E}\cdot\mathrm{d}\vec{l}=0\\[6pt]
\displaystyle\oint_S \vec{J}\cdot\mathrm{d}\vec{S}=0
\end{cases}
\]
\end{formulaBox}
\begin{itemize}
  \item $\nabla\times\vec{E}=0$：恒定电场是无旋场，可引入电位 $\vec{E}=-\nabla\varphi$；
  \item $\nabla\cdot\vec{J}=0$：恒定电场是无源场\emph{（就电流线而言）}，但注意 $\nabla\cdot\vec{D}=\rho$ 仍可能成立（见下款）；
  \item 辅助方程（有电介质时）：$\nabla\cdot\vec{D}=\rho$，$\vec{D}=\varepsilon\vec{E}$，与静电场相同。
\end{itemize}

\subsection{分界面上的衔接条件}
设 $\vec{e}_n$ 为分界面上由媒质 1 指向媒质 2 的单位法向，由积分形式的基本方程可导出四个衔接条件（前两个是\emph{场量}形式，后两个是\emph{电位}形式）：
\begin{formulaBox}
\[
\begin{cases}
E_{1t}=E_{2t}\ \ \text{或}\ \ \vec{e}_n\times(\vec{E}_2-\vec{E}_1)=0 & \text{（切向：}E_t\text{连续）}\\[4pt]
J_{1n}=J_{2n}\ \ \text{或}\ \ \vec{e}_n\cdot(\vec{J}_2-\vec{J}_1)=0 & \text{（法向：}J_n\text{连续）}\\[4pt]
D_{2n}-D_{1n}=\sigma\ \ \text{或}\ \ \vec{e}_n\cdot(\vec{D}_2-\vec{D}_1)=\sigma & \text{（法向：}D_n\text{跃变）}\\[4pt]
\varphi_1=\varphi_2,\quad \gamma_1\dfrac{\partial\varphi_1}{\partial n}=\gamma_2\dfrac{\partial\varphi_2}{\partial n} & \text{（电位形式）}
\end{cases}
\]
\end{formulaBox}
\textbf{条件与含义}：$E_t$ 连续来自 $\oint\vec{E}\cdot\mathrm{d}\vec{l}=0$（沿跨界面小矩形回路）；$J_n$ 连续来自 $\oint\vec{J}\cdot\mathrm{d}\vec{S}=0$（跨界面小扁盒）；$D_n$ 的跃变量恰为分界面上的自由面电荷密度 $\sigma$。注意 $J_n$ 连续\textbf{不意味着} $D_n$ 连续——因为两侧的 $\varepsilon/\gamma$ 不同。

综合 $E_{1t}=E_{2t}$、$J_{1n}=J_{2n}$ 与 $\vec{J}=\gamma\vec{E}$，得\textbf{恒定电场的折射定律}：
\begin{formulaBox}
\[
\frac{\tan\theta_1}{\tan\theta_2}=\frac{\gamma_1}{\gamma_2}
\]
\end{formulaBox}
其中 $\theta_1,\theta_2$ 分别为两侧电流密度（或电场强度）与界面法线的夹角。

\subsection{两类特殊分界面的讨论}
\begin{enumerate}
  \item \textbf{良导体—不良导体}（$\gamma_2\gg\gamma_1$）：由折射定律 $\tan\theta_1=(\gamma_1/\gamma_2)\tan\theta_2\to 0$，故 $\theta_1\approx0$。结论：从不良导体一侧看，电流线（及 $\vec{E}$ 线）近似\emph{垂直于}分界面；良导体一侧的界面可近似当作\textbf{等位面}。
  \item \textbf{导体—理想介质}（如架空线与空气，$\gamma_2=0$）：由 $J_{2n}=J_{1n}$ 得 $J_{2n}=0$，即介质中无电流，导体侧电流只能沿表面流动（$J_1=J_{1t}$）；同时
  \[
  D_{2n}-D_{1n}=\sigma \Rightarrow D_{2n}=\sigma,\qquad E_{1t}=E_{2t},\ E_1=E_{1t}
  \]
  结论：分界面上\textbf{必有面电荷分布}；电场切向分量不为零，介质中紧贴导体表面的 $\vec{E}$ \emph{不再}垂直于导体表面。但通常 $E_{2n}\gg E_{2t}$，故工程上仍近似认为导体表面外的场垂直于表面。
\end{enumerate}

\corepoint{本节最易错的三处：(1) 把"$J_n$ 连续"写成"$D_n$ 连续"；(2) 把"$E_t$ 连续"与"$D_t$ 连续"混淆（$D_t$ 一般不连续）；(3) 直接套用静电场"导体表面场垂直于表面"的结论——在恒定电场中它只是近似。}

\subsection{导电媒质中恒定电荷的分布}
\begin{enumerate}
  \item \textbf{均匀媒质}（$\gamma$、$\varepsilon$ 均为常数）：由 $\nabla\cdot\vec{J}=0$ 得
  \begin{formulaBox}
  \[
  \nabla\cdot(\gamma\vec{E})=\nabla\cdot\Big(\frac{\gamma}{\varepsilon}\vec{D}\Big)=0
  \ \Longrightarrow\ \rho=0
  \]
  \end{formulaBox}
  即均匀导电媒质内部\textbf{无体电荷}，电荷只能分布在分界面上或电极表面。
  \item \textbf{不均匀媒质}（$\gamma$ 随位置变化）：
  \begin{formulaBox}
  \[
  \rho=\nabla\cdot\vec{D}=\nabla\cdot\Big(\frac{\varepsilon}{\gamma}\vec{J}\Big)
  =\vec{J}\cdot\nabla\Big(\frac{\varepsilon}{\gamma}\Big)
  \]
  \end{formulaBox}
  即媒质内部\textbf{可以存在体电荷}，其分布取决于 $\varepsilon/\gamma$ 的梯度方向与电流方向的关系。
  \item \textbf{分界面上的面电荷}：
  \begin{formulaBox}
  \[
  \sigma=D_{2n}-D_{1n}=\varepsilon_2E_{2n}-\varepsilon_1E_{1n}
  =\Big(\frac{\varepsilon_2}{\gamma_2}-\frac{\varepsilon_1}{\gamma_1}\Big)J_{1n}
  \]
  \end{formulaBox}
  对两种\emph{金属}导体的分界面，经典电子论给出各金属的 $\varepsilon$ 近似相同，故 $\sigma\approx0$；而对金属—介质、不同电解液的分界面一般 $\sigma\neq0$。
\end{enumerate}

\subsection{恒定电场的边值问题}
由 $\nabla\times\vec{E}=0$ 引入 $\vec{E}=-\nabla\varphi$，代入 $\nabla\cdot(\gamma\vec{E})=0$。均匀媒质中 $\gamma$ 为常数，得
\begin{formulaBox}
\[
\nabla^2\varphi=0
\]
\end{formulaBox}
即恒定电场的电位满足\textbf{拉普拉斯方程}，配合衔接条件 $\varphi_1=\varphi_2$、$\gamma_1\partial\varphi_1/\partial n=\gamma_2\partial\varphi_2/\partial n$ 与场域边界条件（给定电位的第一类边界、给定法向电流密度的第二类边界），构成恒定电场的\textbf{边值问题}。它与静电场边值问题同型，故解法（分离变量、镜像、有限差分）完全可借用——这正是下一节"比拟"的依据。

% ---------------------------------------------------------
\section{导电媒质中的恒定电场与静电场的比拟}

\subsection{比拟的依据与条件}
比较\emph{电源外}导电媒质中的恒定电场与\emph{无电荷分布区}的静电场：
\[
\text{静电场：}\ \nabla\times\vec{E}=0,\ \nabla\cdot\vec{D}=0
\qquad
\text{恒定电场：}\ \nabla\times\vec{E}=0,\ \nabla\cdot\vec{J}=0
\]
两者的\emph{无旋}方程完全相同，\emph{无源}方程形式相同，本构方程 $\vec{D}=\varepsilon\vec{E}$ 与 $\vec{J}=\gamma\vec{E}$ 一一对应，衔接条件也一一对应。因此，只要两个场的\textbf{边界条件相同}，其电位函数 $\varphi$ 的解答就完全相同，场图（等位面、场线）一致。

\textbf{适用条件}（缺一不可）：
\begin{enumerate}
  \item 恒定电场须取\emph{电源外}区域（否则 $\oint\vec{E}\cdot\mathrm{d}\vec{l}\neq0$，无旋性不成立）；
  \item 静电场须取\emph{无自由电荷}区域（$\rho=0$），否则 $\nabla\cdot\vec{D}\neq0$；
  \item 若媒质分区不均匀，还须满足 $\varepsilon_1/\varepsilon_2=\gamma_1/\gamma_2$（即各分区内 $\varepsilon/\gamma$ 为同一常数），否则分界面折射情况不同，比拟不成立。
\end{enumerate}

\subsection{对应量对照表}
\begin{center}
\small
\begin{tabularx}{\linewidth}{@{}l X X@{}}
\toprule
\textbf{物理量/关系} & \textbf{静电场（$\rho=0$）} & \textbf{恒定电场（电源外）} \\
\midrule
基本方程（无旋） & $\nabla\times\vec{E}=0$，$\oint\vec{E}\cdot\mathrm{d}\vec{l}=0$ & 同左 \\
基本方程（无源） & $\nabla\cdot\vec{D}=0$，$\oint\vec{D}\cdot\mathrm{d}\vec{S}=0$ & $\nabla\cdot\vec{J}=0$，$\oint\vec{J}\cdot\mathrm{d}\vec{S}=0$ \\
本构方程 & $\vec{D}=\varepsilon\vec{E}$ & $\vec{J}=\gamma\vec{E}$ \\
场量 & $\vec{E}$，$\vec{D}$ & $\vec{E}$，$\vec{J}$ \\
位函数 & 电位 $\varphi$，$\vec{E}=-\nabla\varphi$ & 电位 $\varphi$，$\vec{E}=-\nabla\varphi$ \\
控制方程 & $\nabla^2\varphi=0$ & $\nabla^2\varphi=0$ \\
衔接条件 & $\varphi_1=\varphi_2$，$\varepsilon_1\partial\varphi_1/\partial n=\varepsilon_2\partial\varphi_2/\partial n$ & $\varphi_1=\varphi_2$，$\gamma_1\partial\varphi_1/\partial n=\gamma_2\partial\varphi_2/\partial n$ \\
源量 & 电荷 $q$，面电荷 $\sigma$ & 电流 $I$，面电流 $K$ \\
集总参数 & 电容 $C=q/U$ & 电导 $G=I/U$ \\
参数关系 & $C=\dfrac{\varepsilon\int_S\vec{E}\cdot\mathrm{d}\vec{S}}{\int_A^B\vec{E}\cdot\mathrm{d}\vec{l}}$ & $G=\dfrac{\gamma\int_S\vec{E}\cdot\mathrm{d}\vec{S}}{\int_A^B\vec{E}\cdot\mathrm{d}\vec{l}}$ \\
能量 & 静电能量 $W=\frac12\int\vec{D}\cdot\vec{E}\,\mathrm{d}V$ & 功率损耗 $P=\int\vec{J}\cdot\vec{E}\,\mathrm{d}V$ \\
\bottomrule
\end{tabularx}
\end{center}
由表中两式相比立得核心关系式
\begin{formulaBox}
\[
\frac{C}{G}=\frac{\varepsilon}{\gamma}
\qquad\Longleftrightarrow\qquad
G=\frac{\gamma}{\varepsilon}C
\]
\end{formulaBox}
\textbf{用法}：已知（或易求）电容 $C$，把 $\varepsilon\to\gamma$ 即得电导 $G$；反之亦然。这是计算电导最省力的途径。

\subsection{比拟法的三种用法}
\begin{enumerate}
  \item \textbf{参数置换}：由已知的电容公式直接写出电导公式（$C\to G$，$\varepsilon\to\gamma$）；
  \item \textbf{镜像法}：导电媒质分界面（如土壤—空气）中的线电流问题，可用静电场中点电荷对介质分界面的镜像结果换算（$\varepsilon\to\gamma$，$q\to I$）；
  \item \textbf{模拟实验}：用导电溶液（电解槽）中的恒定电流场模拟难以直接测量的静电场（如高压套管、电缆头、绝缘子电场），测电位分布即得原静电场分布，称为\textbf{电解槽模拟}。
\end{enumerate}

\corepoint{比拟是"同型方程 + 同边界条件"的数学对应，不是物理等价。使用时必须声明三条适用条件，尤其是"电源外"和"$\varepsilon/\gamma$ 为常数"。若媒质不均匀且不满足 $\varepsilon_1/\varepsilon_2=\gamma_1/\gamma_2$，只能逐区求解，不能整体置换。}

% ---------------------------------------------------------
\section{电导和部分电导}

\subsection{电导与电阻的定义}
\textbf{电导} $G$ 定义为流经导电媒质的电流与两端电压之比（电阻 $R=1/G$）：
\begin{formulaBox}
\[
G=\frac{I}{U}
=\frac{\displaystyle\iint_S \vec{J}\cdot\mathrm{d}\vec{S}}
     {\displaystyle\int_A^B \vec{E}\cdot\mathrm{d}\vec{l}}
=\frac{\gamma\displaystyle\iint_S \vec{E}\cdot\mathrm{d}\vec{S}}
     {\displaystyle\int_A^B \vec{E}\cdot\mathrm{d}\vec{l}}
\qquad(\mathrm{S})
\]
\end{formulaBox}
可见 $G$ 只与两电极的形状、尺寸、相互位置及媒质电导率有关，与所加电压、电流大小无关；\emph{电阻（电导）的计算实质上是电场的计算}。

\textbf{三条计算路线}：
\begin{enumerate}
  \item 假设电流 $I$：$I\to\vec{J}\to\vec{E}=\vec{J}/\gamma\to U=\int\vec{E}\cdot\mathrm{d}\vec{l}\to G=I/U$（有对称性时最常用）；
  \item 假设电压 $U$：$U\to\vec{E}\to\vec{J}=\gamma\vec{E}\to I=\iint\vec{J}\cdot\mathrm{d}\vec{S}\to G$（不常用）；
  \item 解拉普拉斯方程求 $\varphi$，再积分求 $I$、$U$（一般情形）；
  \item 静电比拟：先求同形电极系统的 $C$，再由 $G=(\gamma/\varepsilon)C$ 换算。
\end{enumerate}

\textbf{三个必须记住的结果}（均可由比拟法立即得到）：
\begin{formulaBox}
\[
G_{\text{平板}}=\frac{\gamma S}{d},\qquad
G_{\text{同轴}}=\frac{2\pi\gamma l}{\ln(R_2/R_1)},\qquad
G_{\text{同心球}}=\frac{4\pi\gamma R_1R_2}{R_2-R_1}
\]
\end{formulaBox}
同轴电缆的\textbf{绝缘电阻}即 $R=1/G=\dfrac{\ln(R_2/R_1)}{2\pi\gamma l}$，它是用绝缘材料隔开的两个导体之间的直流电阻，由绝缘缺陷（气隙、开裂、受潮、老化）引起的微弱泄漏电流决定，可用于判断绝缘状态。

\subsection{部分电导}
当导电媒质中有三个及以上良导体电极（可视为等位体）时，任意两电极间的电流不仅取决于它们自身的电压，还受其它电极电压影响，单一电导不足以描述，需引入\textbf{部分电导}。

设 $(n+1)$ 个电极（0 号取作参考电极），线性各向同性媒质中由叠加原理可写出各电极电流与电压的关系：
\begin{formulaBox}
\[
\begin{cases}
I_1=G_{10}U_{10}+G_{12}U_{12}+\cdots+G_{1n}U_{1n}\\
I_2=G_{21}U_{21}+G_{20}U_{20}+\cdots+G_{2n}U_{2n}\\
\qquad\vdots
\end{cases}
\qquad U_{kj}=\varphi_k-\varphi_j
\]
\end{formulaBox}
\begin{itemize}
  \item $G_{k0}$ 称为\textbf{自有部分电导}（电极 $k$ 与参考电极之间）；
  \item $G_{kj}$（$k\neq j$）称为\textbf{互有部分电导}（电极 $k$、$j$ 之间）；
  \item 全部部分电导均为正值（互有部分电导描述两电极间的直接漏电通道，故与"互有电导"式的负系数不同）；
  \item $(n+1)$ 个电极的系统共有 $\dfrac{n(n+1)}{2}$ 个部分电导，与部分电容一一对应（$C_{kj}\leftrightarrow G_{kj}$，$\varepsilon\leftrightarrow\gamma$）。
\end{itemize}
工程上"三极法"测接地电阻、多电极系统的漏电分析都用部分电导描述。

\subsection{接地电阻}
把电气设备的一部分与大地连接称为\textbf{接地}；埋入地中并与设备相连的导体（系统）称为\textbf{接地体}，接地体与接地线合称接地装置。\textbf{接地电阻}是电流由接地装置流入大地、再经大地向远处扩散所遇到的电阻，包括接地线与接地体本身电阻、接地体与大地间的接触电阻、大地电阻，其中\textbf{以前者为主的是大地电阻}（前三项远小于最后一项）。

计算时把接地体看作电极、取无穷远为零电位，且因电流在电极附近密度最大，接地电阻主要集中于电极附近，故远端电流如何汇聚影响很小。

\textbf{深埋导体球}（半径 $a$，可不计地面影响）：$J=I/(4\pi r^2)$，$E=I/(4\pi\gamma r^2)$，
\begin{formulaBox}
\[
R=\int_a^\infty \frac{I}{4\pi\gamma r^2}\,\mathrm{d}r=\frac{1}{4\pi\gamma a}
\]
\end{formulaBox}
\textbf{紧靠地面的半球形接地体}：用镜像法，实际电极与其镜像构成完整球，系统中流出的电流为实际的两倍，故实际接地电阻为完整球的两倍：
\begin{formulaBox}
\[
R=2\cdot\frac{1}{4\pi\gamma a}=\frac{1}{2\pi\gamma a}
\]
\end{formulaBox}

\subsection{跨步电压与危险区}
接地体附近有大的入地电流时，地面上相距一个步距的两点间存在电位差，称为\textbf{跨步电压}。对半径为 $a$ 、入地电流为 $I$ 的半球形接地体，距球心 $r$ 处
\begin{formulaBox}
\[
J=\frac{I}{2\pi r^2},\qquad E=\frac{I}{2\pi\gamma r^2},\qquad
\varphi(r)=\frac{I}{2\pi\gamma r}
\]
\end{formulaBox}
设 $A$、$B$ 两点与球心距离分别为 $l$、$l+b$（$b$ 为步距），则跨步电压
\begin{formulaBox}
\[
U_{BA}=\varphi_A-\varphi_B=\frac{I}{2\pi\gamma}\Big(\frac{1}{l}-\frac{1}{l+b}\Big)
=\frac{I\,b}{2\pi\gamma\,l(l+b)}
\]
\end{formulaBox}
若对人体有危险的临界电压为 $U_0$，令 $U_{BA}=U_0$ 解出的 $l$ 即\textbf{危险区}边界，危险区是以接地体中心为圆心、半径 $l$ 的圆面积。当 $l\gg b$ 时可近似
\begin{formulaBox}
\[
U_{BA}\approx\frac{I\,b}{2\pi\gamma l^2}
\ \Longrightarrow\
l\approx\sqrt{\frac{I\,b}{2\pi\gamma U_0}}
\]
\end{formulaBox}
真正危及生命的不是电压而是通过人体的电流（工频电流超过约 30\,mA 即有生命危险），故 $U_0$ 应按人体允许电流与人体电阻折算。

\corepoint{接地电阻常见三个错误：把"接地体本身的电阻"当接地电阻（它只占极小部分）；忘记半球形接地体要\emph{乘 2}（镜像法带来的倍数）；把跨步电压写成 $\varphi_A+\varphi_B$ 或漏掉步距 $b$ 的影响。}

% ==================== 第3章 恒定磁场 ====================
% 子文件：仅正文，由 1.电磁场笔记.tex \input 引入
% 教材：《工程电磁场导论》第3章 §3-1 ~ §3-9
% 公式统一用 formulaBox（本目录 3.矢量分析.tex 中定义）
% ==========================================================

\chapter{恒定磁场}

\textbf{研究对象}：恒定电流（不随时间变化）所激发的磁场，场量满足 $\partial(\bullet)/\partial t=0$。

\textbf{结构对照}：静电场由 $\nabla\times\vec{E}=0$ 引入标量电位、由 $\nabla\cdot\vec{D}=\rho$ 得泊松方程；恒定磁场恰好相反——由 $\nabla\cdot\vec{B}=0$ 引入\emph{矢量}磁位 $\vec{A}$、由 $\nabla\times\vec{H}=\vec{J}$ 得泊松方程。这个"叉与点互换"是本章公式的总来源。

% ---------------------------------------------------------
\section{磁感应强度}

\subsection{安培力定律与毕奥—沙伐定律}
真空中两细导线回路 $l$、$l'$（电流 $I$、$I'$）间的作用力为
\begin{formulaBox}
\[
\vec{F}=\frac{\mu_0}{4\pi}\oint_l\oint_{l'}\frac{I\,\mathrm{d}\vec{l}\times\big(I'\,\mathrm{d}\vec{l}'\times\vec{e}_R\big)}{R^2}
\]
\end{formulaBox}
$\mu_0=4\pi\times10^{-7}\ \mathrm{H/m}$ 为真空磁导率，$\vec{e}_R$ 由 $I'\mathrm{d}\vec{l}'$ 指向 $I\mathrm{d}\vec{l}$。它与库仑定律地位相当，但含双重矢量积，方向关系复杂——根源是电流元的矢量性；且只适用于\emph{闭合}回路（恒定电流必须闭合）。

把上式括号内的量理解为回路 $l'$ 在 $I\mathrm{d}\vec{l}$ 处产生的效应，定义\textbf{磁感应强度}（磁通密度）：
\begin{formulaBox}
\[
\vec{B}=\frac{\mu_0}{4\pi}\oint_{l'}\frac{I'\,\mathrm{d}\vec{l}'\times\vec{e}_R}{R^2}
=\frac{\mu_0}{4\pi}\int_{V'}\frac{\vec{J}\times\vec{e}_R}{R^2}\,\mathrm{d}V'
=\frac{\mu_0}{4\pi}\int_{S'}\frac{\vec{K}\times\vec{e}_R}{R^2}\,\mathrm{d}S'
\qquad(\mathrm{T})
\]
\end{formulaBox}
（依次对应线、体、面电流，即\textbf{毕奥—沙伐定律}。）于是安培力定律写成一般形式
\begin{formulaBox}
\[
\mathrm{d}\vec{f}=I\,\mathrm{d}\vec{l}\times\vec{B},
\qquad
\vec{F}=\oint_l I\,\mathrm{d}\vec{l}\times\vec{B}
\qquad(\text{均匀场中 }\vec{F}=I\vec{l}\times\vec{B})
\]
\end{formulaBox}
三个可直接引用的积分结论（$\rho$ 为场点到导线的垂距，$K$ 为面电流线密度）：
\begin{formulaBox}
\[
B_{\text{无限长直导线}}=\frac{\mu_0 I}{2\pi\rho},
\qquad
B_{\text{长 }2L\text{ 直导线中点垂距}\rho}=\frac{\mu_0 I L}{2\pi\rho\sqrt{\rho^2+L^2}},
\qquad
B_{\text{无限大电流片}}=\frac{\mu_0 K}{2}
\]
\end{formulaBox}

\subsection{洛伦兹力与磁感应线}
\begin{formulaBox}
\[
\vec{F}=q\vec{v}\times\vec{B},
\qquad
\vec{B}\times\mathrm{d}\vec{l}=0
\ \Big(\text{即}\ \frac{\mathrm{d}x}{B_x}=\frac{\mathrm{d}y}{B_y}=\frac{\mathrm{d}z}{B_z}\Big)
\]
\end{formulaBox}
前者为运动电荷受的\textbf{洛伦兹力}（静止电荷不受磁场力；力恒垂直于速度，只改方向不改量值，故\textbf{洛伦兹力不做功}）；后者是\textbf{磁感应线}的微分方程。

\subsection{磁通与磁通连续性原理}
\begin{formulaBox}
\[
\Phi=\iint_S \vec{B}\cdot\mathrm{d}\vec{S}\ (\mathrm{Wb}),
\qquad
\oint_S \vec{B}\cdot\mathrm{d}\vec{S}=0
\qquad\Longleftrightarrow\qquad
\nabla\cdot\vec{B}=0
\]
\end{formulaBox}
磁感应线闭合、无始无终，自然界不存在与电荷对应的磁荷，故穿出任意闭合面的净磁通为零——\textbf{磁通连续性原理}：恒定磁场是\emph{无散场}，穿出等于穿入，同一闭合曲线为周界的各曲面上磁通相同。

\corepoint{$\nabla\cdot\vec{B}=0$ 只说明"无磁荷源"，不说明磁场没有源——磁场的涡旋源是电流。另注意 $\Phi$ 是穿过某个\emph{面}的通量，$\oint_S\vec{B}\cdot\mathrm{d}\vec{S}=0$ 只对\emph{闭合}面成立。}

% ---------------------------------------------------------
\section{安培环路定律}

\subsection{真空中的安培环路定律}
\begin{formulaBox}
\[
\oint_l \vec{B}\cdot\mathrm{d}\vec{l}=\mu_0\sum I_i
\qquad\Longleftrightarrow\qquad
\nabla\times\vec{B}=\mu_0\vec{J}_c
\]
\end{formulaBox}
电流方向与回路绕向符合\textbf{右手螺旋关系}时取正。磁场是\textbf{有旋场}，$\vec{B}$ 的环流一般不恒为零，故磁场非保守场、一般无"磁势"；磁感线与电流垂直并相互交链。

\textbf{用它求 $\vec{B}$ 的条件}：磁场具有某种对称性，且能选到这样的积分路径——路径上各点 $\vec{B}$ 与 $\mathrm{d}\vec{l}$ 夹角相同、$B$ 值相等，从而把 $\oint$ 化为代数式。典型应用如载流圆柱（$\rho<R$ 时 $B=\mu_0I\rho/(2\pi R^2)$，$\rho>R$ 时 $B=\mu_0I/(2\pi\rho)$）、同轴电缆、螺环等。

\subsection{媒质的磁化与磁化强度}
分子（原子）中电子的轨道运动与自旋都等效于电流，其对外磁效应的总可用一个等效圆电流表示，称\textbf{分子电流}（束缚电流），它不引起电荷迁移，但能激发磁场。分子电流具有\textbf{分子磁矩} $\vec{m}=I\vec{S}$（$\vec{S}$ 方向与电流绕向服从右手螺旋）。无外场时分子磁矩因热运动随机排列、总磁矩为零；放入外场后受转矩 $\vec{T}=\vec{m}\times\vec{B}$ 而趋于与外场同向，物质显磁性——即\textbf{磁化}，其宏观后果是出现\textbf{磁化电流}，附加磁场改变原场分布。

\textbf{磁化强度}（单位体积内分子磁矩的矢量和）与磁化电流的关系为
\begin{formulaBox}
\[
\vec{M}=\lim_{\Delta V\to 0}\frac{\sum_i\vec{m}_i}{\Delta V}\ (\mathrm{A/m}),
\qquad
\vec{J}_m=\nabla\times\vec{M},
\qquad
\oint_l\vec{M}\cdot\mathrm{d}\vec{l}=I_m,
\qquad
\vec{K}_m=\vec{M}\times\vec{e}_n
\]
\end{formulaBox}
于是"有磁介质时的磁场"可化为"把磁化电流与自由电流一起在真空中求场"。

\subsection{磁场强度与一般形式的安培环路定律}
由 $\nabla\times\vec{B}=\mu_0(\vec{J}_c+\vec{J}_m)$、$\vec{J}_m=\nabla\times\vec{M}$ 移项得
\begin{formulaBox}
\[
\nabla\times\Big(\frac{\vec{B}}{\mu_0}-\vec{M}\Big)=\vec{J}_c
\ \Longrightarrow\
\vec{H}=\frac{\vec{B}}{\mu_0}-\vec{M},
\qquad
\nabla\times\vec{H}=\vec{J}_c
\ \Longleftrightarrow\
\oint_l\vec{H}\cdot\mathrm{d}\vec{l}=\sum I
\]
\end{formulaBox}
$\vec{H}$ 称\textbf{磁场强度}（$\mathrm{A/m}$）。\textbf{物理意义}：$\vec{H}$ 的\emph{环流}只由自由电流决定、与磁化电流（即媒质分布）无关；但 $\vec{H}$ 的\emph{分布}仍与媒质有关，两者不可混同。

\subsection{磁介质的物性方程}
\begin{formulaBox}
\[
\vec{B}=\mu_0(\vec{H}+\vec{M}),
\quad
\vec{M}=\chi_m\vec{H},
\quad
\vec{B}=\mu_0(1+\chi_m)\vec{H}=\mu_0\mu_r\vec{H}=\mu\vec{H}
\]
\end{formulaBox}
$\chi_m$ 为磁化率，$\mu_r=1+\chi_m$ 为相对磁导率，$\mu=\mu_0\mu_r$ 为磁导率（$\mathrm{H/m}$）。顺磁质 $\chi_m>0$、抗磁质 $\chi_m<0$（均极小），铁磁质 $\mu_r$ 可达 $10^3\sim10^5$ 且非线性、有磁滞与饱和。注意 $\vec{B}=\mu\vec{H}$ \textbf{仅适用于各向同性线性}媒质，而 $\vec{B}=\mu_0(\vec{H}+\vec{M})$ 无此限制。

\textbf{推论}：若电流周围无限充满均匀导磁媒质，则各点 $\vec{B}$ 方向同真空、量值增大 $\mu_r$ 倍（把 $\mu_0$ 换成 $\mu$）。

\subsection{电介质极化与磁介质磁化的对照}
\begin{center}
\small
\begin{tabularx}{\linewidth}{@{}l X X@{}}
\toprule
 & \textbf{电介质极化} & \textbf{磁介质磁化} \\
\midrule
基本单元 & 电偶极矩 $\vec{p}=q\vec{d}$ & 分子磁矩 $\vec{m}=I\vec{S}$ \\
强度量 & $\vec{P}=\lim\dfrac{\sum\vec{p}_i}{\Delta V}$ & $\vec{M}=\lim\dfrac{\sum\vec{m}_i}{\Delta V}$ \\
附加源 & $\sigma_p=\vec{P}\cdot\vec{e}_n$，$\rho_p=-\nabla\cdot\vec{P}$ & $\vec{K}_m=\vec{M}\times\vec{e}_n$，$\vec{J}_m=\nabla\times\vec{M}$ \\
辅助场量 & $\vec{D}=\varepsilon_0\vec{E}+\vec{P}$ & $\vec{H}=\vec{B}/\mu_0-\vec{M}$ \\
媒质方程 & $\vec{D}=\varepsilon\vec{E}$ & $\vec{B}=\mu\vec{H}$ \\
散/旋 & $\nabla\cdot\vec{D}=\rho$，$\nabla\times\vec{E}=0$ & $\nabla\cdot\vec{B}=0$，$\nabla\times\vec{H}=\vec{J}_c$ \\
\bottomrule
\end{tabularx}
\end{center}

\corepoint{最易写错的两组关系：\emph{电}——$\vec{D}$ 的散度由自由电荷决定、$\vec{E}$ 无旋；\emph{磁}——$\vec{B}$ 的散度恒为零、$\vec{H}$ 的旋度由自由电流决定。}

% ---------------------------------------------------------
\section{恒定磁场的基本方程与分界面上的衔接条件}

\subsection{恒定磁场的基本方程组}
\begin{formulaBox}
\[
\begin{cases}
\displaystyle\oint_S\vec{B}\cdot\mathrm{d}\vec{S}=0\\[6pt]
\displaystyle\oint_l\vec{H}\cdot\mathrm{d}\vec{l}=\sum I
\end{cases}
\qquad
\begin{cases}
\nabla\cdot\vec{B}=0\\[4pt]
\nabla\times\vec{H}=\vec{J}_c
\end{cases}
\qquad
\vec{B}=\mu\vec{H}
\]
\end{formulaBox}
$\nabla\cdot\vec{B}=0$：无源场，磁感应线无头无尾、无磁荷；$\nabla\times\vec{H}=\vec{J}_c$：有旋场，电流密度是涡旋源；$\vec{B}=\mu\vec{H}$：各向同性线性媒质的构成方程（铁磁质不成立）。

\subsection{分界面上的衔接条件}
设 $\vec{e}_n$ 由媒质 1 指向媒质 2：
\begin{formulaBox}
\[
\begin{cases}
H_{1t}-H_{2t}=K\ \ \text{或}\ \ \vec{e}_n\times(\vec{H}_2-\vec{H}_1)=\vec{K} & \text{（切向：}H_t\text{跃变）}\\[4pt]
B_{1n}=B_{2n}\ \ \text{或}\ \ \vec{e}_n\cdot(\vec{B}_2-\vec{B}_1)=0 & \text{（法向：}B_n\text{连续）}
\end{cases}
\]
\end{formulaBox}
$H_t$ 的跃变量等于界面上的\textbf{面电流线密度} $K$（由跨界面小矩形回路的 $\oint\vec{H}\cdot\mathrm{d}\vec{l}=\sum I$ 得到）；$B_n$ 连续由跨界面小扁盒的 $\oint\vec{B}\cdot\mathrm{d}\vec{S}=0$ 得到。当界面上无面电流（$K=0$）时 $H_t$ 也连续，但 $B_t$ 仍不连续（$B_t=\mu H_t$，两侧 $\mu$ 不同）。理论上只有理想导体、超导体表面才可能有面电流；工程上仅当集肤效应十分显著时才可近似视为面电流。

\subsection{磁场折射定律}
分界面上无面电流、两侧均为各向同性线性媒质时
\begin{formulaBox}
\[
\frac{\tan\theta_1}{\tan\theta_2}=\frac{\mu_1}{\mu_2}
\]
\end{formulaBox}
若媒质 2 为铁磁质（$\mu_2\gg\mu_1$），除 $\theta_1=90^\circ$ 外 $\theta_2\to90^\circ$：铁磁质内 $B$ 线几乎平行于界面，且 $B_2\gg B_1$，磁通高度集中。即"\textbf{$B$ 线沿铁走}"，这是磁路概念的物理基础。

% ---------------------------------------------------------
\section{磁矢位·恒定磁场的边值问题}

\subsection{磁矢位的引出与库仑规范}
因 $\nabla\cdot\vec{B}=0$，可引入矢量函数使
\begin{formulaBox}
\[
\vec{B}=\nabla\times\vec{A}\qquad(\mathrm{Wb/m})
\]
\end{formulaBox}
代入 $\nabla\times\vec{H}=\vec{J}$、$\vec{B}=\mu\vec{H}$，用 $\nabla\times\nabla\times\vec{A}=\nabla(\nabla\cdot\vec{A})-\nabla^2\vec{A}$：
\begin{formulaBox}
\[
\nabla(\nabla\cdot\vec{A})-\nabla^2\vec{A}=-\mu\vec{J}
\]
\end{formulaBox}
唯一确定矢量场须同时给定旋度与散度，故再规定
\begin{formulaBox}
\[
\nabla\cdot\vec{A}=0\ \text{（库仑规范）}
\qquad\Longrightarrow\qquad
\nabla^2\vec{A}=-\mu\vec{J}\ \text{（矢量泊松方程）}
\]
\end{formulaBox}
直角坐标下即三个形式相同的标量泊松方程 $\nabla^2A_x=-\mu J_x$ 等。当电流分布于有限空间、取无限远处 $\vec{A}=0$ 时
\begin{formulaBox}
\[
\vec{A}=\frac{\mu_0}{4\pi}\int_{V'}\frac{\vec{J}}{R}\mathrm{d}V'
=\frac{\mu_0}{4\pi}\int_{S'}\frac{\vec{K}}{R}\mathrm{d}S'
=\frac{\mu_0}{4\pi}\oint_{l'}\frac{I\,\mathrm{d}\vec{l}'}{R}
\]
\end{formulaBox}
每个元电流产生的 $\vec{A}$ 与该元电流\emph{同方向}（与标量电位的根本差别），故可分量计算；$\vec{A}$ 无物理意义，但便于求 $\vec{B}$、磁通与电感。

\subsection{磁矢位与磁通}
由斯托克斯定理
\begin{formulaBox}
\[
\Phi=\iint_S\vec{B}\cdot\mathrm{d}\vec{S}=\iint_S(\nabla\times\vec{A})\cdot\mathrm{d}\vec{S}=\oint_l\vec{A}\cdot\mathrm{d}\vec{l}
\]
\end{formulaBox}
即 $\vec{A}$ 沿闭合回路的环量等于穿过以该回路为周界的曲面的磁通。对\textbf{平行平面磁场}（$\vec{A}=A_z\vec{e}_z$），$B$ 线即等 $A$ 线，等 $A$ 线密处 $B$ 大。

\subsection{磁矢位的衔接条件与边值问题}
由跨界面小矩形回路的 $\Phi=\oint\vec{A}\cdot\mathrm{d}\vec{l}$（高度趋于零）得切向分量连续，由 $\nabla\cdot\vec{A}=0$（跨界面小扁圆柱）得法向分量连续，故
\begin{formulaBox}
\[
\vec{A}_1=\vec{A}_2\ \text{（界面上 $\vec{A}$ 连续）}
\]
\end{formulaBox}
另一条件由 $\vec{e}_n\times(\vec{H}_2-\vec{H}_1)=\vec{K}$ 换算。对平行平面磁场简化为
\begin{formulaBox}
\[
A_1=A_2,
\qquad
\frac{1}{\mu_1}\frac{\partial A_1}{\partial n}-\frac{1}{\mu_2}\frac{\partial A_2}{\partial n}=K
\]
\end{formulaBox}
"$\nabla^2\vec{A}=-\mu\vec{J}$（无电流处 $\nabla^2\vec{A}=0$）$+$ 衔接条件 $+$ 场域边界条件"即恒定磁场的\textbf{边值问题}（磁矢位表述）。

\subsection{标量磁位}
在\textbf{传导电流处处为零}的区域（$\vec{J}_c=0$），$\nabla\times\vec{H}=0$，可有条件引入\textbf{标量磁位}：
\begin{formulaBox}
\[
\vec{H}=-\nabla\varphi_m,
\quad
\varphi_{mP}=\int_P^Q\vec{H}\cdot\mathrm{d}\vec{l}\ (\mathrm{A}),
\quad
U_{mAB}=\int_A^B\vec{H}\cdot\mathrm{d}\vec{l}=\varphi_{mA}-\varphi_{mB}
\]
\end{formulaBox}
均匀媒质中 $\nabla^2\varphi_m=0$（拉普拉斯方程）；衔接条件为
\begin{formulaBox}
\[
\varphi_{m1}=\varphi_{m2},
\qquad
\mu_1\frac{\partial\varphi_{m1}}{\partial n}=\mu_2\frac{\partial\varphi_{m2}}{\partial n}
\]
\end{formulaBox}
\textbf{多值性}：两点间磁压与积分路径有关，若回路绕电流 $k$ 次则差 $kI$，故即使选定参考点 $\varphi_m$ 仍是多值函数。消除办法是规定积分路径不穿过电流回路所限定的面（\textbf{磁屏障面}），使磁位单值。多值性不影响 $\vec{B}$、$\vec{H}$ 的计算。

\textbf{与电位的三点区别}：(1) 存在条件不同——$\varphi_m$ 只在 $\vec{J}_c=0$ 区可引入；(2) 多值性来源不同——电位只由参考点决定，磁位还受路径影响；(3) 物理意义不同——$\varphi_m$ 纯为计算方便，无物理意义。另：因 $H_t\approx0$，磁导率很大的材料表面近似为\textbf{等磁位面}，此结论在磁路计算中反复使用。

\corepoint{$\vec{A}$ 与 $\varphi_m$ 的分工：有电流的区域只能用 $\vec{A}$（矢量泊松方程），无电流区域才用 $\varphi_m$（标量拉普拉斯方程）。解题第一步先判断区域内有无电流，选错位函数则方程不成立。}

% ---------------------------------------------------------
\section{镜像法}

\subsection{依据}
由\textbf{唯一性定理}，若所求区域内方程相同、边界条件相同，解唯一。故可用镜像电流代替分界面上的（磁化）束缚电流，把分区问题化为均匀空间的场叠加——即\textbf{磁镜像法}。

\subsection{线电流对无限大媒质分界面的镜像}
媒质 1、2 的磁导率为 $\mu_1$、$\mu_2$，媒质 1 中有平行于界面的无限长直导线（电流 $I$）。求媒质 1 的场时，认为全空间充满 $\mu_1$，场由 $I$ 与对称位置的像电流 $I'$ 共同产生；求媒质 2 的场时，认为全空间充满 $\mu_2$，场由原位置的像电流 $I''$ 单独产生。由 $H_{2t}=H_{1t}$、$B_{1n}=B_{2n}$ 联立得
\begin{formulaBox}
\[
I'=\frac{\mu_2-\mu_1}{\mu_2+\mu_1}I,
\qquad
I''=\frac{2\mu_1}{\mu_2+\mu_1}I
\]
\end{formulaBox}
$I'$、$I''$ 的\textbf{参考方向均规定与 $I$ 相同}（正负由此体现）；$I''$ 恒为正，$I'$ 的正负取决于 $(\mu_2-\mu_1)$。用法记忆：\emph{媒质 1 场} $=I+I'$（全在 $\mu_1$ 中）；\emph{媒质 2 场} $=I''$（全在 $\mu_2$ 中），两者不可混用于同一区域。

\subsection{两种典型特例}
\begin{enumerate}
  \item \textbf{导线在空气中，另一侧为铁磁质}（$\mu_1=\mu_0$，$\mu_2\to\infty$）：$I'\approx I$，$I''\approx0$。铁磁质内 $H_2\approx0$，但\emph{不能}认为 $B_2$ 也为零——由 $B_n$ 连续，磁通仍穿入铁磁质。
  \item \textbf{导线在铁磁质中，另一侧为空气}（$\mu_1\to\infty$，$\mu_2=\mu_0$）：$I'\approx-I$，$I''\approx2I$。空气中的磁感应强度比全空间为空气时\textbf{增大一倍}（电流相同）。
\end{enumerate}

\subsection{适用条件与易错点}
\begin{enumerate}
  \item 界面须为无限大平面，两媒质各自均匀、各向同性、线性；
  \item 镜像电流位置由对称确定，符号约定须统一（都相对 $I$ 的参考方向）；
  \item 只适用于线电流及可等效为线电流的情形，有限长导线与线圈需另行处理。
\end{enumerate}

\corepoint{磁镜像与电镜像"符号相反"：静电场中电荷靠近理想导体（$\varepsilon_2\to\infty$）时镜像电荷异号（$-q$）；磁场中电流靠近铁磁质（$\mu_2\to\infty$）时镜像电流同号（$+I$）。记不住就回到 $B_n$ 连续这一条件重推——铁磁质"导通磁通"，导体"终止电力线"，性质相反。}

% ---------------------------------------------------------
\section{电感}

\subsection{自感}
线性媒质中磁链与电流成正比：
\begin{formulaBox}
\[
\psi=LI,
\qquad
L=\frac{\psi}{I}\ (\mathrm{H}),
\qquad
L=L_i+L_o
\]
\end{formulaBox}
$L$ 只与回路尺寸、几何形状及媒质分布有关，与电流、磁链数值无关。
\begin{itemize}
  \item \textbf{内自感} $L_i=\psi_i/I$：内磁通在导线内部闭合，只与\emph{部分}电流交链，求元内磁链时交链电流按 $(I'/I)=(\rho/R)^2$ 折算，故 $\mathrm{d}\psi_i=(\rho^2/R^2)\mathrm{d}\Phi$；
  \item \textbf{外自感} $L_o=\psi_o/I$：外磁通在导线外部闭合，与全部电流交链，$\mathrm{d}\psi_o=\mathrm{d}\Phi$。
\end{itemize}
圆截面导线的内自感为
\begin{formulaBox}
\[
L_i=\frac{\mu_0 l}{8\pi}
\qquad\text{（只与长度 }l\text{ 有关，与半径无关）}
\]
\end{formulaBox}

\subsection{两个必须掌握的自感结果}
\textbf{同轴电缆}（长 $l$，内导体半径 $R_1$，外导体内半径 $R_2$，外壳很薄，磁导率均为 $\mu_0$）：
\begin{formulaBox}
\[
L=\frac{\mu_0 l}{8\pi}+\frac{\mu_0 l}{2\pi}\ln\frac{R_2}{R_1}
\]
\end{formulaBox}
\textbf{二线传输线}（半径 $R$ 的两长直导线，轴线间距 $D\gg R$）：
\begin{formulaBox}
\[
L=\frac{\mu_0 l}{\pi}\Big[\ln\frac{D}{R}+\frac14\Big]
\]
\end{formulaBox}
（$\ln(D/R)$ 为外自感，$1/4$ 为两根导线内自感之和。）

\textbf{计算步骤}：设电流 $I\to$ 由安培环路定律求 $\vec{H}$、$\vec{B}=\mu\vec{H}\to$ 取面元求 $\mathrm{d}\Phi=B\mathrm{d}S\to$ 按交链电流比例得 $\mathrm{d}\psi$、积分得 $\psi=\psi_i+\psi_o\to L=\psi/I$。

\subsection{互感}
\begin{formulaBox}
\[
\psi_{21}=M_{21}I_1,
\qquad
M_{21}=\frac{\psi_{21}}{I_1},
\qquad
M_{12}=M_{21}=M\ (\mathrm{H})
\]
\end{formulaBox}
互感不仅与两回路的形状、尺寸、磁导率有关，还与\textbf{相互位置}有关。对细导线回路的一般公式（\textbf{诺伊曼公式}）为
\begin{formulaBox}
\[
M=\frac{\mu_0}{4\pi}\oint_{l_1}\oint_{l_2}\frac{\mathrm{d}\vec{l}_1\cdot\mathrm{d}\vec{l}_2}{R}
\qquad
\Big(\text{紧绕 }N_1,N_2\text{ 匝时乘 }N_1N_2\Big)
\]
\end{formulaBox}
它说明互感只取决于两回路的几何，无需先求场；但双重线积分通常较繁，工程上多用磁链法直接算 $\psi_{21}$。

\subsection{用磁场能量求电感}
对单个载流回路 $W_m=\frac12 LI^2$，故
\begin{formulaBox}
\[
L=\frac{2W_m}{I^2}
\]
\end{formulaBox}
当 $\vec{B}$、$\vec{H}$ 已用数值方法求出时，能量法求电感比磁链法更方便，是工程数值计算的常用做法。

% ---------------------------------------------------------
\section{磁场能量与力}

\subsection{磁场能量}
设磁场建立过程缓慢、媒质线性、无辐射与其它损耗，则外源做功全部转为磁场能量：
\begin{formulaBox}
\[
W_m=\frac12\sum_{k=1}^{n}I_k\psi_k,
\qquad
W_m=\frac12\int_V\vec{A}\cdot\vec{J}\,\mathrm{d}V
=\frac12\int_V\vec{B}\cdot\vec{H}\,\mathrm{d}V
\]
\end{formulaBox}
单回路为 $W_m=\frac12LI^2$；两回路为
\begin{formulaBox}
\[
W_m=\frac12L_1I_1^2+\frac12L_2I_2^2+MI_1I_2
\]
\end{formulaBox}
前两项为\textbf{自有能}（恒正，只与各自回路的电流、自感有关）；$MI_1I_2$ 为\textbf{互有能}（可正可负：两回路电流同时从同名端流入时取正）。可见磁场能量只与最终状态有关，与建立过程无关。

\subsection{磁场能量密度}
\begin{formulaBox}
\[
w_m=\frac12\vec{B}\cdot\vec{H}=\frac{B^2}{2\mu}=\frac12\mu H^2
\qquad(\mathrm{J/m^3})
\]
\end{formulaBox}
磁场能量并非只存在于电流回路内，而是\textbf{分布在整个磁场空间}——场是能量的载体，与静电场 $w_e=\frac12\vec{D}\cdot\vec{E}$ 完全对应。

\subsection{虚位移法求磁场力}
设 $n$ 个载流回路接于电压源，仅第 $p$ 号回路可沿广义坐标 $g$ 移动。能量守恒为
\begin{formulaBox}
\[
\mathrm{d}W=\mathrm{d}W_m+f\,\mathrm{d}g,
\qquad
\mathrm{d}W=\sum_k I_k\,\mathrm{d}\psi_k
\]
\end{formulaBox}
\begin{enumerate}
  \item \textbf{电流不变}（各回路接恒流源）：$\mathrm{d}W=2\mathrm{d}W_m$，外源供能一半储能、一半作机械功，
  \begin{formulaBox}
  \[
  f=\left.\frac{\partial W_m}{\partial g}\right|_{I=\text{常数}}
  \]
  \end{formulaBox}
  \item \textbf{磁链不变}（回路开路、无感应电流）：$\mathrm{d}W=0$，机械功只能取自磁场能量的减少，
  \begin{formulaBox}
  \[
  f=-\left.\frac{\partial W_m}{\partial g}\right|_{\psi=\text{常数}}
  \]
  \end{formulaBox}
\end{enumerate}
两式在当时状态下给出同一个力，可任选其一。求相互作用力时，写出相互能 $MI_1I_2$ 后对 $g$ 求偏导即可。

\textbf{电磁铁起重力}：气隙磁场近似均匀、铁芯内储能可忽略，每个气隙储能 $W_m=\dfrac{B^2}{2\mu_0}Sl$，对气隙长度 $l$ 求偏导：
\begin{formulaBox}
\[
F=\frac{B^2S}{2\mu_0}
\qquad\text{（负号表示力使气隙缩短，表现为吸力）}
\]
\end{formulaBox}

\subsection{磁场力面密度与法拉第观点}
\begin{formulaBox}
\[
f_0=\frac12\frac{B^2}{\mu_0}=\frac12\mu_0H^2
\qquad\text{（等于该处磁场能量体密度）}
\]
\end{formulaBox}
\textbf{法拉第观点}：每束磁感应线形成的"磁感应管"沿轴向受\textbf{纵张力}、垂直方向受\textbf{侧压力}，二者数值相等：
\begin{formulaBox}
\[
\frac12BH=\frac12\mu H^2=\frac{B^2}{2\mu}
\]
\end{formulaBox}
据此可定性判断受力（气隙中 $B$ 管沿轴向有收缩趋势，故磁极间表现为吸力）。对两种媒质的分界面：
\begin{formulaBox}
\[
f_0=\frac{\mu_2-\mu_1}{2\mu_1\mu_2}\big(B_{1n}^2+\mu_1\mu_2H_{1t}^2\big)
\]
\end{formulaBox}
无论磁场方向如何，此力\textbf{总垂直于界面}，且\textbf{总是由磁导率较大的媒质指向较小的媒质}。

\corepoint{虚位移法三处高频错误：(1) 恒流与恒磁链的符号用反（恒流取 $+\partial W_m/\partial g$，恒磁链取 $-\partial W_m/\partial g$）；(2) 忘记恒流时"外源供能是储能增量的两倍"；(3) 把 $f_0=B^2/(2\mu_0)$ 当总力——它只是单位面积上的力。}

% ---------------------------------------------------------
\section{磁路及其计算}

\subsection{磁路的概念}
铁磁材料 $\mu\gg\mu_0$，由折射定律知 $B$ 线几乎平行于铁磁质表面、并高度集中于内部（"$B$ 线沿铁走"）。把铁磁质做成闭合或基本闭合的形状，磁通便主要集中在铁心内，与电流集中在导体内相似。这种\textbf{能使磁通集中通过的路径}称为\textbf{磁路}。
\begin{itemize}
  \item \textbf{主磁通} $\Phi$：沿磁路（含气隙）闭合的磁通；
  \item \textbf{漏磁通}：穿出铁心、经磁路周围非铁磁质闭合的磁通。
\end{itemize}
绕在铁心上的线圈通以较小的励磁电流即可获得较强磁场，且磁场差不多被约束在磁路内，周围非铁磁质中很弱。

\subsection{磁路定律}
对无分支闭合铁心，沿一条闭合磁力线应用安培环路定律，并由 $\Phi=BS$、$H=\Phi/(\mu S)$ 得
\begin{formulaBox}
\[
\oint_l\vec{H}\cdot\mathrm{d}\vec{l}=NI
\ \Longrightarrow\
\Phi\int_l\frac{\mathrm{d}l}{\mu S}=NI
\]
\end{formulaBox}
与电阻 $R=\int\mathrm{d}l/(\gamma S)$ 对比，定义\textbf{磁阻}与\textbf{磁动势}
\begin{formulaBox}
\[
R_m=\int_l\frac{\mathrm{d}l}{\mu S}
\ \ \Big(\text{均匀段 }R_m=\frac{l}{\mu S}\Big),
\qquad
F_m=NI\ (\mathrm{A})
\]
\end{formulaBox}
得\textbf{无分支闭合磁路的欧姆定律}
\begin{formulaBox}
\[
\Phi R_m=F_m
\]
\end{formulaBox}
对有分支磁路，由磁通连续性与安培环路定律得\textbf{磁路定律}（对应电路的两个基尔霍夫定律）：
\begin{formulaBox}
\[
\sum\Phi=0\ (\text{节点}),
\qquad
\sum\Phi R_m=\sum F_m\ (\text{回路})
\]
\end{formulaBox}
磁阻的串并联与电阻同形：\textbf{串联}（各段 $\Phi$ 相同）$R_m=\sum_k R_{mk}$，铁心带气隙即两段磁阻串联；\textbf{并联}（各分支磁压相同）$\dfrac{1}{R_m}=\sum_k\dfrac{1}{R_{mk}}$。
\begin{center}
\small
\begin{tabularx}{\linewidth}{@{}l X X@{}}
\toprule
 & \textbf{电路} & \textbf{磁路} \\
\midrule
"流量" & 电流 $I$（连续） & 磁通 $\Phi$（连续） \\
"势" & 电动势 $\varepsilon$（电源） & 磁动势 $F_m=NI$（励磁线圈） \\
"阻力" & 电阻 $R=\dfrac{l}{\gamma S}$ & 磁阻 $R_m=\dfrac{l}{\mu S}$ \\
物性参数 & 电导率 $\gamma$ & 磁导率 $\mu$ \\
定律 & $IR=\varepsilon$，$\sum I=0$，$\sum IR=\sum\varepsilon$ & $\Phi R_m=F_m$，$\sum\Phi=0$，$\sum\Phi R_m=\sum F_m$ \\
"绝缘" & 理想绝缘体 $\gamma=0$，无漏电流 & \textbf{不存在磁绝缘体}，总有漏磁 \\
\bottomrule
\end{tabularx}
\end{center}
磁路定律是在"不计漏磁、$B$ 线沿铁心走向、截面上 $B$ 均匀"等近似下得到的，本质是\textbf{估算}，但对电机、变压器、电磁铁与仪表设计十分必要（含永磁体的磁路更复杂，此处从略）。

由 $\psi=N\Phi$ 还可得电感与磁阻的关系
\begin{formulaBox}
\[
L=\frac{\psi}{I}=\frac{N\Phi}{I}=\frac{N^2}{R_m}
\]
\end{formulaBox}
即\textbf{磁阻越小、电感越大}——这解释了铁心线圈为何电感远大于空心线圈，也解释了开气隙后电感为何急剧下降。

\subsection{磁路计算}
恒定磁通磁路（$\Phi$ 不随时间变化）的计算目的是找出 $\Phi$ 与 $F_m=NI$ 的关系，分两类问题：
\begin{enumerate}
  \item \textbf{已知 $\Phi$ 求 $NI$}（正面问题）：$\Phi$ 与各段截面 $\to$ 各段 $B$；查材料的 $B$—$H$ 曲线得各段 $H$（线性段可直接 $H=B/\mu$）；求各段磁压 $Hl$；相加得 $NI=\sum Hl$。
  \item \textbf{已知 $NI$ 求 $\Phi$}（反面问题）：铁磁质非线性使 $\mu$ 随 $H$ 变化，不能一次算出，须用\textbf{试探法}——先假设 $\Phi$（例如粗略认为全部磁动势都降在气隙上），算出对应 $NI$ 与给定值比较后修正，反复数次直至吻合；实质是"已知 $\Phi$ 求 $NI$"的多次重复。
\end{enumerate}
\textbf{气隙的影响}：$R_\delta=\dfrac{\delta}{\mu_0 S}$，因 $\mu_0\ll\mu_{\text{铁}}$，即使 $\delta$ 很小，气隙磁压降也常占全部磁动势的绝大部分——\emph{气隙决定磁路工作状态}，这是磁路计算第一条经验。

\textbf{边缘效应}：气隙处 $B$ 线向外扩散，实际有效面积略大于铁心截面，气隙磁阻比 $\delta/(\mu_0S)$ 略小。工程上用"有效面积"（如按 $(a+\delta)(b+\delta)$ 放大气隙截面）或大于 1 的卡氏系数修正，$\delta$ 越小越可忽略。

\subsection{磁屏蔽}
用高 $\mu$ 铁磁质做屏蔽罩，把需保护的空间罩起来。用磁阻并联解释：外界磁通遇到"铁壳壁"与"空腔"两条并联路径，空腔的 $\mu$ 远小于罩壁、磁阻远大于罩壁，故绝大部分 $B$ 线沿罩壁通过，进入空腔的很少。

\begin{itemize}
  \item 与静电屏蔽不同：闭合导体空腔内静电场为零，而外磁场中闭合铁磁空腔内的磁场\textbf{并不为零}，磁屏蔽效果远不如静电屏蔽；
  \item 改善办法：加厚罩壁或多层铁壳，把漏入的残余磁通逐次屏蔽；
  \item 不宜用于高频交变磁场，否则罩内产生很大的铁损（涡流、磁滞损耗）。
\end{itemize}

\corepoint{磁路计算三个易错点：(1) 把磁阻写成 $\mu S/l$（应为 $l/(\mu S)$，$\mu$ 对应 $\gamma$ 而非 $1/\gamma$）；(2) 漏算气隙——气隙磁阻常是铁心磁阻的几十上百倍，漏算会使 $NI$ 严重偏小；(3) 对铁磁段直接用 $H=B/\mu$ 而不查 $B$—$H$ 曲线（铁磁质 $\mu$ 不是常数）。}

% ============================================================
% 时变电磁场（子文件，只放内容：不重复 \usepackage、不改字号）
% 对应教材：冯慈璋·马西奎《工程电磁场导论》第 4 章
% 节名与教材 §4-1 ～ §4-6 一一对齐，章号、节号由主文档自动生成
% ============================================================

% formulaBox 原本由 3.矢量分析.tex 定义；这里做一次“未定义才定义”的保护，
% 避免那个文件被改写后本文件整章报 Environment formulaBox undefined。
\makeatletter
\@ifundefined{formulaBox}{%
  \newenvironment{formulaBox}{%
    \par\begin{adjustbox}{max width=\linewidth}%
    \begin{tcolorbox}[colback=blue!8!white,colframe=transparent,boxrule=0pt,halign=center,boxsep=2pt,left=4pt,right=4pt,top=3pt,bottom=3pt]%
  }{%
    \end{tcolorbox}\end{adjustbox}\par
  }%
}{}
\makeatother

\chapter{时变电磁场}

本章讨论随时间变化的电磁场：场量既是空间坐标的函数，又是时间的函数，电场与磁场互为源、构成
统一的电磁场。全章六节对应教材 \S 4-1～\S 4-6：电磁感应定律和全电流定律 → 电磁场基本方程组·
分界面上的衔接条件 → 动态位及其积分解 → 电磁功率流和坡印亭矢量 → 正弦电磁场 → 电磁辐射。
前四节是“建方程 + 求位 + 算功率”的主干，后两节给出最常用的复数化工具与最简单的辐射模型。

\section{电磁感应定律和全电流定律}

本节是全章的起点：电磁感应与全电流两条定律，分别回答“变化的磁场如何产生电场”和
“变化的电场如何产生磁场”。

\subsection{电磁感应定律及其三种表现形式}

当穿过闭合导体回路的磁通（不论由于什么原因）发生变化时，回路中出现感应电流，说明回路中
必然存在感应电动势。法拉第总结出：闭合回路中的感应电动势与穿过该回路的磁通随时间的变化率
成正比，

\begin{formulaBox}
\[
\mathcal{E}=-\frac{\mathrm{d}\Phi_m}{\mathrm{d}t}
=-\frac{\mathrm{d}}{\mathrm{d}t}\int_{S}\vec{B}\cdot \mathrm{d}\vec{S}
\]
\end{formulaBox}

式中 $S$ 是由闭合回路 $l$ 所限定的面积。\textbf{参考方向的规定}：感应电动势的参考方向与
穿过回路的磁通的参考方向符合右手螺旋关系；面积元的正法线方向与 $l$ 的绕向也符合右手螺旋关系。
负号即楞次定律：感应电流产生的磁场总要阻碍原磁场的变化。

按磁通变化的原因，式中的 $-{\mathrm{d}\Phi_m}/{\mathrm{d}t}$ 可拆成三种情形：

\textbf{（1）感生电动势。}$\vec{B}$ 随时间变化而回路（相对媒质）不动：
\begin{formulaBox}
\[
\mathcal{E}=\oint_{l}\vec{E}_j\cdot \mathrm{d}\vec{l}
=-\int_{S}\frac{\partial \vec{B}}{\partial t}\cdot \mathrm{d}\vec{S}
\]
\end{formulaBox}
变压器即依此工作，故又称\textbf{变压器电动势}。

\textbf{（2）动生电动势。}$\vec{B}$ 恒定而回路整体或局部以速度 $\vec{v}$ 相对媒质运动：
\begin{formulaBox}
\[
\mathcal{E}=\oint_{l}(\vec{v}\times \vec{B})\cdot \mathrm{d}\vec{l}
\]
\end{formulaBox}
这是发电机的工作原理，故又称\textbf{发电机电动势}。其微观解释是导体中的自由电荷受洛伦兹
磁力 $q\vec{v}\times\vec{B}$ 驱动。

\textbf{（3）两者同时存在。}感应电动势为两者的叠加：
$\mathcal{E}=\oint_{l}(\vec{E}+\vec{v}\times\vec{B})\cdot \mathrm{d}\vec{l}$，
右端分别是感生项 $-\int_S\partial\vec{B}/\partial t\cdot \mathrm{d}\vec{S}$ 与动生项
$\oint_l(\vec{v}\times\vec{B})\cdot \mathrm{d}\vec{l}$。

两个层次要分清：\textbf{感应电动势}只取决于磁通变化率，与回路材料无关，可推广到任意假想回路；
\textbf{感应电流}还要由电阻决定，回路不闭合时有电动势而无电流。故感应电动势更本质。

\subsection{感应电场（涡旋电场）}

电动势是非保守电场的环路积分，回路中存在感应电动势，说明回路中存在非保守电场。麦克斯韦据此
假设：\textbf{除了电荷产生电场外，变化的磁场也总要在空间产生电场}，这种电场称为
\textbf{感应电场}，记作 $\vec{E}_j$。

这一假设的关键在于它\textbf{不限于导体回路}：不论空间有无导体、有无回路，不论在真空中还是
媒质中，时变磁场与 $\vec{E}_j$ 都相伴存在。于是

\begin{formulaBox}
\[
\oint_{l}\vec{E}_j\cdot \mathrm{d}\vec{l}
=-\int_{S}\frac{\partial \vec{B}}{\partial t}\cdot \mathrm{d}\vec{S}
\quad\Longrightarrow\quad
\nabla\times \vec{E}_j=-\frac{\partial \vec{B}}{\partial t},
\qquad \nabla\cdot \vec{E}_j=0
\]
\end{formulaBox}

感应电场的环量不为零，是\textbf{非保守场}；其力线是无头无尾的闭合曲线，故又称
\textbf{涡旋电场}；它无源（$\nabla\cdot\vec{E}_j=0$），不始于电荷、不止于电荷。

空间中的总电场是库仑场与感应电场的叠加 $\vec{E}=\vec{E}_C+\vec{E}_j$：

\begin{center}
\begin{tabular}{@{}>{\raggedright\arraybackslash}p{0.45\linewidth}
>{\raggedright\arraybackslash}p{0.45\linewidth}@{}}
\toprule
\textbf{库仑场（电荷激发）} & \textbf{感应电场（变化磁场激发）} \\
\midrule
$\nabla\times \vec{E}_C=0$，$\nabla\cdot \vec{E}_C=\rho/\varepsilon_0$
& $\nabla\times \vec{E}_j=-\partial \vec{B}/\partial t$，$\nabla\cdot \vec{E}_j=0$ \\
有源无旋场，保守场，可引入标量电位 $\vec{E}_C=-\nabla\varphi$
& 有旋无源场（涡旋电场），非保守场，\textbf{不能}引入单值标量电位 \\
力线起于正电荷、止于负电荷 & 力线闭合，无起点无终点 \\
\bottomrule
\end{tabular}
\end{center}

总电场的两个方程为
\begin{formulaBox}
\[
\nabla\times \vec{E}=-\frac{\partial \vec{B}}{\partial t},
\qquad
\nabla\cdot \vec{E}=\frac{\rho}{\varepsilon_0}
\]
\end{formulaBox}
其旋度方程即\textbf{麦克斯韦第二方程}的微分形式，包含“时变磁场激励电场”这一重要信息，
是电磁感应定律的微分形式。

\subsection{全电流定律与位移电流}

把恒定磁场的安培环路定律用到含电容的交流电路上会出现矛盾：取同一闭合曲线 $l$ 所限定的两个
不同曲面 $S_1$（穿过导线）与 $S_2$（穿过电容器极板间），得

\[
\oint_{l}\vec{H}\cdot \mathrm{d}\vec{l}=\int_{S_1}\vec{J}\cdot \mathrm{d}\vec{S}=i,
\qquad
\oint_{l}\vec{H}\cdot \mathrm{d}\vec{l}=\int_{S_2}\vec{J}\cdot \mathrm{d}\vec{S}=0
\]

两式矛盾，根源是\textbf{传导电流在电容器处不连续}，故安培环路定律不能直接用于非恒定电流。

麦克斯韦注意到：极板处传导电流中断，但同时极板上电荷量在变化，极板间电场也在变化。他把电
位移（电通密度）的时间变化率看作一种等效的电流密度，称为\textbf{位移电流密度}：

\begin{formulaBox}
\[
\vec{J}_D=\frac{\partial \vec{D}}{\partial t},
\qquad
\vec{J}=\vec{J}_C+\vec{J}_D=\vec{J}_C+\frac{\partial \vec{D}}{\partial t}
\]
\end{formulaBox}

两者之和称为\textbf{全电流}。由电荷守恒定律 $\oint_S\vec{J}_C\cdot \mathrm{d}\vec{S}
=-\dfrac{\mathrm{d}q}{\mathrm{d}t}$ 与高斯定律 $\oint_S\vec{D}\cdot \mathrm{d}\vec{S}=q$ 可得

\[
\nabla\cdot \vec{J}=\nabla\cdot\left(\vec{J}_C+\frac{\partial \vec{D}}{\partial t}\right)
=\nabla\cdot \vec{J}_C+\frac{\partial \rho}{\partial t}=0
\]

即\textbf{全电流是连续的}。这与电荷守恒定律是等价的（互为充要）。

把安培环路定律修正为全电流定律：

\begin{formulaBox}
\[
\oint_{l}\vec{H}\cdot \mathrm{d}\vec{l}
=\int_{S}\left(\vec{J}_C+\frac{\partial \vec{D}}{\partial t}\right)\cdot \mathrm{d}\vec{S},
\qquad
\nabla\times \vec{H}=\vec{J}_C+\frac{\partial \vec{D}}{\partial t}
\]
\end{formulaBox}

后者即\textbf{麦克斯韦第一方程}的微分形式。其物理内容为：不但传导电流能激发磁场，位移电流
也\textbf{以相同的方式}激发磁场。这就是“变化的电场产生磁场”，与“变化的磁场产生电场”一起，
使电磁场成为相互激发、不可分割的统一体。

\subsection{位移电流与传导电流的对比}

\begin{center}
\begin{tabular}{@{}p{0.22\linewidth}p{0.34\linewidth}p{0.34\linewidth}@{}}
\toprule
\textbf{比较项} & \textbf{传导电流} $\vec{J}_C=\gamma\vec{E}$ & \textbf{位移电流} $\vec{J}_D=\partial\vec{D}/\partial t$ \\
\midrule
起因 & 带电粒子的宏观定向运动 & 电场（电位移）随时间变化 \\
激发磁场 & 相同规律，都满足 $\nabla\times\vec{H}=\vec{J}$ 的形式 & 相同规律，都对磁场有贡献 \\
热效应 & 产生焦耳热 $p=J_C^2/\gamma$ & 真空中不产生焦耳热；介质中由 $\varepsilon$ 变化引起的部分有热效应，但规律不同 \\
连续性 & 一般不连续（电容器处中断） & 与传导电流之和连续 \\
\bottomrule
\end{tabular}
\end{center}

\corepoint{最常错的两处：（1）位移电流\textbf{不是}电荷的运动，真空中它只对应电场的变化，因此
\textbf{不能}写成 $\vec{J}_D=\rho\vec{v}$，也\textbf{不}服从欧姆定律；（2）“全电流连续”说的是
$\vec{J}_C+\vec{J}_D$ 连续，\textbf{不能}因此说传导电流连续。判据上还要记住：导体中
$J_D/J_C=\omega\varepsilon/\gamma$，良导体中这个比值极小，所以低频下导体内可以只算传导电流
（这正是第 5 章磁准静态场的入口）。}

\section{电磁场基本方程组·分界面上的衔接条件}

\subsection{麦克斯韦方程组（积分形式与微分形式）}

把前几章的结论与位移电流假设合起来，就得到概括一切宏观电磁现象的四个方程，即
\textbf{麦克斯韦方程组}：

\begin{formulaBox}
积分形式：
\[
\begin{cases}
\displaystyle\oint_{l}\vec{H}\cdot \mathrm{d}\vec{l}
=\int_{S}\left(\vec{J}_C+\dfrac{\partial \vec{D}}{\partial t}\right)\cdot \mathrm{d}\vec{S}\\[10pt]
\displaystyle\oint_{l}\vec{E}\cdot \mathrm{d}\vec{l}
=-\int_{S}\dfrac{\partial \vec{B}}{\partial t}\cdot \mathrm{d}\vec{S}\\[10pt]
\displaystyle\oint_{S}\vec{B}\cdot \mathrm{d}\vec{S}=0\\[10pt]
\displaystyle\oint_{S}\vec{D}\cdot \mathrm{d}\vec{S}=\int_{V}\rho_f\,\mathrm{d}V
\end{cases}
\qquad\xrightarrow[\text{高斯、斯托克斯}]{\text{散度、旋度定理}}\qquad
\begin{cases}
\nabla\times \vec{H}=\vec{J}_C+\dfrac{\partial \vec{D}}{\partial t}\\[8pt]
\nabla\times \vec{E}=-\dfrac{\partial \vec{B}}{\partial t}\\[8pt]
\nabla\cdot \vec{B}=0\\[8pt]
\nabla\cdot \vec{D}=\rho_f
\end{cases}
\]
\end{formulaBox}

有媒质存在时方程组尚不完备，还须补充\textbf{构成关系（辅助方程）}。对各向同性线性媒质：

\begin{formulaBox}
\[
\vec{D}=\varepsilon\vec{E},\qquad \vec{B}=\mu\vec{H},\qquad \vec{J}_C=\gamma\vec{E}
\]
\end{formulaBox}

与洛伦兹力公式一起构成宏观电磁场理论的完整基础：
$\vec{f}=q\vec{E}+q(\vec{v}\times\vec{B})$。

四个方程的物理意义分别是：第 I 方程 = 全电流定律（传导电流与变化电场都激发磁场）；第 II 方程
= 推广的电磁感应定律（变化磁场激发有旋电场）；第 III 方程 = 磁通连续性（磁力线无头无尾）；
第 IV 方程 = 高斯定律（电荷以发散方式产生电场）。方程组\textbf{不对称}的根源是磁场的无源性：
不存在磁荷，也没有磁流密度。

\subsection{关于方程组的几点说明}

\begin{enumerate}
  \item \textbf{两个旋度方程是核心，第 III 方程不独立。}由第 II 方程取散度：
  $\nabla\cdot(\nabla\times\vec{E})=-\dfrac{\partial}{\partial t}(\nabla\cdot\vec{B})=0$，
  得 $\nabla\cdot\vec{B}=C$；该式在整个时间轴上对任意情况都成立，故 $C=0$。

  \item \textbf{方程组隐含电荷守恒定律。}由第 I、IV 方程取散度：
  $\nabla\cdot(\nabla\times\vec{H})=\nabla\cdot\vec{J}_C
  +\dfrac{\partial}{\partial t}(\nabla\cdot\vec{D})=0$，
  即 $\nabla\cdot\vec{J}_C+\dfrac{\partial\rho_f}{\partial t}=0$。

  \item \textbf{方程组是线性方程组}，满足叠加定理，这是分解为各频率正弦场分别求解的依据。

  \item \textbf{静电场、恒定电场、恒定磁场都是它的特例}：令 $\partial/\partial t=0$ 即退化为
  $\nabla\times\vec{H}=\vec{J}_C$、$\nabla\times\vec{E}=0$、$\nabla\cdot\vec{B}=0$、
  $\nabla\cdot\vec{D}=\rho_f$。

  \item 方程组与物性方程一起\textbf{完备、自洽}；给定电荷电流分布及初始、边界条件，场的
  变化被唯一确定（电磁场的唯一性定理）。

  \item \textbf{电磁场的波动性。}无源自由空间中 $\vec{J}_C=0$、$\rho_f=0$，方程组成为
  $\nabla\times\vec{H}=\partial\vec{D}/\partial t$、$\nabla\times\vec{E}=-\partial\vec{B}/\partial t$、
  $\nabla\cdot\vec{B}=0$、$\nabla\cdot\vec{D}=0$。电场与磁场互为源、相互激发，能量在两者之间
  连续转换，使电磁场像波浪一样由场源向远处传播——这就是\textbf{电磁波}，是无线电通信的基础。
  对各向同性、线性、均匀媒质，可导出
  \textbf{广义波动方程}
  \begin{formulaBox}
  \[
  \nabla^2\vec{E}-\mu\gamma\frac{\partial \vec{E}}{\partial t}
  -\mu\varepsilon\frac{\partial^2\vec{E}}{\partial t^2}=0,
  \qquad
  \nabla^2\vec{H}-\mu\gamma\frac{\partial \vec{H}}{\partial t}
  -\mu\varepsilon\frac{\partial^2\vec{H}}{\partial t^2}=0
  \]
  \end{formulaBox}
  理想介质中 $\gamma=0$，得到\textbf{齐次波动方程}
  \begin{formulaBox}
  \[
  \nabla^2\vec{E}-\frac{1}{v^2}\frac{\partial^2\vec{E}}{\partial t^2}=0,
  \qquad
  \nabla^2\vec{H}-\frac{1}{v^2}\frac{\partial^2\vec{H}}{\partial t^2}=0,
  \qquad v=\frac{1}{\sqrt{\mu\varepsilon}}
  \]
  \end{formulaBox}
  真空中 $v=c=1/\sqrt{\mu_0\varepsilon_0}=3\times10^{8}\,\mathrm{m/s}$；波动是电磁场的基本
  运动形态。
\end{enumerate}

\subsection{分界面上的衔接条件}

把方程组的积分形式用于跨越分界面的扁平小盒与狭长小回路，即得衔接条件。设 $\vec{n}_{12}$ 为
由媒质 1 指向媒质 2 的单位法向矢量，$\vec{K}$ 为分界面上的面电流线密度，$\sigma_f$ 为面电荷密度：

\begin{formulaBox}
\[
\begin{cases}
\vec{n}_{12}\times(\vec{H}_2-\vec{H}_1)=\vec{K}
&\Longleftrightarrow\quad H_{1t}-H_{2t}=K\\[6pt]
\vec{n}_{12}\times(\vec{E}_2-\vec{E}_1)=0
&\Longleftrightarrow\quad E_{1t}=E_{2t}\\[6pt]
\vec{n}_{12}\cdot(\vec{B}_2-\vec{B}_1)=0
&\Longleftrightarrow\quad B_{1n}=B_{2n}\\[6pt]
\vec{n}_{12}\cdot(\vec{D}_2-\vec{D}_1)=\sigma_f
&\Longleftrightarrow\quad D_{2n}-D_{1n}=\sigma_f
\end{cases}
\]
\end{formulaBox}

四种衔接条件恰好对应四个方程：\textbf{切向 $\vec{H}$ 的跃变等于面电流}、\textbf{切向 $\vec{E}$
连续}、\textbf{法向 $\vec{B}$ 连续}、\textbf{法向 $\vec{D}$ 的跃变等于面电荷}。连续介质
内部的方程组加上分界面上的衔接条件，才构成任意区域完整的电磁场方程。

\subsection{两类常见分界面与折射定律}

\textbf{（1）两种理想介质分界面}（$\sigma_f=0$，$\vec{K}=0$）：四个量都连续
\begin{formulaBox}
\[
H_{1t}=H_{2t},\qquad E_{1t}=E_{2t},\qquad B_{1n}=B_{2n},\qquad D_{1n}=D_{2n}
\]
\end{formulaBox}
由此可导出\textbf{折射定律}：设 $\alpha_1,\alpha_2$ 为电场与法线的夹角，$\beta_1,\beta_2$ 为
磁场与法线的夹角，则
\begin{formulaBox}
\[
\frac{\tan\alpha_1}{\tan\alpha_2}=\frac{\varepsilon_1}{\varepsilon_2},
\qquad
\frac{\tan\beta_1}{\tan\beta_2}=\frac{\mu_1}{\mu_2}
\]
\end{formulaBox}

\textbf{（2）理想导体表面}（导体内部无时变场，$\vec{E}_1=\vec{B}_1=0$）：
\begin{formulaBox}
\[
\vec{n}_{12}\times\vec{H}_2=\vec{K},\qquad
\vec{n}_{12}\times\vec{E}_2=0,\qquad
\vec{n}_{12}\cdot\vec{B}_2=0,\qquad
\vec{n}_{12}\cdot\vec{D}_2=\sigma_f
\]
\end{formulaBox}
即：理想导体表面\textbf{切向电场为零、法向磁场为零}，电力线垂直于导体表面，磁力线平行于导体
表面，电流以面电流形式在表面流动。

\corepoint{衔接条件三个易错点：（1）$\vec{n}_{12}$ 的方向必须“由 1 指向 2”，反过来写会让
$\sigma_f$、$K$ 的符号全错；（2）连续的是 $E_t$ 与 $B_n$，\textbf{不是} $E_n$ 与 $B_t$
（$D_n$ 只在无面电荷时连续）；（3）时变场里即使导体“理想”，其表面仍可有 $\sigma_f$ 与 $\vec{K}$，
不能把“导体内部无场”误推成“导体表面无电荷”。}

\section{动态位及其积分解}

\subsection{位函数的引入}

直接解 $\vec{E}$、$\vec{B}$ 的耦合方程组很困难，因此引入辅助位函数把问题化为波动方程的求解。
由 $\nabla\cdot\vec{B}=0$ 可令

\begin{formulaBox}
\[
\vec{B}=\nabla\times\vec{A},
\qquad
\vec{E}=-\frac{\partial \vec{A}}{\partial t}-\nabla\varphi
\]
\end{formulaBox}

$\vec{A}$ 称为\textbf{矢量位函数}（动态矢量位），$\varphi$ 称为\textbf{标量位函数}
（动态标量位）。注意 $\vec{E}$ 的表达式比静电场多出 $-\partial\vec{A}/\partial t$ 一项：
时变场中电场不再单由电荷产生，\textbf{不能}再写 $\vec{E}=-\nabla\varphi$。

\subsection{达朗贝尔方程与洛伦兹条件}

把上式代入第 IV、第 I 方程，在线性、各向同性、均匀媒质中得

\begin{formulaBox}
\[
\nabla^2\varphi+\frac{\partial}{\partial t}(\nabla\cdot\vec{A})=-\frac{\rho_f}{\varepsilon},
\qquad
\nabla^2\vec{A}-\mu\varepsilon\frac{\partial^2\vec{A}}{\partial t^2}
=-\mu\vec{J}_C+\nabla\left(\nabla\cdot\vec{A}+\mu\varepsilon\frac{\partial\varphi}{\partial t}\right)
\]
\end{formulaBox}

两式通过 $\nabla\cdot\vec{A}$ 与 $\partial\varphi/\partial t$ 耦合，联立求解困难。由于位函数
本身有选择的自由度（规范不变性），可选取\textbf{洛伦兹条件}
\begin{formulaBox}
\[
\nabla\cdot\vec{A}+\mu\varepsilon\frac{\partial\varphi}{\partial t}=0
\]
\end{formulaBox}
使方程解耦，得到一对形式完全对称的\textbf{达朗贝尔方程}：
\begin{formulaBox}
\[
\begin{cases}
\displaystyle\nabla^2\vec{A}-\mu\varepsilon\frac{\partial^2\vec{A}}{\partial t^2}=-\mu\vec{J}_C\\[10pt]
\displaystyle\nabla^2\varphi-\mu\varepsilon\frac{\partial^2\varphi}{\partial t^2}
=-\frac{\rho_f}{\varepsilon}
\end{cases}
\]
\end{formulaBox}

达朗贝尔方程是\textbf{非齐次波动方程}：等号左边的 $\mu\varepsilon\,\partial^2/\partial t^2$ 项
正是波动性的来源，去掉它就是泊松方程（即恒定场的情形）。

\subsection{达朗贝尔方程的特解——推迟位}

在无限大均匀媒质中，达朗贝尔方程的特解为

\begin{formulaBox}
\[
\vec{A}(\vec{r},t)=\frac{\mu}{4\pi}\int_{V'}
\frac{\vec{J}_C\left(\vec{r}',\,t-\dfrac{R}{v}\right)}{R}\,\mathrm{d}V',
\qquad
\varphi(\vec{r},t)=\frac{1}{4\pi\varepsilon}\int_{V'}
\frac{\rho_f\left(\vec{r}',\,t-\dfrac{R}{v}\right)}{R}\,\mathrm{d}V'
\]
\end{formulaBox}
$R=|\vec{r}-\vec{r}'|$ 为场点到源点的距离，$v=1/\sqrt{\mu\varepsilon}$ 为电磁波的传播速度。

这两个解称为\textbf{推迟位}（滞后位）。其物理意义：场点 $P$ 在时刻 $t$ 的位，不是由 $t$ 时刻
的源分布决定，而是由 $t'=t-R/v$ 时刻的源分布决定；时间差 $\Delta t=R/v$ 正是电磁作用以速度
$v$ 传播距离 $R$ 所需的时间。可见电磁扰动以\textbf{有限速度}传播，不存在超距作用；空间各点的
位总是滞后于激励源，电磁波由近及远向外传播，波速由媒质特性决定。

\corepoint{洛伦兹条件与库仑规范别混：洛伦兹条件
$\nabla\cdot\vec{A}+\mu\varepsilon\partial\varphi/\partial t=0$ 是为解耦达朗贝尔方程而选，
适用于时变场；库仑规范 $\nabla\cdot\vec{A}=0$ 适用于恒定磁场（此时
$\nabla^2\vec{A}=-\mu\vec{J}_C$、$\nabla^2\varphi=-\rho_f/\varepsilon$ 退化为泊松方程）。
另外“推迟位”只对无限大均匀媒质成立，有边界面时必须叠加满足边界条件的解。}

\section{电磁功率流和坡印亭矢量}

\subsection{坡印亭定理}

设均匀各向同性媒质中，电磁场的能量密度为电场能量密度与磁场能量密度之和
\begin{formulaBox}
\[
w=w_e+w_m=\frac{1}{2}\vec{E}\cdot\vec{D}+\frac{1}{2}\vec{B}\cdot\vec{H}
\quad(\mathrm{J/m^3}),
\qquad
W=\int_{V}w\,\mathrm{d}V
\]
\end{formulaBox}

求 $W$ 的时间变化率，用物性方程与第 I、II 方程消去 $\partial\vec{D}/\partial t$、
$\partial\vec{B}/\partial t$，再用恒等式
$\nabla\cdot(\vec{E}\times\vec{H})=\vec{H}\cdot(\nabla\times\vec{E})
-\vec{E}\cdot(\nabla\times\vec{H})$ 与高斯散度定理，得\textbf{坡印亭定理}：

\begin{formulaBox}
\[
-\frac{\partial W}{\partial t}
=\oint_{A}(\vec{E}\times\vec{H})\cdot \mathrm{d}\vec{A}
+\int_{V}\vec{E}\cdot\vec{J}_C\,\mathrm{d}V
\]
\end{formulaBox}

它表示：区域内单位时间减少的电磁能量 $-\partial W/\partial t$，一部分通过包围该区域的闭合面
$A$ 流出（$\oint_A\vec{S}\cdot \mathrm{d}\vec{A}$），一部分在区域内转化为焦耳热
（$\int_V\vec{E}\cdot\vec{J}_C\,\mathrm{d}V$，即 $\int_V J_C^2/\gamma\,\mathrm{d}V$）。
这是\textbf{能量守恒与转换定律在电磁场中的具体表现}。

\subsection{含电源的坡印亭定理}

若区域内还有电源（局外场 $\vec{E}_e$ 作用），则要加入电源提供的功率项：

\begin{formulaBox}
\[
\oint_{A}(\vec{E}\times\vec{H})\cdot \mathrm{d}\vec{A}
=-\frac{\partial W}{\partial t}
-\int_{V}\frac{J_C^{2}}{\gamma}\,\mathrm{d}V
+\int_{V}\vec{J}\cdot\vec{E}_e\,\mathrm{d}V
\]
\end{formulaBox}

式中各项依次为：向外输送的电磁功率、储能的变化、焦耳热损耗、电源提供的功率。注意“电源”项
对应非静电力做功，与 $-\partial W/\partial t$ 符号方向相反。

\subsection{坡印亭矢量与典型能流}

定义\textbf{坡印亭矢量}（能流密度矢量、功率流密度矢量）
\begin{formulaBox}
\[
\vec{S}=\vec{E}\times\vec{H}\qquad(\mathrm{W/m^2})
\]
\end{formulaBox}
大小表示单位时间内通过垂直于能量传播方向的单位面积的电磁能量，方向即能量流动方向。

\textbf{（1）长直同轴电缆。}内外导体均为理想导体时，两导体之间为真空，设电压 $U$、电流 $I$，
则
\[
\vec{E}=\frac{U}{\rho\ln(b/a)}\vec{e}_\rho,\qquad
\vec{H}=\frac{I}{2\pi\rho}\vec{e}_\varphi
\quad\Longrightarrow\quad
\vec{S}=\frac{UI}{2\pi\rho^{2}\ln(b/a)}\vec{e}_z
\]
沿横截面 $A$ 积分得
$P=\int_{A}\vec{S}\cdot \mathrm{d}\vec{A}=UI$，即传输功率恰等于电路理论中的 $UI$。结论：
\textbf{电磁能量在导体之间的空间沿轴线流动}，导体（理想）本身不传输能量，其作用是引导场、
约束能量的流动方向。

\textbf{（2）导体非理想时。}内导体内部存在轴向电场
$E_z=I/(\pi a^{2}\gamma)$，于是坡印亭矢量有指向导体内部的径向分量 $S_\rho<0$，这部分能量
进入导体并转化为焦耳热，单位长度损耗功率为 $I^{2}r$。这正是“损耗
功率是从场流入导体的”这一图景。

\textbf{（3）圆柱形平板电容器缓慢充电。}内部能量增加所需的功率是从\textbf{侧面}由场流入的，
而非经极板导线流入——导线不是能量载体，能量由场传输。

\corepoint{坡印亭矢量的两个常见误用：（1）$\vec{S}$ 是\textbf{能流密度}，某点的 $\vec{S}$ 并不
等于该点“消耗的功率”，判断某区域获得/失去功率必须做面积分；（2）$\vec{S}=\vec{E}\times\vec{H}$
不代表能量一定沿导线方向流动——在同轴电缆中能量走的是导体\textbf{之间}的空间，在电阻器中
能量从侧面\textbf{流入}电阻体。}

\section{正弦电磁场}

\subsection{复数表示法}

以一定频率作正弦变化的场称为\textbf{正弦电磁场}。一般非正弦时变场可用傅里叶级数分解为稳态
分量与各频率正弦分量分别研究，因此正弦场是基本工具。

设电场各分量同频率、不同初相：
\[
\vec{E}(x,y,z;t)=E_{xm}\cos(\omega t+\phi_x)\vec{e}_x
+E_{ym}\cos(\omega t+\phi_y)\vec{e}_y
+E_{zm}\cos(\omega t+\phi_z)\vec{e}_z
\]
引入\textbf{电场强度复矢量}（有效值相量）
\begin{formulaBox}
\[
\vec{E}(x,y,z;t)=\mathrm{Re}\left[\sqrt{2}\,\dot{\vec{E}}(x,y,z)e^{j\omega t}\right],
\qquad
\dot{\vec{E}}=\frac{1}{\sqrt{2}}\left(E_{xm}e^{j\phi_x}\vec{e}_x
+E_{ym}e^{j\phi_y}\vec{e}_y+E_{zm}e^{j\phi_z}\vec{e}_z\right)
\]
\end{formulaBox}
其中 $\sqrt{2}$ 的约定使 $|\dot{\vec{E}}|$ 为\textbf{有效值}。复数法的最大好处是对时间的求导
化为乘 $j\omega$：
\begin{formulaBox}
\[
\frac{\partial}{\partial t}\ \longleftrightarrow\ j\omega,
\qquad
\frac{\partial^2}{\partial t^2}\ \longleftrightarrow\ -\omega^{2}
\]
\end{formulaBox}

\subsection{麦克斯韦方程组的复数形式}

把 $\partial/\partial t\to j\omega$ 代入微分形式方程组，得
\begin{formulaBox}
\[
\begin{cases}
\nabla\times\dot{\vec{H}}=\dot{\vec{J}}_C+j\omega\dot{\vec{D}}\\[6pt]
\nabla\times\dot{\vec{E}}=-j\omega\dot{\vec{B}}\\[6pt]
\nabla\cdot\dot{\vec{B}}=0\\[6pt]
\nabla\cdot\dot{\vec{D}}=\dot{\rho}_f
\end{cases}
\qquad
\dot{\vec{D}}=\varepsilon\dot{\vec{E}},\quad
\dot{\vec{B}}=\mu\dot{\vec{H}},\quad
\dot{\vec{J}}_C=\gamma\dot{\vec{E}}
\]
\end{formulaBox}
第一式可写成
$\nabla\times\dot{\vec{H}}=(\gamma+j\omega\varepsilon)\dot{\vec{E}}$：括号内的
$\gamma+j\omega\varepsilon$ 把传导电流与位移电流合在一起，是分析有损耗媒质的出发点
（第 5 章判断“良导体/良介质”就是比较这两项的大小）。

\subsection{复坡印亭矢量与坡印亭定理的复数形式}

定义\textbf{复坡印亭矢量}
\begin{formulaBox}
\[
\dot{\vec{S}}=\dot{\vec{E}}\times\dot{\vec{H}}^{*},
\qquad
\vec{S}_{av}=\mathrm{Re}\left[\dot{\vec{E}}\times\dot{\vec{H}}^{*}\right]
=\frac{1}{T}\int_{0}^{T}\vec{S}(t)\,\mathrm{d}t
\]
\end{formulaBox}
其实部是坡印亭矢量的平均值（有功功率密度，表示能量流动），虚部是无功功率密度（表示电磁能量
的交换）。对 $\dot{\vec{S}}$ 取散度并利用复数形式的第 I、II 方程，再体积分、用散度定理，得
\textbf{坡印亭定理的复数形式}：
\begin{formulaBox}
\[
-\oint_{A}\left(\dot{\vec{E}}\times\dot{\vec{H}}^{*}\right)\cdot \mathrm{d}\vec{A}
=\int_{V}\dot{\vec{J}}_C^{*}\cdot\dot{\vec{E}}\,\mathrm{d}V
+j\omega\int_{V}\left(\dot{\vec{B}}\cdot\dot{\vec{H}}^{*}
-\dot{\vec{E}}\cdot\dot{\vec{D}}^{*}\right)\mathrm{d}V
\]
\end{formulaBox}
与电路中的复功率 $\tilde{S}=P+j(Q_L-Q_C)$ 对照，若体积内不含电源：
\begin{formulaBox}
\[
\tilde{S}=-\oint_{A}\left(\dot{\vec{E}}\times\dot{\vec{H}}^{*}\right)\cdot \mathrm{d}\vec{A}
=P+jQ
\]
\[
P=\mathrm{Re}[\tilde{S}]=\int_{V}\dot{\vec{J}}_C^{*}\cdot\dot{\vec{E}}\,\mathrm{d}V
\ \ \text{（焦耳热损耗的平均值，即有功功率）}
\]
\[
Q=\mathrm{Im}[\tilde{S}]
=\omega\int_{V}\left(\dot{\vec{B}}\cdot\dot{\vec{H}}^{*}
-\dot{\vec{E}}\cdot\dot{\vec{D}}^{*}\right)\mathrm{d}V
\ \ \text{（电磁能量的交换，即无功功率）}
\]
\end{formulaBox}

按“等效电路”的观点令 $P=I^{2}R$、$Q=I^{2}X$，可得时变场中体积 $V$ 内媒质的\textbf{等效电路
参数}
\begin{formulaBox}
\[
R=\frac{1}{I^{2}}\,\mathrm{Re}\left[-\oint_{A}
\left(\dot{\vec{E}}\times\dot{\vec{H}}^{*}\right)\cdot \mathrm{d}\vec{A}\right],
\qquad
X=\frac{1}{I^{2}}\,\mathrm{Im}\left[-\oint_{A}
\left(\dot{\vec{E}}\times\dot{\vec{H}}^{*}\right)\cdot \mathrm{d}\vec{A}\right]
\]
\end{formulaBox}
$R$ 反映介质/导体的损耗，$X$ 反映体内磁场储能（内电抗）。这两个式子是第 5 章“导体的交流
内阻抗”的计算依据。

\subsection{达朗贝尔方程的复数形式与滞后位}

正弦稳态下达朗贝尔方程化为\textbf{非齐次亥姆霍兹方程}
\begin{formulaBox}
\[
\begin{cases}
\nabla^{2}\dot{\vec{A}}+\beta^{2}\dot{\vec{A}}=-\mu\dot{\vec{J}}_C\\[6pt]
\nabla^{2}\dot{\varphi}+\beta^{2}\dot{\varphi}=-\dfrac{\dot{\rho}_f}{\varepsilon}
\end{cases}
\qquad
\beta=\omega\sqrt{\mu\varepsilon}\ \ (\mathrm{rad/m})
\]
\end{formulaBox}
$\beta$ 称为\textbf{相位常数}（波数）。时间上推迟 $\Delta t=R/v$ 相当于相位上滞后
$\omega\Delta t=\omega R/v=\beta R$，故解（\textbf{滞后位}）为
\begin{formulaBox}
\[
\dot{\vec{A}}=\frac{\mu}{4\pi}\int_{V'}
\frac{\dot{\vec{J}}_C\,e^{-j\beta R}}{R}\,\mathrm{d}V',
\qquad
\dot{\varphi}=\frac{1}{4\pi\varepsilon}\int_{V'}
\frac{\dot{\rho}_f\,e^{-j\beta R}}{R}\,\mathrm{d}V'
\]
\end{formulaBox}
因子 $e^{-j\beta R}$ 表示场点比源点滞后 $\beta R$ 的相位，\textbf{它的存在标志着场的波动性}；
若 $\beta R\ll 1$（即 $R\ll\lambda$），可令 $e^{-j\beta R}\approx1$，滞后位退化为静态场的
位函数表达式。

场量与位的关系、洛伦兹条件在复数形式下为
\begin{formulaBox}
\[
\dot{\vec{B}}=\nabla\times\dot{\vec{A}},
\qquad
\dot{\vec{E}}=-j\omega\dot{\vec{A}}-\nabla\dot{\varphi}
=-j\omega\dot{\vec{A}}+\frac{\nabla(\nabla\cdot\dot{\vec{A}})}{j\omega\mu\varepsilon},
\qquad
\nabla\cdot\dot{\vec{A}}+j\omega\mu\varepsilon\dot{\varphi}=0
\]
\end{formulaBox}
即只要求出 $\dot{\vec{A}}$，就能算出 $\dot{\vec{E}}$ 与 $\dot{\vec{H}}$。

\corepoint{三个高频错点：（1）本书用\textbf{有效值}相量，所以
$\vec{S}_{av}=\mathrm{Re}[\dot{\vec{E}}\times\dot{\vec{H}}^{*}]$ \textbf{不带} $1/2$；
若改用幅值相量则必须补 $1/2$，两种约定不能混用。（2）复坡印亭定理不是由时域形式逐项“翻译”
来的，它与时域形式没有一一对应关系：时域形式讲能量守恒与转换，复数形式讲有功、无功与复功率
守恒。（3）只有实部是平均功率，虚部是能量交换，\textbf{不能}把 $|\dot{\vec{S}}|$ 当成平均功率
密度。}

\section{电磁辐射}

\subsection{电磁辐射与单元偶极子天线}

电磁能量脱离场源而单独存在于空间中的现象称为\textbf{电磁辐射}。只要有随时间变化的电流、
电荷，就会产生辐射；辐射过程形成电磁波并以有限速度传播。天线就是一种专门的辐射器。

实际的线天线可看作许多\textbf{单元偶极子}（电基本振子）串联而成：一段载流细导线，其长度与
横截面尺寸都远小于波长，且观察点到导线上各点距离近似相同。故其上电流可视为等幅同相，
$\dot{I}$ 为有效值。

单元偶极子位于坐标原点、沿 $z$ 轴放置时，其矢量位为
\begin{formulaBox}
\[
\dot{\vec{A}}(\vec{r})=\frac{\mu_0\dot{I}\Delta l}{4\pi r}e^{-j\beta r}\vec{e}_z
\]
\end{formulaBox}
代入 $\dot{\vec{H}}=\nabla\times\dot{\vec{A}}/\mu_0$、
$\dot{\vec{E}}=-j\omega\dot{\vec{A}}+\nabla(\nabla\cdot\dot{\vec{A}})/(j\omega\mu_0\varepsilon_0)$，
得自由空间中的辐射场（球坐标）：三个分量 $\dot{E}_r$、$\dot{E}_\theta$、$\dot{H}_\varphi$
各含 $1/(\beta r)$、$1/(\beta r)^{2}$、$1/(\beta r)^{3}$ 三类项，且
$\dot{E}_\varphi=\dot{H}_r=\dot{H}_\theta=0$。$\beta r$ 很小时高次幂项占主导，$\beta r$ 很大时
$1/(\beta r)$ 项占主导，故必须分区讨论（\textbf{轻推导、重结论}）。

\subsection{近区场与远区场}

\textbf{（1）近区场（$r\ll\lambda$，即 $\beta r\ll1$）。}高次幂项占主导，$e^{-j\beta r}\approx1$：
\begin{formulaBox}
\[
\dot{E}_r=\frac{\dot{p}}{2\pi\varepsilon_0 r^{3}}\cos\theta,
\qquad
\dot{E}_\theta=\frac{\dot{p}}{4\pi\varepsilon_0 r^{3}}\sin\theta,
\qquad
\dot{H}_\varphi=\frac{\dot{I}\Delta l}{4\pi r^{2}}\sin\theta,
\qquad
\dot{p}=\frac{\dot{I}\Delta l}{j\omega}
\]
\end{formulaBox}
特点：磁场与用毕奥—沙伐定律算出的电流元磁场相同，电场与用库仑定律算出的电偶极子电场相同，
场与源\textbf{同相位}——即推迟效应可忽略，时变场的分布规律与恒定场完全相同。从能量看，
电场与磁场相位相差 $90^{\circ}$，复坡印亭矢量只有虚部，平均能流为零：近区内只有电能与磁能的
交换和振荡，\textbf{不向外净辐射能量}（该结论是忽略高次项的近似结果）。

\textbf{（2）远区场（$r\gg\lambda$，即 $\beta r\gg1$）。}只保留 $1/(\beta r)$ 项：
\begin{formulaBox}
\[
\dot{E}_\theta=j\sqrt{\frac{\mu_0}{\varepsilon_0}}
\frac{\dot{I}\Delta l\sin\theta}{2\lambda r}e^{-j\beta r},
\qquad
\dot{H}_\varphi=j\frac{\dot{I}\Delta l\sin\theta}{2\lambda r}e^{-j\beta r},
\qquad
\dot{E}_r=\dot{E}_\varphi=\dot{H}_r=\dot{H}_\theta=0
\]
\end{formulaBox}
远区场的四个特征：
\begin{itemize}
  \item \textbf{行波性}：相位随 $r$ 增大不断滞后，推迟效应不可忽略，场具有波的性质；
  \item \textbf{球面波}：等相面是以源为中心的球面；
  \item \textbf{横电磁波}：$\vec{E}\perp\vec{H}\perp$ 传播方向，且满足右手螺旋关系，且
  \item \textbf{波阻抗为实数}：$Z_0=\dot{E}_\theta/\dot{H}_\varphi
  =\sqrt{\mu_0/\varepsilon_0}\approx377\,\Omega$，电场与磁场同相位，复坡印亭矢量只有实部，
  沿 $r$ 方向不断向外辐射能量。故远区又称\textbf{辐射区}，其中的场称\textbf{辐射场}。
\end{itemize}

\subsection{辐射功率、辐射电阻与方向性}

穿过半径为 $r$ 的球面向外辐射的总平均功率为
\begin{formulaBox}
\[
P=\oint_{A}\mathrm{Re}(\dot{\vec{S}})\cdot \mathrm{d}\vec{A}
=80\pi^{2}\left(\frac{\Delta l}{\lambda}\right)^{2}I^{2}=I^{2}R_e,
\qquad
R_e=80\pi^{2}\left(\frac{\Delta l}{\lambda}\right)^{2}
\]
\end{formulaBox}
$P$ 与 $r$ 无关，说明能量不在空间停留，而是不断由波源向外辐射。与 $P=I^{2}R$ 对照，把
$R_e$ 称为\textbf{辐射电阻}，它表征天线辐射电磁能量的能力：$R_e$ 越大，辐射能力越强。

远区场的振幅还与观察方向有关：$\dot{E}_\theta,\dot{H}_\varphi\propto\sin\theta$，即
\textbf{方向性因子} $F(\theta)=\sin\theta$。$\theta=90^{\circ}$（垂直于偶极子轴）辐射最强，
$\theta=0^{\circ}$ 或 $180^{\circ}$（沿轴线）无辐射；与 $\varphi$ 无关，故 H 面方向图是圆。
以 $F(\theta)$ 绘制的图形称\textbf{方向图}；多个相同单元按规律排列构成\textbf{天线阵}，其
方向性函数为单元因子与阵因子之积（方向性相乘原理），可获得更强方向性。

\corepoint{辐射部分三个易错点：（1）\textbf{辐射电阻不是损耗电阻}，它不代表发热，而是把辐射出去
的功率“等效”为电路中的电阻消耗；（2）近区“平均能流为零”是忽略了 $1/(\beta r)$ 等高次项的
\textbf{近似}结论，真正的辐射恰恰来自这些被忽略的项，故不能说“近区绝对不辐射”；（3）波阻抗
$377\,\Omega$ 只对\textbf{远区、自由空间}成立，近区 $\dot{E}_\theta/\dot{H}_\varphi$ 是复数且
随 $r$ 变化，不能套用。}

% ============================================================
% 准静态电磁场（子文件，只放内容：不重复 \usepackage、不改字号）
% 对应教材：冯慈璋·马西奎《工程电磁场导论》第 5 章 §5-1 ～ §5-7
% 本章课件无单独讲义，内容以教材为准
% ============================================================

% formulaBox 由 3.矢量分析.tex 定义；这里做一次“未定义才定义”的保护。
\makeatletter
\@ifundefined{formulaBox}{%
  \newenvironment{formulaBox}{%
    \par\begin{adjustbox}{max width=\linewidth}%
    \begin{tcolorbox}[colback=blue!8!white,colframe=transparent,boxrule=0pt,halign=center,boxsep=2pt,left=4pt,right=4pt,top=3pt,bottom=3pt]%
  }{%
    \end{tcolorbox}\end{adjustbox}\par
  }%
}{}
\makeatother

\chapter{准静态电磁场}

时变场中，若感应电场远小于库仑电场（可忽略 $\partial\vec{B}/\partial t$），称\textbf{电准静态
场}（EQS）；若位移电流远小于传导电流（可忽略 $\partial\vec{D}/\partial t$），称\textbf{磁
准静态场}（MQS）。两者统称\textbf{准静态电磁场}：\textbf{属时变场，却保留静态场与源“瞬时
对应”的性质}（略去了波动性与推迟效应）。电气工程中的问题绝大部分属 MQS。

\section{电准静态场与磁准静态场}

\subsection{电准静态场（EQS）}

忽略电磁感应项 $\partial\vec{B}/\partial t$，即 $\vec{E}_j\ll\vec{E}_C$ 时，电场无旋：
$\nabla\times\vec{E}\approx\nabla\times\vec{E}_C=0$。其基本方程组为

\begin{formulaBox}
\[
\nabla\times\vec{E}=0,
\qquad
\nabla\cdot\vec{D}=\rho_f,
\qquad
\nabla\times\vec{H}=\vec{J}_C+\frac{\partial\vec{D}}{\partial t},
\qquad
\nabla\cdot\vec{B}=0
\]
\end{formulaBox}

与静态情形相比，\textbf{只有磁场方程变了}（计入位移电流产生的磁场），电场方程与静电场完全
相同；只是 $\vec{E}$、$\vec{D}$ 随时间变化，但与源\textbf{瞬时对应}。故已知电荷分布即可用静电
场公式求 $\vec{E}$、$\vec{D}$，再由全电流定律求 $\vec{H}$。典型对象：高压电场、电子器件附近
电场、低频交流平板电容器内的电场。

放宽情形：若 $\vec{E}_j$ 不小但其\textbf{旋度}很小，式仍成立，也可按 EQS 处理（如低频电感
线圈导线中的电场可按恒定电场处理）。

\subsection{磁准静态场（MQS）}

忽略位移电流项 $\partial\vec{D}/\partial t$，磁场按恒定磁场处理：

\begin{formulaBox}
\[
\nabla\times\vec{H}=\vec{J}_C,
\qquad
\nabla\cdot\vec{B}=0,
\qquad
\nabla\times\vec{E}=-\frac{\partial\vec{B}}{\partial t},
\qquad
\nabla\cdot\vec{D}=\rho_f
\]
\end{formulaBox}

与静态情形相比，\textbf{只有电场方程变了}（计入电磁感应），磁场方程与恒定磁场相同。与恒定
磁场一样，$\vec{B}=\nabla\times\vec{A}$、$\vec{E}=-\partial\vec{A}/\partial t-\nabla\varphi$，
且在库仑规范 $\nabla\cdot\vec{A}=0$ 下

\begin{formulaBox}
\[
\nabla^{2}\vec{A}=-\mu\vec{J}_C,
\qquad
\nabla^{2}\varphi=-\frac{\rho_f}{\varepsilon}
\]
\end{formulaBox}

即 $\vec{A}$、$\varphi$ 随时间变却遵循静态场（泊松方程）的规律：已知电流、电荷分布即可用静态
公式求位。即可略去场的波动性，认为场与源具有类似静态场的“瞬时”对应关系，故 MQS 又称
\textbf{似稳场}。典型对象：电机、变压器、感应加热、磁悬浮的磁场，及电工中的涡流问题。

\subsection{似稳条件}

\textbf{（1）导体中的时变场}：若 $\dfrac{\omega\varepsilon}{\gamma}\ll1$，位移电流可略，导体内
时变场按 MQS 处理。满足此条件的导体称\textbf{良导体}，其中的 MQS 又称\textbf{涡流场}；对纯
金属该条件一直到紫外波段都成立。

\textbf{（2）理想介质中的时变场}：与源相距 $R$ 处的相位因子为 $e^{-j\beta R}$，要使其近似为 1，
须
\[
\beta R=\frac{2\pi R}{\lambda}\ll1
\quad\Longleftrightarrow\quad
R\ll\lambda
\]
满足该条件的区域称\textbf{近区}（似稳区），其中场与源同步变化、无相位移。两条件统称
\textbf{MQS 近似条件}（似稳条件）。它给出重要判据：\textbf{能否用“路”的方法处理，取决于尺寸
与波长之比，而非绝对尺寸或频率高低}——工频 $50\,\mathrm{Hz}$ 时 $\lambda\approx6000\,\mathrm{km}$，
只有数百公里的输电才需计波动；$3\,\mathrm{GHz}$ 时 $\lambda=10\,\mathrm{cm}$，手掌大的系统
就须计波动。似稳区是\textbf{相对}概念。

\corepoint{\textbf{EQS 丢掉的是“变化磁场激发电场”这条通路（$\partial\vec{B}/\partial t$），
保留位移电流；MQS 丢掉的是“变化电场激发磁场”这条通路（$\partial\vec{D}/\partial t$），保留
电磁感应。}两者\textbf{不能同时}用，否则退化为静态场、丢掉全部耦合与波动信息。}

\section{磁准静态场和电路}

MQS 方程是交流电路的理论基础，基尔霍夫两条定律均可由它导出。

\textbf{（1）基尔霍夫电流定律。}对 $\nabla\times\vec{H}=\vec{J}_C$ 取散度得
$\oint_S\vec{J}_C\cdot \mathrm{d}\vec{S}=0$，即 MQS 中传导电流\textbf{连续}；用于包围节点的
闭合面即得 $\sum i_k=0$。

\textbf{（2）基尔霍夫电压定律。}由 $\vec{J}_C=\gamma(\vec{E}+\vec{E}_e)$ 得
$\vec{E}=\vec{J}_C/\gamma-\vec{E}_e$，代入
$\vec{E}=-\partial\vec{A}/\partial t-\nabla\varphi$ 沿导线由 $A$ 到 $B$ 积分，四项依次对应
电动势、电阻压降、电感压降与电容电压：

\begin{formulaBox}
\[
e(t)=\int_{A}^{B}\vec{E}_e\cdot \mathrm{d}\vec{l}
=i\left(R+r+R_i\right)+L\frac{\mathrm{d}i}{\mathrm{d}t}+u_C
\]
\end{formulaBox}
即“沿闭合回路的电动势等于各段电压降之和”。可见电路理论不过是 $R\ll\lambda$、
$\omega\varepsilon\ll\gamma$ 条件下麦克斯韦理论的近似。

\corepoint{时变场中“电压”只在\textbf{同一横截面上的两点间}才有意义：在纵剖面上取闭合路径，
$\oint_l\vec{E}\cdot \mathrm{d}\vec{l}\neq0$（所围面上有磁通）；只有取在同一横截面内的闭合
路径，才因所围面无磁通而等于零。}

\section{电准静态场与电荷弛豫}

在 $\gamma$、$\varepsilon$ 均匀的导体内，由电荷守恒定律、$\vec{J}_C=\gamma\vec{E}$ 与
$\nabla\cdot\vec{E}=\rho/\varepsilon$ 得

\begin{formulaBox}
\[
\frac{\partial\rho}{\partial t}+\frac{\gamma}{\varepsilon}\rho=0
\quad\Longrightarrow\quad
\rho(t)=\rho_0e^{-t/\tau},
\qquad
\tau=\frac{\varepsilon}{\gamma}
\]
\end{formulaBox}

自由电荷体密度按指数规律衰减，$\tau$ 称\textbf{弛豫时间}。良导体的 $\tau$ 极小，故可认为
\textbf{良导体内部无自由电荷积累}（$\rho=0$），电荷只能分布在表面——与静电场结论一致，但这里
是\textbf{动态}得到的。

EQS 下仍可定义 $\vec{E}=-\nabla\varphi$，代入高斯定律得

\begin{formulaBox}
\[
\nabla^{2}\varphi=-\frac{\rho}{\varepsilon}
=-\frac{\rho_0}{\varepsilon}e^{-t/\tau},
\qquad
\varphi(\vec{r},t)=\varphi(\vec{r},0)\,e^{-t/\tau}
\]
\end{formulaBox}
（后式为无限大均匀导电媒质中的解。）即电位分布按同一指数规律衰减，快慢同取决于 $\tau$；对点
电荷 $\varphi=\dfrac{q_0}{4\pi\varepsilon R}e^{-t/\tau}$，而磁场处处为零。

\textbf{分片均匀导体。}电荷趋于聚积在两导体分界面上，界面两侧须同时满足衔接条件与连续性条件：
\[
E_{1t}=E_{2t},
\qquad
D_{2n}-D_{1n}=\sigma,
\qquad
J_{1n}=J_{2n}\ \Longrightarrow\ \gamma_1E_{1n}=\gamma_2E_{2n}
\]
以双层有损介质平板电容器接直流电压为例：$t=0^+$ 时界面电荷来不及建立，电场按\textbf{介电
常数}分配（如两层理想介质）；$t\to\infty$ 稳态仅有传导电流，电场按\textbf{电导率}分配（如两层
导电媒质）；其间按指数过渡。工程结论：低频（工频）交流下若位移电流远大于漏电流，电场按
$\varepsilon$ 分布，属静电场问题；直流下按 $\gamma$ 分布，属恒定电场问题，两种稳态一般不同。

\corepoint{判据两种写法等价：$\omega\varepsilon/\gamma\ll1$ 与 $\omega\tau\ll1$
（$\tau=\varepsilon/\gamma$），前者从“位移电流/传导电流”比值得出，后者从“周期/弛豫时间”得出。
$\gamma$ 是与 $\omega$、$\varepsilon$ \textbf{相对}比较的：同一材料低频是良导体，到光频可能
变成良介质。}

\section{集肤效应}

MQS 下略去位移电流，对 $\nabla\times\vec{H}=\vec{J}_C$ 取旋度并利用矢量恒等式与物性方程，得

\begin{formulaBox}
\[
\nabla^{2}\vec{H}=\mu\gamma\frac{\partial\vec{H}}{\partial t},
\quad
\nabla^{2}\vec{E}=\mu\gamma\frac{\partial\vec{E}}{\partial t},
\quad
\nabla^{2}\vec{J}_C=\mu\gamma\frac{\partial\vec{J}_C}{\partial t}
\]
\[
\text{正弦稳态：}\quad
\nabla^{2}\dot{\vec{H}}=j\omega\mu\gamma\dot{\vec{H}},
\quad
\nabla^{2}\dot{\vec{E}}=j\omega\mu\gamma\dot{\vec{E}},
\quad
\nabla^{2}\dot{\vec{J}}_C=j\omega\mu\gamma\dot{\vec{J}}_C
\]
\end{formulaBox}
它们也称\textbf{涡流方程}或\textbf{磁扩散方程}（时域形式描述磁场扩散过程）。取 $x>0$ 为导体、
正弦电流沿 $y$ 方向且只随 $x$ 变，则
$\mathrm{d}^{2}\dot{J}_y/\mathrm{d}x^{2}=k^{2}\dot{J}_y$，舍去随 $x$ 增大的项，得

\begin{formulaBox}
\[
k=\frac{1+j}{\delta},
\qquad
\dot{J}_y=\dot{J}_0e^{-kx}=\dot{J}_0e^{-x/\delta}e^{-jx/\delta},
\qquad
\dot{E}_y=\dot{E}_0e^{-kx},
\qquad
\dot{H}_z=-\frac{jk}{\omega\mu}\dot{E}_0e^{-kx}
\]
\[
\delta=\frac{1}{\alpha}=\sqrt{\frac{2}{\omega\mu\gamma}}
=\frac{1}{\sqrt{\pi f\mu\gamma}}
\]
\end{formulaBox}

$\delta$ 为\textbf{透入深度}：场量振幅衰减到表面值的 $1/e$（约 36.8\%）所经距离。场量振幅沿纵深
按 $e^{-x/\delta}$ 衰减且\textbf{相位随之改变}，愈深愈小——这种场量集中在导体表面附近的现象即
\textbf{集肤效应}。由 $\dot{H}_z/\dot{E}_y$ 还可见：导体内磁场\textbf{落后}电场 $45^{\circ}$、
幅值比为 $\sqrt{\omega\mu/\gamma}$（与理想介质中“同相、比值 $\sqrt{\mu/\varepsilon}$”不同）。

要点：$f$、$\gamma$、$\mu$ 越大则 $\delta$ 越小、集肤越显著（铜在 $50\,\mathrm{Hz}$ 时
$\delta\approx9.4\,\mathrm{mm}$）；经 $4.6\delta$ 振幅降为 $1\%$；深度大于数个 $\delta$ 后场量
近似为零但\textbf{不等于零}；$\delta$ 由无限大平面导体导出，只要表面曲率半径远大于 $\delta$
仍可近似使用。工程应用：高频表面淬火；高频线圈用空心铜管或镀银线。

\section{涡流及其损耗}

大块导体处在变化磁场中时，内部感应出\textbf{在大块导体内部自成闭合回路、呈旋涡状流动}的电流，
称\textbf{涡旋电流}（涡流）。它有两个效应：\textbf{热效应}（损耗发热）与\textbf{去磁效应}
（涡流磁场反抗外加场的变化）。工业上用热效应加热、冶炼金属，用去磁效应制电磁阻尼、电磁闸；
变压器铁心中的涡流损耗则须设法减小。MQS 下涡流场同样遵守上节的涡流方程。

把铁心叠片视为厚度 $d$ 的薄导电平板，外磁场沿板面法向、场量与 $y$、$z$ 无关，取板中心为
$x=0$，则

\begin{formulaBox}
\[
\dot{H}_z(x)=\dot{H}_0\frac{\operatorname{ch}(kx)}
{\operatorname{ch}(kd/2)},
\qquad
\dot{J}_y(x)=-\dot{H}_0k\,
\frac{\operatorname{sh}(kx)}{\operatorname{ch}(kd/2)},
\qquad
k=\frac{1+j}{\delta}
\]
\end{formulaBox}

磁场在\textbf{板中心最小}（涡流去磁），涡流密度\textbf{反对称于中心}、中心为零、表面最大，
愈深愈小，同样呈集肤效应。对电工钢片（$\mu_r\sim1000$、$d=0.5\,\mathrm{mm}$），$50\,\mathrm{Hz}$
时 $B$ 可视为沿截面均匀；$2000\,\mathrm{Hz}$ 时板中 $B$ 已远小于表面，须改用更薄钢片，更高频
则用粉状材料压制的铁心。体积 $V$ 中的平均损耗及低频近似为

\begin{formulaBox}
\[
P=\int_{V}\frac{1}{\gamma}\left|\dot{J}_y\right|^{2}\mathrm{d}V,
\qquad
P=\frac{1}{12}\gamma\omega^{2}d^{2}B_{av}^{2}V
\quad\left(\frac{d}{\delta}\ll1\right)
\]
\end{formulaBox}
$B_{av}$ 为磁感应强度在板厚上的平均值。减损三途径：\textbf{减小厚度}（$P\propto d^{2}$，故用
彼此绝缘的硅钢片叠装，片间绝缘切断涡流回路）、\textbf{减小电导率}（掺硅提高电阻率）、
\textbf{降低频率或磁密}。

\corepoint{低频下 $P\propto f^{2}B_{av}^{2}d^{2}\gamma$。三处反直觉：（1）$\gamma$ 越\textbf{大}
损耗越\textbf{大}，故靠“掺硅降 $\gamma$”减损，而非提高导电性；（2）$P\propto d^{2}$ 只在
$d\ll\delta$ 时成立，高频下须用 $P=\int|J|^{2}/\gamma\,\mathrm{d}V$ 精算；（3）硅钢片若片间
不绝缘，涡流回路重新闭合，减损效果全部丧失。}

\section{导体的交流内阻抗}

交流下集肤效应使电流与场集中于表面附近，\textbf{实际载流截面减小}，故电阻与内电感不同于直流。
设通有总电流 $\dot{I}$（有效值），等效参数 $Z=R+jX$：一方面 $\dot{I}\dot{I}^{*}Z=I^{2}(R+jX)$，
另一方面由坡印亭定理流入导体的复功率为
$-\oint_S(\dot{\vec{E}}\times\dot{\vec{H}}^{*})\cdot \mathrm{d}\vec{S}$，两者相等得

\begin{formulaBox}
\[
Z=\frac{-\oint_{S}\left(\dot{\vec{E}}\times\dot{\vec{H}}^{*}\right)
\cdot \mathrm{d}\vec{S}}{I^{2}}=R+jX
\]
\[
R=\frac{1}{I^{2}}\mathrm{Re}\left[-\oint_{S}
\left(\dot{\vec{E}}\times\dot{\vec{H}}^{*}\right)\cdot \mathrm{d}\vec{S}\right]
\]
\[
X=\frac{1}{I^{2}}\mathrm{Im}\left[-\oint_{S}
\left(\dot{\vec{E}}\times\dot{\vec{H}}^{*}\right)\cdot \mathrm{d}\vec{S}\right]
\]
\end{formulaBox}
$Z$ 只计及\textbf{导体内部}场引起的阻抗，故称\textbf{交流内阻抗}：$R$ 为交流电阻（焦耳损耗），
$X=\omega L_i$ 为内电抗（体内磁场储能）。

\textbf{（1）半无限大导体。}取底面积 $h\times a$ 的柱体，由 $\dot{I}=a\gamma\dot{E}_0/k$ 与坡印亭
矢量只在 $x=0$ 面有通量，得

\begin{formulaBox}
\[
Z=R+jX=\frac{h}{a\gamma\delta}(1+j),
\qquad
R=\frac{h}{a\gamma\delta},
\qquad
L_i=\frac{X}{\omega}=\frac{h}{a\gamma\delta\omega}
\]
\end{formulaBox}
意义：有效值 $I$ 的交流电流给出的有功功率，等于同值直流电流\textbf{集中在厚度 $\delta$ 的层内}
的功率——透入深度的另一物理含义；因 $\delta$ 很小，$Z$ 也称\textbf{表面阻抗}。

\textbf{（2）圆导线。}半径 $a\gg\delta$ 时可视为厚度无限大、宽为周长 $2\pi a$ 的平面导体：

\begin{formulaBox}
\[
R=\frac{1}{2\pi a\gamma\delta},
\qquad
R_{dc}=\frac{1}{\gamma\pi a^{2}},
\qquad
\frac{R}{R_{dc}}=\frac{a}{2\delta}
\]
\end{formulaBox}
即交流电阻相当于厚度为 $\delta$ 的圆管的直流电阻。故同截面下减小交流电阻的唯一办法是\textbf{增大
表面面积}——用相互绝缘的多股细线（绞线），高频时用空心铜管。

\corepoint{三处易错：（1）$\dfrac{R}{R_{dc}}=\dfrac{a}{2\delta}$ 要求 $a\gg\delta$（强集肤）；
低频时 $\delta\gg a$，有 $R\approx R_{dc}$，不可套用。（2）$L_i=X/\omega$ 随频率变化（$f$ 升高
$\delta$ 减小、$L_i$ 减小），非常数。（3）式中 $I$ 是\textbf{有效值}。}

\section{邻近效应和磁屏蔽}

\textbf{（1）邻近效应。}相互靠近的导体通交变电流时，每一导体还处于其它导体电流的场中，电流
分布因而改变，称\textbf{邻近效应}；\textbf{频率愈高、靠得愈近愈显著}，与集肤效应共存。二线
传输线中电流集中于两导线\textbf{相对的内侧}：能量主要经两线之间的空间以电磁波形式传给负载，
内侧能量密度大、传入导线的功率大；通以反向电流的平行汇流排，内侧面电流密度也最大。

\textbf{（2）电磁屏蔽。}用来减弱某些源在空间某区域（不含这些源）产生的电磁场的结构，称
\textbf{电磁屏蔽}，是抑制邻近效应的常用措施，屏蔽体多用铜、铝、钢。它利用\textbf{导体内涡流的
去磁作用}抵制外加场，故又称\textbf{涡流屏蔽}；有效屏蔽要求屏蔽层厚度接近透入深度的 3～6 倍：

\begin{formulaBox}
\[
h\ \geq\ (3\sim6)\,\delta,
\qquad
S=\frac{E_0}{E}\ \ \text{或}\ \ S=\frac{H_0}{H}
\]
\end{formulaBox}
$S$ 称\textbf{屏蔽系数}（无屏蔽体时该区场强与有屏蔽体时之比）。$1\,\mathrm{MHz}$ 时铝的
$\delta\approx82\,\mu\mathrm{m}$，薄铝罩即可屏蔽射频场；\textbf{高频不宜用铁磁材料}（损耗大、
发热重），而低频恒定磁场恰要靠高 $\mu$ 铁磁材料“短路”磁通来屏蔽。

\textbf{（3）磁扩散过程。}薄壁导体圆管（内半径 $a$、壁厚 $\Delta$）外突然建立轴向场 $H_0$ 时，
感应电流的去磁作用使管内 $H_i$ 不能突变，随后感应电流衰减、磁场逐渐透入，最后均匀分布：

\begin{formulaBox}
\[
\frac{\mathrm{d}H_i}{\mathrm{d}t}+\frac{1}{\tau_m}H_i=\frac{H_0}{\tau_m},
\qquad
H_i(t)=H_0\left(1-e^{-t/\tau_m}\right),
\qquad
\tau_m=\frac{\mu_0\gamma\Delta a}{2}
\]
\[
\text{正弦稳态：}\quad
\dot{H}_i=\frac{H_0}{1+j\omega\tau_m},
\qquad
\omega\gg\frac{1}{\tau_m}=\frac{2}{\mu_0\gamma\Delta a}
\ \text{时}\ H_i\ll H_0
\]
\end{formulaBox}
$\tau_m$ 称\textbf{磁扩散时间}。可见薄壁导体圆管中感应电流的作用是\textbf{去磁}的；若管内再套
一管，则内管磁场极弱——这正是多层屏蔽更有效的原因。

\corepoint{电磁屏蔽三处易错：（1）$h\gtrsim(3\sim6)\delta$ 而 $\delta\propto1/\sqrt{f}$，故
\textbf{频率越高屏蔽层越薄}；对低频（尤其直流）磁场，薄金属罩几乎无效，必须用高 $\mu$ 铁磁
材料。（2）屏蔽靠涡流去磁，屏蔽体必须\textbf{闭合、导电良好}，开缝开孔会显著降低效能。
（3）须避开\textbf{谐振}：被屏蔽场频率接近屏蔽体固有频率时效能急剧下降，甚至使防护区场强
反而增大。}

\end{document}